% =====================================================================
%  Anatomía de un LLM — Fase 0 — Módulo 0.3
%  Números, álgebra y sumatorios
%  Compilar: latexmk -r ../latexmkrc P03-algebra.tex
% =====================================================================
\documentclass[11pt,a4paper,twoside,openany]{book}
\usepackage{../estilo-apuntes}

\newcommand{\ApuntesCurso}{Anatomía de un LLM}
\newcommand{\ApuntesModulo}{Módulo 0.3}

% Anclas de hiperenlace únicas por tema para los enunciados
\makeatletter
\@for\nombre:=definicion,teorema,proposicion,lema,corolario,ejemplo,observacion\do{%
  \expandafter\xdef\csname theH\nombre\endcsname{\noexpand\thechapter.\noexpand\arabic{definicion}}}
\makeatother

\newcommand{\N}{\mathbb{N}}
\newcommand{\Z}{\mathbb{Z}}
\newcommand{\Q}{\mathbb{Q}}
\newcommand{\fl}{\operatorname{fl}}
\newcommand{\sgn}{\operatorname{sgn}}
\newcommand{\bits}[1]{\texttt{#1}}

\hypersetup{pdftitle={Módulo 0.3. Números, álgebra y sumatorios},pdfauthor={Anatomía de un LLM}}

\begin{document}
\frontmatter
% ------------------------------------------------------------------ Portada
\begin{titlepage}
\centering
\vspace*{2.5cm}
{\large\scshape Anatomía de un LLM}\par
\vspace{0.4cm}
{\normalsize Fundamentos matemáticos de los modelos de lenguaje}\par
\vspace{2.2cm}
\rule{\textwidth}{0.8pt}\par\vspace{0.5cm}
{\small\scshape Fase 0 \textperiodcentered\ Preliminares matemáticos}\par\vspace{0.3cm}
{\Huge\bfseries Módulo 0.3\par}\vspace{0.3cm}
{\LARGE Números, álgebra y sumatorios\par}
\vspace{0.5cm}\rule{\textwidth}{0.4pt}\par
\vspace{1.6cm}
{\large Apuntes del curso}\par
\vfill
{\small Edición de octubre de 2026}\par
\end{titlepage}

% ------------------------------------------------------- Presentación
\chapter*{Presentación del módulo}
\addcontentsline{toc}{chapter}{Presentación del módulo}
\markboth{Presentación del módulo}{}

Casi todos los cálculos del curso se reducen, en último término, a manipular números reales: desigualdades que acotan errores, polinomios que se anulan, sumas de logaritmos de probabilidades, series geométricas que aparecen en las medias móviles de los optimizadores y números binarios con los que la máquina representa todo lo anterior. Este módulo repasa esas herramientas desde el nivel de bachillerato con el rigor de un primer curso universitario: cada propiedad se enuncia con precisión, los resultados importantes se demuestran y se proponen muchos ejercicios, tanto de cálculo como de demostración.

\section*{Objetivos}
Al finalizar el módulo, el estudiante debe ser capaz de:
\begin{enumerate}[label=\arabic*.]
  \item deducir las propiedades aritméticas y de orden de los números reales a partir de los axiomas de cuerpo ordenado, y demostrar la irracionalidad de números como $\sqrt{2}$;
  \item manejar el valor absoluto, resolver inecuaciones y demostrar y aplicar las desigualdades triangular, de las medias, de Bernoulli y de Cauchy--Schwarz;
  \item dividir y factorizar polinomios, resolver ecuaciones algebraicas, aplicar las fórmulas de Vieta y resolver sistemas lineales pequeños;
  \item manipular sumatorios y productorios (cambio de índice, sumas aritméticas, geométricas y telescópicas, sumas dobles) y demostrar el binomio de Newton;
  \item operar con exponenciales y logaritmos, resolver ecuaciones con ellos y calcular de forma estable la función log-suma-exp;
  \item representar enteros y reales en binario, hexadecimal, complemento a dos y coma flotante, y estimar los errores de redondeo.
\end{enumerate}

\section*{Requisitos previos}
Matemáticas de bachillerato: operaciones con fracciones, potencias y radicales, ecuaciones de primer y segundo grado, y la idea de función. Se emplean las técnicas de demostración (directa, por contrarrecíproco, por reducción al absurdo y por inducción) de los módulos anteriores de esta fase; el principio de inducción se recuerda en el \cref{tema:4}.

\section*{Organización}
\begin{center}\small
\begin{tabular}{@{}cll@{}}
\toprule
Tema & Contenido & Fuentes principales \\
\midrule
1 & Los números reales & \cite{spivak}, \cite{openstax} \\
2 & Valor absoluto y desigualdades & \cite{spivak}, \cite{openstax}, \cite{hammack} \\
3 & Polinomios, ecuaciones y sistemas & \cite{openstax} \\
4 & Sumatorios, productorios y binomio de Newton & \cite{gkp}, \cite{mcs}, \cite{openstax}, \cite{hammack} \\
5 & Exponenciales y logaritmos & \cite{openstax}, \cite{ct}, \cite{blanchard} \\
6 & Sistemas de numeración y coma flotante & \cite{csapp}, \cite{goldberg}, \cite{ieee754} \\
\bottomrule
\end{tabular}
\end{center}

Cada tema termina con una lista de lecturas y con ejercicios graduados en tres niveles: $\star$ comprensión, $\star\star$ cálculo y $\star\star\star$ demostración. Los ejercicios marcados con \emph{Fuente} están adaptados de cursos y libros de texto universitarios; el enunciado se ha redactado de nuevo y, en algunos casos, se han cambiado los datos para que puedan resolverse a mano. Las soluciones completas, paso a paso, están en el \cref{ap:soluciones}.

\section*{Sobre las fuentes}
El texto es de redacción propia. La construcción axiomática de los números reales y el tratamiento de las desigualdades siguen los capítulos~1 y~2 de Spivak~\cite{spivak}. Las técnicas algebraicas de bachillerato (polinomios, ecuaciones, sistemas, exponenciales y logaritmos) siguen a OpenStax \emph{College Algebra}~\cite{openstax}. Los sumatorios siguen el capítulo~2 de Graham, Knuth y Patashnik~\cite{gkp} y el capítulo~14 de Lehman, Leighton y Meyer~\cite{mcs}, texto del curso 6.042J del MIT~\cite{mit6042}; las demostraciones por inducción, el capítulo~10 de Hammack~\cite{hammack}. La representación de números en la máquina sigue el capítulo~2 de Bryant y O'Hallaron~\cite{csapp}, texto de los cursos de sistemas de Carnegie Mellon, y el artículo de Goldberg~\cite{goldberg}.

\section*{Dónde se usa en el curso}
\begin{center}\small
\begin{tabularx}{\textwidth}{@{}>{\raggedright\arraybackslash}p{5.2cm}>{\raggedright\arraybackslash}X@{}}
\toprule
Concepto de este módulo & Módulos del curso que lo emplean \\
\midrule
Axiomas de cuerpo y orden; fallo de la asociatividad en coma flotante & M11 (precisión mixta), M17 (cotas de error, reducciones en paralelo) \\
Valor absoluto, desigualdad triangular & M02 (normas y distancias), M07 (convergencia), M17 (cotas de error) \\
Medias aritmética y geométrica, Cauchy--Schwarz & M02 (similitud coseno), M06 (desigualdad de Gibbs), M01 (perplejidad como media geométrica) \\
Ecuación de segundo grado, fórmulas de Vieta & M03 (polinomio característico, traza y determinante), M07 (optimización de cuadráticas) \\
Sistemas lineales pequeños & M02, M03 (eliminación gaussiana), M08 (ecuaciones normales) \\
Sumatorios, sumas dobles, productorios & M01 (probabilidad de una secuencia), M05 (esperanzas), M13 (coste de la atención causal) \\
Serie geométrica & M07 (medias exponenciales de Adam), M21 (retorno descontado) \\
Binomio de Newton y números combinatorios & M05 (distribución binomial), M22 (pass@$k$) \\
Logaritmos, bits y nats & M01 (log-verosimilitud, perplejidad), M06 (entropía y códigos), M16 (leyes de potencia en escala logarítmica) \\
Log-suma-exp & M01 y M04 (softmax), M08 (entropía cruzada), M13 (atención) \\
Binario, hexadecimal, bytes & M12 (UTF-8 y BPE a nivel de byte) \\
Complemento a dos y coma flotante & M11 (fp16, bf16), M17 (cuantización int8, memoria de los modelos) \\
\bottomrule
\end{tabularx}
\end{center}

\section*{Notación}
\begin{center}
\begin{tabular}{@{}ll@{}}
\toprule
Símbolo & Significado \\
\midrule
$\N=\{1,2,3,\dots\}$, $\N_0=\N\cup\{0\}$ & números naturales, con y sin el cero \\
$\Z$, $\Q$, $\R$ & enteros, racionales y reales \\
$\abs{x}$ & valor absoluto de $x$ \\
$\sum_{k=m}^{n}a_k$, $\prod_{k=m}^{n}a_k$ & sumatorio y productorio \\
$n!$, $\binom{n}{k}$ & factorial y número combinatorio \\
$\ln$, $\log_2$, $\log_{10}$ & logaritmos natural, binario y decimal \\
$\LSE(z)=\ln\sum_i e^{z_i}$ & función log-suma-exp \\
$(d_{m}\dots d_1d_0)_b$, \texttt{0x}\dots & número escrito en base $b$; prefijo hexadecimal \\
$\lfloor x\rfloor$, $\lceil x\rceil$ & parte entera por defecto y por exceso \\
$\fl(x)$, $\varepsilon$ & redondeo a coma flotante y épsilon de la máquina \\
\bottomrule
\end{tabular}
\end{center}

\tableofcontents

\mainmatter
% =====================================================================
\chapter{Los números reales}
\label{tema:1}
% =====================================================================

Este tema fija las reglas de juego de todo el módulo: qué se puede hacer con los números reales y por qué. En lugar de dar por sabidas las propiedades de la aritmética, se parte de una lista corta de axiomas, la de un \emph{cuerpo ordenado}, y se deducen de ella las reglas habituales. La presentación sigue los capítulos~1 y~2 de Spivak~\cite{spivak}; las secciones de potencias y radicales siguen los §1.1--1.3 de OpenStax~\cite{openstax}.

\section{Conjuntos numéricos}

Se trabaja con la cadena de inclusiones $\N\subset\Z\subset\Q\subset\R$, donde $\N=\{1,2,3,\dots\}$ son los naturales, $\Z$ los enteros, $\Q=\{p/q: p,q\in\Z,\ q\neq 0\}$ los racionales y $\R$ los reales. Un real que no es racional se llama \emph{irracional}.

\begin{proposicion}\label{prop:decimal-periodico}
Un número real es racional si y solo si su expresión decimal es finita o periódica.
\end{proposicion}

\begin{proof}[Idea de la demostración]
Si $x=p/q$ con $q>0$, la división larga de $p$ entre $q$ solo puede producir los restos $0,1,\dots,q-1$; en a lo sumo $q$ pasos un resto se repite y, desde ese momento, las cifras se repiten. Recíprocamente, si $x=0{,}\overline{d_1\dots d_k}$, entonces $10^{k}x-x=(d_1\dots d_k)$ es un entero, luego $x=(d_1\dots d_k)/(10^{k}-1)\in\Q$; el caso de un anteperiodo se reduce a este multiplicando por una potencia de~10.
\end{proof}

\begin{ejemplo}\label{ej:periodico}
Sea $x=0{,}\overline{36}$. Entonces $100x=36{,}\overline{36}$ y $100x-x=36$, de donde $x=36/99=4/11$.
\end{ejemplo}

\section{Axiomas de cuerpo}

\begin{definicion}[Cuerpo]\label{def:cuerpo}
Un \emph{cuerpo} es un conjunto $K$ con dos operaciones, suma $(a,b)\mapsto a+b$ y producto $(a,b)\mapsto ab$, que cumplen, para todos $a,b,c\in K$:
\begin{enumerate}[label=(C\arabic*)]
  \item asociatividad: $(a+b)+c=a+(b+c)$ y $(ab)c=a(bc)$;
  \item conmutatividad: $a+b=b+a$ y $ab=ba$;
  \item elementos neutros: existen $0,1\in K$, con $0\neq 1$, tales que $a+0=a$ y $a\cdot 1=a$;
  \item opuesto: para cada $a$ existe $-a$ con $a+(-a)=0$;
  \item inverso: para cada $a\neq 0$ existe $a^{-1}$ con $aa^{-1}=1$;
  \item distributividad: $a(b+c)=ab+ac$.
\end{enumerate}
Se definen la resta $a-b=a+(-b)$ y, si $b\neq 0$, la división $a/b=ab^{-1}$.
\end{definicion}

$\Q$ y $\R$ son cuerpos; $\Z$ no lo es, porque $2$ no tiene inverso en $\Z$. Todas las reglas de la aritmética escolar se deducen de (C1)--(C6).

\begin{proposicion}[Consecuencias de los axiomas de cuerpo]\label{prop:cuerpo}
Para todos $a,b,x$ de un cuerpo:
\begin{enumerate}
  \item si $a+x=a$, entonces $x=0$ (unicidad del neutro); si $a+x=0$, entonces $x=-a$ (unicidad del opuesto);
  \item $a\cdot 0=0$;
  \item $(-a)b=-(ab)$ y $(-a)(-b)=ab$;
  \item si $ab=0$, entonces $a=0$ o $b=0$.
\end{enumerate}
\end{proposicion}

\begin{proof}
(i) Sumando $-a$ a ambos lados de $a+x=a$ y usando (C1), (C4) y (C3): $x=0+x=(-a+a)+x=-a+(a+x)=-a+a=0$. El segundo apartado es análogo: $x=(-a+a)+x=-a+(a+x)=-a+0=-a$.

(ii) Por (C3) y (C6), $a\cdot 0=a(0+0)=a\cdot 0+a\cdot 0$; por (i) aplicado a $a\cdot 0$, se obtiene $a\cdot 0=0$.

(iii) $ab+(-a)b=(a+(-a))b=0\cdot b=0$, luego, por la unicidad del opuesto, $(-a)b=-(ab)$. Aplicándolo dos veces, $(-a)(-b)=-(a(-b))=-(-(ab))=ab$, donde $-(-c)=c$ porque $c+(-c)=0$ dice que $c$ es el opuesto de $-c$.

(iv) Si $a\neq 0$, se multiplica $ab=0$ por $a^{-1}$: $b=(a^{-1}a)b=a^{-1}(ab)=a^{-1}\cdot 0=0$.
\end{proof}

El apartado~(iv) es el que justifica resolver $(x-2)(x+3)=0$ igualando a cero cada factor.

\section{Axiomas de orden}

\begin{definicion}[Cuerpo ordenado]\label{def:orden}
Un cuerpo $K$ es \emph{ordenado} si existe un subconjunto $P\subset K$, el de los \emph{positivos}, tal que:
\begin{enumerate}[label=(O\arabic*)]
  \item tricotomía: para cada $a\in K$ se cumple exactamente una de $a=0$, $a\in P$, $-a\in P$;
  \item si $a,b\in P$, entonces $a+b\in P$;
  \item si $a,b\in P$, entonces $ab\in P$.
\end{enumerate}
Se escribe $a<b$ (o $b>a$) si $b-a\in P$, y $a\le b$ si $a<b$ o $a=b$.
\end{definicion}

\begin{proposicion}[Reglas de las desigualdades]\label{prop:orden}
En un cuerpo ordenado, para todos $a,b,c$:
\begin{enumerate}
  \item si $a<b$ y $b<c$, entonces $a<c$;
  \item si $a<b$, entonces $a+c<b+c$;
  \item si $a<b$ y $c>0$, entonces $ac<bc$;
  \item si $a\neq 0$, entonces $a^2>0$; en particular, $1>0$.
\end{enumerate}
\end{proposicion}

\begin{proof}
(i) $c-a=(c-b)+(b-a)$ es suma de dos positivos, luego positivo por (O2). (ii) $(b+c)-(a+c)=b-a\in P$. (iii) $bc-ac=(b-a)c$ es producto de dos positivos, positivo por (O3). (iv) Por (O1), $a\in P$ o $-a\in P$. En el primer caso $a^2=a\cdot a\in P$; en el segundo, $a^2=(-a)(-a)\in P$ por la \cref{prop:cuerpo}(iii). Como $1=1^2$ y $1\neq 0$, se tiene $1>0$.
\end{proof}

\begin{observacion}[La completitud]
$\Q$ y $\R$ son ambos cuerpos ordenados; lo que los distingue es el \emph{axioma del supremo}: todo subconjunto no vacío de $\R$ acotado superiormente tiene una cota superior mínima~\cite[cap.~8]{spivak}. Gracias a él existen en $\R$ números como $\sqrt 2$, que no están en $\Q$ (\cref{teo:raiz2}). En este módulo el axioma solo se usa implícitamente, al admitir que existen las raíces $n$-ésimas y los logaritmos de los reales positivos.
\end{observacion}

\section{Números irracionales}

\begin{lema}\label{lem:par}
Sea $n\in\Z$. Si $n^2$ es par, entonces $n$ es par.
\end{lema}

\begin{proof}
Por contrarrecíproco. Si $n$ es impar, $n=2k+1$ con $k\in\Z$, y $n^2=4k^2+4k+1=2(2k^2+2k)+1$ es impar.
\end{proof}

\begin{teorema}\label{teo:raiz2}
$\sqrt2$ es irracional.
\end{teorema}

\begin{proof}
Por reducción al absurdo, supóngase $\sqrt2=p/q$ con $p,q\in\Z$, $q\neq0$, y la fracción irreducible. Elevando al cuadrado, $p^2=2q^2$, luego $p^2$ es par y, por el \cref{lem:par}, $p=2r$. Sustituyendo, $4r^2=2q^2$, es decir, $q^2=2r^2$, y por el mismo lema $q$ es par. Entonces $p$ y $q$ son ambos pares, en contra de que la fracción sea irreducible.
\end{proof}

\section{Potencias, radicales e intervalos}

Para $a\in\R$ y $n\in\N$ se define $a^n=a\cdots a$ ($n$ factores), $a^0=1$ si $a\neq 0$, y $a^{-n}=1/a^n$. Para $a>0$ se define $a^{1/n}=\sqrt[n]{a}$ como el único real positivo cuya potencia $n$-ésima es $a$, y $a^{m/n}=(a^{1/n})^m$.

\begin{proposicion}[Leyes de los exponentes]\label{prop:exponentes}
Para $a,b>0$ y $r,s\in\Q$: $a^{r}a^{s}=a^{r+s}$, $(a^{r})^{s}=a^{rs}$, $(ab)^{r}=a^{r}b^{r}$ y $a^{-r}=1/a^{r}$. Para exponentes enteros, las tres primeras valen también con $a,b\neq 0$.
\end{proposicion}

Para exponentes naturales, las leyes se prueban por inducción sobre el exponente; para exponentes racionales se reducen a ese caso usando la unicidad de la raíz $n$-ésima~\cite[§1.2--1.3]{openstax}.

\begin{definicion}[Intervalos]
Para $a<b$: $[a,b]=\{x: a\le x\le b\}$, $(a,b)=\{x: a<x<b\}$, y análogamente $[a,b)$, $(a,b]$, $(a,\infty)$, $(-\infty,b]$, etc. Los símbolos $\pm\infty$ no son números reales y nunca llevan corchete cerrado.
\end{definicion}

\begin{observacion}[Errores frecuentes]
\begin{enumerate}[label=(\alph*)]
\item $\sqrt{a^2}=\abs{a}$, no $a$: con $a=-3$, $\sqrt{9}=3$.
\item $(a+b)^2\neq a^2+b^2$ y $\sqrt{a+b}\neq\sqrt a+\sqrt b$ en general.
\item Al multiplicar una desigualdad por un número negativo, su sentido se invierte (\cref{eje:1-negativo}).
\item No se puede dividir por una expresión que puede valer cero: de $x^2=x$ no se sigue $x=1$, porque se pierde la solución $x=0$.
\end{enumerate}
\end{observacion}

\begin{observacion}[Conexión con el curso]
Los ordenadores no trabajan con $\R$ sino con un conjunto finito de números en coma flotante (\cref{tema:6}). En ese conjunto la suma no es asociativa: con el formato bfloat16, $(256+1)+1=256$ mientras que $256+(1+1)=258$ (\cref{eje:6-espaciado}). Por eso el resultado de sumar los gradientes de muchas GPU depende del orden de la reducción (módulos~11 y~17), y por eso los axiomas de esta sección son una idealización que hay que vigilar en la práctica.
\end{observacion}

\begin{lecturas}
\item Spivak~\cite{spivak}, cap.~1 «Basic Properties of Numbers» (axiomas P1--P12 y sus consecuencias) y cap.~2 «Numbers of Various Sorts» (irracionalidad de $\sqrt2$).
\item OpenStax \emph{College Algebra}~\cite{openstax}, §1.1 «Real Numbers: Algebra Essentials», §1.2 «Exponents and Scientific Notation» y §1.3 «Radicals and Rational Exponents».
\item Hammack~\cite{hammack}, cap.~6 «Proof by Contradiction», para practicar demostraciones de irracionalidad.
\end{lecturas}

\section*{Ejercicios del tema 1}
\addcontentsline{toc}{section}{Ejercicios del tema 1}

\begin{ejercicio}{1}{eje:1-clasificar}
Indique a cuáles de los conjuntos $\N$, $\Z$, $\Q$ y $\R\setminus\Q$ pertenece cada número: $-3$, $0{,}75$, $0{,}\overline{3}$, $\sqrt9$, $\sqrt8$, $\pi$, $1/7$ y $0$.
\end{ejercicio}

\begin{ejercicio}{1}{eje:1-axiomas}
En la siguiente deducción de «si $a+x=a+y$, entonces $x=y$», indique el axioma (C1)--(C6) que justifica cada igualdad:
\[
x\overset{(1)}{=}0+x\overset{(2)}{=}((-a)+a)+x\overset{(3)}{=}(-a)+(a+x)\overset{(4)}{=}(-a)+(a+y)\overset{(5)}{=}((-a)+a)+y\overset{(6)}{=}0+y\overset{(7)}{=}y.
\]
\end{ejercicio}

\begin{ejercicio}{1}{eje:1-vf}
Razone si son verdaderas o falsas: (a) la suma de dos irracionales es irracional; (b) el producto de un racional no nulo por un irracional es irracional; (c) entre dos racionales distintos siempre hay otro racional.
\end{ejercicio}

\begin{ejercicio}{2}{eje:1-fracciones}
Escriba como fracción irreducible: (a) $0{,}1\overline{6}$; (b) $2{,}\overline{45}$; (c) $0{,}\overline{142857}$.
\end{ejercicio}

\begin{ejercicio}{2}{eje:1-potencias}
Simplifique sin calculadora: (a) $\dfrac{2^{-3}\cdot 4^{2}}{8^{-1}}$; (b) $27^{2/3}\cdot 16^{-3/4}$; (c) $\sqrt{50}-\sqrt{18}+\sqrt{8}$.
\fuente{adaptado de OpenStax \emph{College Algebra}~\cite{openstax}, §1.2 y §1.3.}
\end{ejercicio}

\begin{ejercicio}{2}{eje:1-racionalizar}
Racionalice los denominadores: (a) $\dfrac{1}{\sqrt3-\sqrt2}$; (b) $\dfrac{2}{1+\sqrt5}$; (c) $\dfrac{6}{\sqrt[3]{2}}$.
\fuente{adaptado de OpenStax \emph{College Algebra}~\cite{openstax}, §1.3.}
\end{ejercicio}

\begin{ejercicio}{2}{eje:1-ordenar}
Ordene de menor a mayor $\sqrt2+\sqrt3$, $\sqrt{10}$ y $\pi$ sin calculadora, usando solo que $3{,}1415<\pi<3{,}1416$. (Indicación: compare los cuadrados y aplique el \cref{eje:1-potencia-n}.)
\end{ejercicio}

\begin{ejercicio}{2}{eje:1-cientifica}
(a) Calcule $\dfrac{(3{,}2\cdot 10^{5})(4\cdot 10^{-8})}{1{,}6\cdot 10^{-2}}$ en notación científica. (b) Un modelo tiene $7\cdot 10^{10}$ parámetros y cada uno ocupa 2 bytes. ¿Cuántos bytes ocupa? Expréselo en gigabytes ($1\ \text{GB}=10^{9}$ bytes).
\fuente{adaptado de OpenStax \emph{College Algebra}~\cite{openstax}, §1.2 (notación científica).}
\end{ejercicio}

\begin{ejercicio}{2}{eje:1-intervalos}
Sean $A=[-2,3)$ y $B=(1,5]$. Escriba como intervalos o uniones de intervalos: $A\cap B$, $A\cup B$, $A\setminus B$ y $\R\setminus A$.
\end{ejercicio}

\begin{ejercicio}{2}{eje:1-ecuacion-producto}
Resuelva $x^3=x$ en $\R$ justificando cada paso con la \cref{prop:cuerpo}. Explique qué error se comete al «simplificar» dividiendo por $x$.
\end{ejercicio}

\begin{ejercicio}{3}{eje:1-quebrados}
Demuestre, a partir de los axiomas de cuerpo, que si $b\neq0$ y $d\neq0$, entonces
\[
\frac ab+\frac cd=\frac{ad+bc}{bd}\qquad\text{y}\qquad \frac ab\cdot\frac cd=\frac{ac}{bd}.
\]
(Puede usar sin demostración que $(bd)^{-1}=b^{-1}d^{-1}$.)
\fuente{adaptado de Spivak~\cite{spivak}, cap.~1, problemas finales.}
\end{ejercicio}

\begin{ejercicio}{3}{eje:1-negativo}
Demuestre, a partir de los axiomas de orden: (a) si $a<b$ y $c<0$, entonces $ac>bc$; (b) si $a>0$, entonces $1/a>0$; (c) si $0<a<b$, entonces $1/b<1/a$.
\fuente{adaptado de Spivak~\cite{spivak}, cap.~1, problemas finales.}
\end{ejercicio}

\begin{ejercicio}{3}{eje:1-potencia-n}
Demuestre por inducción que, si $0\le a<b$, entonces $a^n<b^n$ para todo $n\in\N$. Deduzca que, si $a,b\ge 0$ y $a^2<b^2$, entonces $a<b$.
\fuente{adaptado de Spivak~\cite{spivak}, cap.~1, problemas finales.}
\end{ejercicio}

\begin{ejercicio}{3}{eje:1-raiz3}
Demuestre que $\sqrt3$ es irracional y, a partir de ello, que $\sqrt2+\sqrt3$ también lo es.
\fuente{adaptado de Spivak~\cite{spivak}, cap.~2, problemas finales.}
\end{ejercicio}

\begin{ejercicio}{3}{eje:1-epsilon}
Sean $x,y\in\R$. Demuestre: (a) si $a<b$, entonces $a<\frac{a+b}{2}<b$; (b) si $x\le y+\varepsilon$ para todo $\varepsilon>0$, entonces $x\le y$.
\end{ejercicio}

% =====================================================================
\chapter{Valor absoluto y desigualdades}
\label{tema:2}
% =====================================================================

Gran parte del trabajo matemático en aprendizaje automático consiste en \emph{acotar}: garantizar que un error es pequeño, que una pérdida no es negativa o que un método converge. Este tema reúne las desigualdades elementales que se usan para ello. El valor absoluto y la desigualdad triangular siguen el capítulo~1 de Spivak~\cite{spivak}; las inecuaciones, el §2.7 de OpenStax~\cite{openstax}; las demostraciones por inducción, el capítulo~10 de Hammack~\cite{hammack}.

\section{Valor absoluto}

\begin{definicion}[Valor absoluto]\label{def:abs}
Para $x\in\R$ se define $\abs{x}=x$ si $x\ge 0$ y $\abs{x}=-x$ si $x<0$. La \emph{distancia} entre $x$ e $y$ es $\abs{x-y}$.
\end{definicion}

\begin{proposicion}\label{prop:abs}
Para todos $x,y\in\R$ y $r\ge 0$:
\begin{enumerate}
  \item $\abs{x}\ge 0$, con igualdad si y solo si $x=0$; $\abs{-x}=\abs{x}$;
  \item $\abs{x}=\max(x,-x)$ y $-\abs{x}\le x\le\abs{x}$;
  \item $\abs{x}=\sqrt{x^2}$ y $\abs{xy}=\abs{x}\abs{y}$;
  \item $\abs{x}\le r$ si y solo si $-r\le x\le r$; $\abs{x}\ge r$ si y solo si $x\le -r$ o $x\ge r$.
\end{enumerate}
\end{proposicion}

\begin{proof}
(i) y (ii) se comprueban separando los casos $x\ge0$ y $x<0$. (iii) $\abs{x}$ es el único número no negativo cuyo cuadrado es $x^2$, luego coincide con $\sqrt{x^2}$; entonces $\abs{xy}=\sqrt{x^2y^2}=\sqrt{x^2}\sqrt{y^2}=\abs{x}\abs{y}$. (iv) Si $x\ge0$, $\abs{x}\le r$ equivale a $x\le r$, y $-r\le x$ es automático; si $x<0$, $\abs{x}\le r$ equivale a $-x\le r$, es decir, $x\ge -r$, y $x\le r$ es automático. La segunda equivalencia es la negación de la primera con desigualdad estricta.
\end{proof}

\section{La desigualdad triangular}

\begin{teorema}[Desigualdad triangular]\label{teo:triangular}
Para todos $x,y\in\R$, $\abs{x+y}\le\abs{x}+\abs{y}$, con igualdad si y solo si $xy\ge0$.
\end{teorema}

\begin{proof}
Ambos miembros son no negativos, así que basta comparar sus cuadrados (\cref{eje:1-potencia-n}):
\[
(\abs{x}+\abs{y})^2-\abs{x+y}^2=x^2+2\abs{x}\abs{y}+y^2-(x^2+2xy+y^2)=2(\abs{xy}-xy)\ge0,
\]
porque $\abs{xy}\ge xy$ por la \cref{prop:abs}(ii). La diferencia es cero si y solo si $\abs{xy}=xy$, es decir, si y solo si $xy\ge0$.
\end{proof}

\begin{corolario}[Desigualdad triangular inversa]\label{cor:triangular-inversa}
Para todos $x,y\in\R$, $\bigl|\abs{x}-\abs{y}\bigr|\le\abs{x-y}$.
\end{corolario}

\begin{proof}
Por el \cref{teo:triangular}, $\abs{x}=\abs{(x-y)+y}\le\abs{x-y}+\abs{y}$, luego $\abs{x}-\abs{y}\le\abs{x-y}$. Intercambiando los papeles, $\abs{y}-\abs{x}\le\abs{y-x}=\abs{x-y}$. Las dos desigualdades juntas dicen que $\abs{x}-\abs{y}\in[-\abs{x-y},\abs{x-y}]$, que es el enunciado por la \cref{prop:abs}(iv).
\end{proof}

La desigualdad triangular se generaliza a $n$ sumandos, $\abs{\sum_i x_i}\le\sum_i\abs{x_i}$ (\cref{eje:2-triangular-n}), y a vectores, donde $\abs{\cdot}$ se sustituye por una norma (módulo~2).

\section{Inecuaciones con valor absoluto}

Por la \cref{prop:abs}(iv), la condición $\abs{x-a}<r$ describe el intervalo $(a-r,a+r)$ de centro $a$ y radio $r$. Las inecuaciones con varios valores absolutos se resuelven separando casos según el signo de cada expresión.

\begin{ejemplo}\label{ej:abs-casos}
Resolver $\abs{x-1}+\abs{x+2}\ge 5$. Los puntos de corte son $-2$ y $1$.
Si $x\ge1$, la inecuación es $(x-1)+(x+2)\ge5$, es decir, $x\ge2$.
Si $-2\le x<1$, es $(1-x)+(x+2)=3\ge5$, imposible.
Si $x<-2$, es $(1-x)-(x+2)\ge5$, es decir, $x\le-3$.
La solución es $(-\infty,-3]\cup[2,\infty)$: los puntos cuya suma de distancias a $1$ y a $-2$ es al menos~5.
\end{ejemplo}

\section{Desigualdad entre las medias aritmética y geométrica}

\begin{proposicion}[Caso de dos números]\label{prop:amgm2}
Si $a,b\ge0$, entonces $\sqrt{ab}\le\dfrac{a+b}{2}$, con igualdad si y solo si $a=b$.
\end{proposicion}

\begin{proof}
$\dfrac{a+b}{2}-\sqrt{ab}=\dfrac{(\sqrt a-\sqrt b)^2}{2}\ge0$, y la igualdad equivale a $\sqrt a=\sqrt b$.
\end{proof}

\begin{lema}\label{lem:producto1}
Si $x_1,\dots,x_n>0$ y $x_1x_2\cdots x_n=1$, entonces $x_1+\dots+x_n\ge n$, con igualdad si y solo si todos valen~1.
\end{lema}

\begin{proof}
Por inducción sobre $n$. Para $n=1$, $x_1=1$. Supóngase cierto para $n$ y sean $x_1,\dots,x_{n+1}>0$ de producto~1. Si todos valen~1, la suma es $n+1$. Si no, hay alguno mayor que~1 y alguno menor que~1; reordenando, $x_n>1>x_{n+1}$. Entonces $(x_n-1)(1-x_{n+1})>0$, que desarrollado da
\[
x_n+x_{n+1}>1+x_nx_{n+1}.
\]
Los $n$ números $x_1,\dots,x_{n-1},\,x_nx_{n+1}$ tienen producto~1, luego, por la hipótesis de inducción, $x_1+\dots+x_{n-1}+x_nx_{n+1}\ge n$. Sumando,
\[
x_1+\dots+x_{n+1}>x_1+\dots+x_{n-1}+x_nx_{n+1}+1\ge n+1.\qedhere
\]
\end{proof}

\begin{teorema}[Desigualdad de las medias]\label{teo:amgm}
Si $a_1,\dots,a_n>0$, entonces
\[
\sqrt[n]{a_1a_2\cdots a_n}\le\frac{a_1+a_2+\dots+a_n}{n},
\]
con igualdad si y solo si $a_1=a_2=\dots=a_n$.
\end{teorema}

\begin{proof}
Sea $G=\sqrt[n]{a_1\cdots a_n}>0$ y $x_i=a_i/G$. Entonces $x_1\cdots x_n=(a_1\cdots a_n)/G^n=1$ y, por el \cref{lem:producto1}, $\sum_i a_i/G\ge n$, es decir, $\frac1n\sum_ia_i\ge G$. La igualdad se da si y solo si todos los $x_i$ valen~1, es decir, si todos los $a_i$ valen~$G$.
\end{proof}

\begin{ejemplo}\label{ej:amgm-optim}
Para $x>0$, $x+\dfrac 9x\ge 2\sqrt{x\cdot\dfrac9x}=6$, con igualdad si y solo si $x=9/x$, es decir, $x=3$. Así se obtiene un mínimo sin derivar.
\end{ejemplo}

\section{Desigualdades de Bernoulli y de Cauchy--Schwarz}

\begin{teorema}[Desigualdad de Bernoulli]\label{teo:bernoulli}
Si $x\ge-1$ y $n\in\N$, entonces $(1+x)^n\ge1+nx$.
\end{teorema}

\begin{proof}
Por inducción sobre $n$. Para $n=1$ hay igualdad. Si $(1+x)^n\ge1+nx$, como $1+x\ge0$ se puede multiplicar sin cambiar el sentido:
\[
(1+x)^{n+1}\ge(1+nx)(1+x)=1+(n+1)x+nx^2\ge1+(n+1)x.\qedhere
\]
\end{proof}

\begin{teorema}[Desigualdad de Cauchy--Schwarz]\label{teo:cs}
Para reales $a_1,\dots,a_n$ y $b_1,\dots,b_n$,
\[
\Bigl(\sum_{i=1}^na_ib_i\Bigr)^2\le\Bigl(\sum_{i=1}^na_i^2\Bigr)\Bigl(\sum_{i=1}^nb_i^2\Bigr).
\]
\end{teorema}

\begin{proof}
Si todos los $a_i$ son nulos, ambos miembros valen~0. En otro caso, sea $A=\sum a_i^2>0$, $B=\sum a_ib_i$ y $C=\sum b_i^2$. Para todo $t\in\R$,
\[
0\le\sum_{i=1}^n(ta_i+b_i)^2=At^2+2Bt+C.
\]
En particular, para $t=-B/A$ se obtiene $0\le B^2/A-2B^2/A+C=C-B^2/A$, y multiplicando por $A>0$, $B^2\le AC$, que es la desigualdad buscada. (Equivalentemente: el discriminante $4B^2-4AC$ de esa cuadrática no negativa es $\le0$; véase el \cref{eje:3-cuadratica-positiva}.)
\end{proof}

\begin{observacion}[Errores frecuentes]
\begin{enumerate}[label=(\alph*)]
\item $\abs{a+b}=\abs a+\abs b$ es falso en general: $\abs{3+(-5)}=2<8$.
\item La desigualdad de las medias exige números no negativos: con $a=b=-1$, $\sqrt{ab}=1>-1=\frac{a+b}2$.
\item Elevar al cuadrado ambos miembros de una desigualdad solo es válido si ambos son no negativos: $-3<2$ pero $9>4$.
\item En $\abs{x}<r$ con $r\le0$ no hay solución; en $\abs{x}>r$ con $r<0$ lo es todo $\R$.
\end{enumerate}
\end{observacion}

\begin{observacion}[Conexión con el curso]
La desigualdad triangular para normas acota errores acumulados (módulos~7 y~17). Cauchy--Schwarz garantiza que la similitud coseno está en $[-1,1]$ (módulo~2). La desigualdad de las medias dice que la media geométrica de las probabilidades que un modelo asigna a un texto nunca supera su media aritmética; la perplejidad del módulo~1 es precisamente el inverso de esa media geométrica. Su generalización, la desigualdad de Jensen, da la desigualdad de Gibbs del módulo~6.
\end{observacion}

\begin{lecturas}
\item Spivak~\cite{spivak}, cap.~1 (valor absoluto y desigualdad triangular) y sus problemas finales, que incluyen la desigualdad de las medias.
\item OpenStax \emph{College Algebra}~\cite{openstax}, §2.7 «Linear Inequalities and Absolute Value Inequalities».
\item Hammack~\cite{hammack}, cap.~10 «Mathematical Induction», para las demostraciones de desigualdades por inducción.
\end{lecturas}

\section*{Ejercicios del tema 2}
\addcontentsline{toc}{section}{Ejercicios del tema 2}

\begin{ejercicio}{1}{eje:2-calcular}
Calcule $\abs{-3}-\abs{2-7}$ y escriba sin barras $\abs{\pi-4}$ y $\abs{x-1}$ cuando $x<1$.
\end{ejercicio}

\begin{ejercicio}{1}{eje:2-vf}
Razone si son verdaderas o falsas para todos $a,b\in\R$: (a) $\abs{a-b}\le\abs a-\abs b$; (b) $\abs{a+b}=\abs a+\abs b$ si y solo si $ab\ge0$; (c) $\sqrt{x^2}=x$.
\end{ejercicio}

\begin{ejercicio}{1}{eje:2-traducir}
Exprese con un valor absoluto: (a) el intervalo $(2,8)$; (b) el intervalo $[-1,4]$; (c) «$x$ está a distancia menor que $0{,}01$ de~$3$».
\end{ejercicio}

\begin{ejercicio}{2}{eje:2-inec-simple}
Resuelva: (a) $\abs{3x-2}<7$; (b) $\abs{2x-1}=\abs{x+3}$.
\fuente{adaptado de OpenStax \emph{College Algebra}~\cite{openstax}, §2.7.}
\end{ejercicio}

\begin{ejercicio}{2}{eje:2-inec-casos}
Resuelva $\abs{x+1}\ge 2x-4$.
\end{ejercicio}

\begin{ejercicio}{2}{eje:2-inec-suma}
Resuelva $\abs{x-1}+\abs{x+2}\le5$ y compare el resultado con el \cref{ej:abs-casos}.
\fuente{adaptado de OpenStax \emph{College Algebra}~\cite{openstax}, §2.7.}
\end{ejercicio}

\begin{ejercicio}{2}{eje:2-amgm-aplic}
Usando la desigualdad de las medias: (a) halle el mínimo de $x+\frac4x$ para $x>0$; (b) entre todos los rectángulos de perímetro~20, halle el de área máxima; (c) halle el mínimo de $x+y+z$ si $x,y,z>0$ y $xyz=8$.
\end{ejercicio}

\begin{ejercicio}{2}{eje:2-bernoulli-num}
Use la desigualdad de Bernoulli para demostrar, sin calcular las potencias, que $1{,}01^{100}\ge2$ y que $0{,}99^{50}\ge\frac12$. Compare con los valores reales.
\end{ejercicio}

\begin{ejercicio}{2}{eje:2-cs-num}
(a) Compruebe la desigualdad de Cauchy--Schwarz para $a=(1,2,3)$ y $b=(4,-1,2)$. (b) Halle el valor máximo de $2x+3y$ sobre los puntos que cumplen $x^2+y^2=13$ y el punto en que se alcanza.
\end{ejercicio}

\begin{ejercicio}{3}{eje:2-triangular-n}
Demuestre por inducción que $\abs{x_1+\dots+x_n}\le\abs{x_1}+\dots+\abs{x_n}$ para todo $n\in\N$.
\fuente{adaptado de Spivak~\cite{spivak}, cap.~2, problemas finales (inducción).}
\end{ejercicio}

\begin{ejercicio}{3}{eje:2-maxmin}
Demuestre que, para todos $x,y\in\R$,
\[
\max(x,y)=\frac{x+y+\abs{x-y}}{2}\qquad\text{y}\qquad\min(x,y)=\frac{x+y-\abs{x-y}}{2}.
\]
\fuente{adaptado de Spivak~\cite{spivak}, cap.~1, problemas finales.}
\end{ejercicio}

\begin{ejercicio}{3}{eje:2-cadena-medias}
Demuestre que si $0<a<b$, entonces $a<\sqrt{ab}<\dfrac{a+b}{2}<b$.
\fuente{adaptado de Spivak~\cite{spivak}, cap.~1, problemas finales.}
\end{ejercicio}

\begin{ejercicio}{3}{eje:2-media-armonica}
Sean $a,b>0$. Demuestre que
\[
\frac{2ab}{a+b}\le\sqrt{ab}\le\frac{a+b}{2}\le\sqrt{\frac{a^2+b^2}{2}},
\]
(medias armónica, geométrica, aritmética y cuadrática) y que cada igualdad se da si y solo si $a=b$.
\end{ejercicio}

\begin{ejercicio}{3}{eje:2-bernoulli-fuerte}
Demuestre por inducción que, si $x\ge0$ y $n\in\N$, entonces $(1+x)^n\ge1+nx+\dfrac{n(n-1)}{2}x^2$. Deduzca que $2^n>\dfrac{n(n-1)}2$ para todo $n\in\N$.
\fuente{adaptado de Hammack~\cite{hammack}, cap.~10 (desigualdades por inducción).}
\end{ejercicio}

\begin{ejercicio}{3}{eje:2-suma-inversos}
Demuestre que si $x_1,\dots,x_n>0$, entonces $\displaystyle\Bigl(\sum_{i=1}^nx_i\Bigr)\Bigl(\sum_{i=1}^n\frac1{x_i}\Bigr)\ge n^2$, y que la igualdad se da si y solo si todos los $x_i$ son iguales.
\end{ejercicio}

% =====================================================================
\chapter{Polinomios, ecuaciones y sistemas}
\label{tema:3}
% =====================================================================

Los polinomios son las funciones más sencillas después de las lineales y aparecen por todas partes: como aproximaciones de Taylor, como funciones de pérdida cuadráticas y como polinomio característico de una matriz. Este tema repasa su aritmética (división, raíces, factorización), la resolución de ecuaciones e inecuaciones y los sistemas lineales pequeños. Sigue los capítulos~2, 5 y~7 de OpenStax \emph{College Algebra}~\cite{openstax}, con las demostraciones que allí se omiten.

\section{Polinomios y división euclídea}

\begin{definicion}[Polinomio]
Un \emph{polinomio} con coeficientes reales es una expresión $P(x)=a_nx^n+\dots+a_1x+a_0$ con $a_i\in\R$. Si $a_n\neq0$, $n$ es el \emph{grado} de $P$, $\deg P$, y $a_n$ su \emph{coeficiente principal}. Un número $r$ es una \emph{raíz} de $P$ si $P(r)=0$.
\end{definicion}

Si $P$ y $Q$ son no nulos, $\deg(PQ)=\deg P+\deg Q$, porque el coeficiente principal del producto es el producto de los coeficientes principales, que no es cero (\cref{prop:cuerpo}(iv)).

\begin{teorema}[División euclídea]\label{teo:division}
Sean $P$ y $D$ polinomios con $D\neq0$. Existen polinomios únicos $Q$ (cociente) y $R$ (resto) tales que
\[
P=DQ+R\qquad\text{y}\qquad R=0\ \text{o}\ \deg R<\deg D.
\]
\end{teorema}

\begin{proof}
\emph{Existencia}, por inducción sobre $\deg P$. Si $P=0$ o $\deg P<\deg D$, se toma $Q=0$ y $R=P$. Si $\deg P=n\ge m=\deg D$, con coeficientes principales $a_n$ y $d_m$, el polinomio $P_1=P-\frac{a_n}{d_m}x^{n-m}D$ tiene grado menor que $n$ (o es nulo), porque se ha cancelado el término de grado $n$. Por la hipótesis de inducción, $P_1=DQ_1+R$ con $R$ como se pide, y entonces $P=D\bigl(\frac{a_n}{d_m}x^{n-m}+Q_1\bigr)+R$.

\emph{Unicidad}. Si $DQ+R=DQ'+R'$, entonces $D(Q-Q')=R'-R$. Si $Q\neq Q'$, el miembro izquierdo tiene grado al menos $\deg D$ y el derecho grado menor que $\deg D$, lo cual es imposible. Luego $Q=Q'$ y, en consecuencia, $R=R'$.
\end{proof}

\begin{teorema}[Teoremas del resto y del factor]\label{teo:resto}
El resto de dividir $P$ entre $x-a$ es el número $P(a)$. En consecuencia, $a$ es raíz de $P$ si y solo si $x-a$ divide a $P$.
\end{teorema}

\begin{proof}
Por el \cref{teo:division}, $P(x)=(x-a)Q(x)+R$ con $R$ constante, porque $\deg R<1$. Evaluando en $x=a$, $P(a)=0\cdot Q(a)+R=R$.
\end{proof}

La \emph{regla de Ruffini} es una forma abreviada de dividir entre $x-a$: se bajan los coeficientes uno a uno, multiplicando el último resultado por $a$ y sumándolo al coeficiente siguiente.

\begin{ejemplo}\label{ej:ruffini}
Dividir $P(x)=2x^3-5x^2+0x+7$ entre $x-2$:
\[
\begin{array}{r|rrrr}
 & 2 & -5 & 0 & 7\\
2 & & 4 & -2 & -4\\\hline
 & 2 & -1 & -2 & 3
\end{array}
\]
Cociente $2x^2-x-2$ y resto $3$; en efecto, $P(2)=16-20+7=3$. El coeficiente nulo de $x$ debe escribirse.
\end{ejemplo}

\begin{corolario}\label{cor:raices}
Un polinomio no nulo de grado $n$ tiene a lo sumo $n$ raíces reales distintas.
\end{corolario}

\begin{proof}
Por inducción sobre $n$. Un polinomio de grado~0 es una constante no nula y no tiene raíces. Si $\deg P=n\ge1$ y $P$ tiene alguna raíz $a$, por el \cref{teo:resto} $P=(x-a)Q$ con $\deg Q=n-1$. Si $b\neq a$ es otra raíz, $0=(b-a)Q(b)$ obliga a $Q(b)=0$. Así, las raíces de $P$ son $a$ y las de $Q$, que por la hipótesis de inducción son a lo sumo $n-1$.
\end{proof}

\begin{teorema}[Raíces racionales]\label{teo:raices-racionales}
Sea $P(x)=a_nx^n+\dots+a_0$ con coeficientes enteros, $a_n\neq0$ y $a_0\neq0$. Si $p/q$, con $p,q$ enteros primos entre sí, es raíz de $P$, entonces $p$ divide a $a_0$ y $q$ divide a $a_n$.
\end{teorema}

\begin{proof}
Multiplicando $P(p/q)=0$ por $q^n$: $a_np^n+a_{n-1}p^{n-1}q+\dots+a_1pq^{n-1}+a_0q^n=0$. Todos los términos salvo el último son múltiplos de $p$, luego $p\mid a_0q^n$; como $p$ y $q$ no tienen factores primos comunes, $p\mid a_0$. Del mismo modo, todos los términos salvo el primero son múltiplos de $q$, luego $q\mid a_np^n$ y $q\mid a_n$.
\end{proof}

\section{La ecuación de segundo grado y las fórmulas de Vieta}

\begin{teorema}[Ecuación de segundo grado]\label{teo:cuadratica}
Sean $a\neq0$ y $\Delta=b^2-4ac$ el \emph{discriminante}. La ecuación $ax^2+bx+c=0$ tiene dos soluciones reales distintas si $\Delta>0$, una (doble) si $\Delta=0$ y ninguna si $\Delta<0$. Cuando $\Delta\ge0$,
\[
x=\frac{-b\pm\sqrt{\Delta}}{2a}.
\]
\end{teorema}

\begin{proof}
Se completa el cuadrado:
\[
ax^2+bx+c=a\Bigl[\Bigl(x+\frac{b}{2a}\Bigr)^2-\frac{\Delta}{4a^2}\Bigr].
\]
Como $a\neq0$, la ecuación equivale a $(x+\frac b{2a})^2=\frac{\Delta}{4a^2}$. El cuadrado de un real es no negativo, luego no hay solución si $\Delta<0$; si $\Delta\ge0$, extrayendo la raíz se obtiene $x+\frac b{2a}=\pm\frac{\sqrt\Delta}{2a}$.
\end{proof}

\begin{proposicion}[Fórmulas de Vieta]\label{prop:vieta}
Si $x_1,x_2$ son las raíces (reales o repetidas) de $ax^2+bx+c$, entonces
\[
ax^2+bx+c=a(x-x_1)(x-x_2),\qquad x_1+x_2=-\frac ba,\qquad x_1x_2=\frac ca.
\]
Más en general, si $P(x)=a_n\prod_{i=1}^n(x-x_i)$, entonces $\sum_ix_i=-a_{n-1}/a_n$ y $\prod_ix_i=(-1)^na_0/a_n$.
\end{proposicion}

\begin{proof}
Si $\Delta\ge0$, la expresión obtenida al completar el cuadrado en la demostración del \cref{teo:cuadratica} es una diferencia de cuadrados:
\[
ax^2+bx+c=a\Bigl(x+\frac b{2a}-\frac{\sqrt\Delta}{2a}\Bigr)\Bigl(x+\frac b{2a}+\frac{\sqrt\Delta}{2a}\Bigr)=a(x-x_1)(x-x_2).
\]
Desarrollando, $a(x-x_1)(x-x_2)=ax^2-a(x_1+x_2)x+ax_1x_2$, e igualando coeficientes se obtienen las fórmulas. El caso general es análogo: en $a_n\prod_i(x-x_i)$, el coeficiente de $x^{n-1}$ es $-a_n\sum_ix_i$ y el término independiente es $a_n\prod_i(-x_i)$.
\end{proof}

\begin{ejemplo}
Para $x^2-7x+10$: $x_1+x_2=7$ y $x_1x_2=10$, lo que sugiere $x_1=2$, $x_2=5$; en efecto, $x^2-7x+10=(x-2)(x-5)$.
\end{ejemplo}

\section{Otras ecuaciones e inecuaciones}

Las ecuaciones \emph{bicuadradas} $ax^4+bx^2+c=0$ se reducen a segundo grado con $t=x^2$. En las ecuaciones \emph{con radicales} se aísla la raíz y se eleva al cuadrado; como elevar al cuadrado no es una operación reversible ($a^2=b^2$ no implica $a=b$), pueden aparecer soluciones \emph{extrañas} y es obligatorio comprobar cada candidata en la ecuación original. En las ecuaciones \emph{racionales} se multiplica por el denominador común y se descartan los valores que lo anulan.

Para resolver una inecuación $P(x)/Q(x)\ge0$ se factorizan numerador y denominador, se marcan en la recta los puntos donde se anula algún factor y se estudia el signo de cada factor en cada intervalo (\emph{tabla de signos}); los ceros del denominador nunca forman parte de la solución.

\begin{ejemplo}\label{ej:tabla-signos}
Resolver $\dfrac{x-1}{x+2}<0$. Los puntos críticos son $-2$ y $1$. En $(-\infty,-2)$ ambos factores son negativos y el cociente es positivo; en $(-2,1)$ el numerador es negativo y el denominador positivo; en $(1,\infty)$ ambos son positivos. La solución es $(-2,1)$.
\end{ejemplo}

\section{Sistemas lineales pequeños}

Un sistema de $m$ ecuaciones lineales con $n$ incógnitas es \emph{compatible determinado} si tiene una única solución, \emph{compatible indeterminado} si tiene infinitas e \emph{incompatible} si no tiene ninguna. Los métodos de sustitución, igualación y reducción de bachillerato son casos particulares de la eliminación gaussiana del módulo~3.

\begin{teorema}[Regla de Cramer para $2\times2$]\label{teo:cramer}
El sistema $ax+by=e$, $cx+dy=f$ tiene solución única si y solo si $\delta=ad-bc\neq0$, y entonces
\[
x=\frac{ed-bf}{\delta},\qquad y=\frac{af-ec}{\delta}.
\]
\end{teorema}

\begin{proof}
Multiplicando la primera ecuación por $d$, la segunda por $b$ y restando: $(ad-bc)x=ed-bf$. Multiplicando la primera por $c$, la segunda por $a$ y restando: $(ad-bc)y=af-ec$. Si $\delta\neq0$, estas igualdades determinan $x$ e $y$, y se comprueba por sustitución que son solución. El caso $\delta=0$ se estudia en el \cref{eje:3-det-cero}.
\end{proof}

\begin{ejemplo}\label{ej:gauss3}
El sistema $x+y+z=6$, $2x-y+z=3$, $x+2y-z=2$ se resuelve por reducción. Restando a la segunda el doble de la primera y a la tercera la primera: $-3y-z=-9$ y $y-2z=-4$. De la última, $y=2z-4$; sustituyendo, $-3(2z-4)-z=-9$, es decir, $-7z=-21$ y $z=3$. Entonces $y=2$ y $x=1$.
\end{ejemplo}

\begin{observacion}[Errores frecuentes]
\begin{enumerate}[label=(\alph*)]
\item Dividir ambos miembros por $x$ o por $(x-a)$ pierde la solución $x=0$ o $x=a$; es preferible factorizar.
\item En las ecuaciones con radicales, no comprobar las soluciones produce soluciones extrañas (\cref{eje:3-ecuaciones}).
\item Multiplicar una inecuación por el denominador sin conocer su signo invierte, o no, el sentido; la tabla de signos evita el problema.
\item En la regla de Ruffini hay que escribir los coeficientes nulos.
\end{enumerate}
\end{observacion}

\begin{observacion}[Conexión con el curso]
Los autovalores de una matriz $2\times2$ son las raíces de su polinomio característico $\lambda^2-(\operatorname{tr}A)\lambda+\det A$; por las fórmulas de Vieta, su suma es la traza y su producto el determinante (módulo~3). Completar el cuadrado es la forma más rápida de minimizar una pérdida cuadrática en una variable, que es el modelo de juguete con el que se estudia la tasa de aprendizaje en el módulo~7. Las ecuaciones normales de la regresión lineal (módulo~8) son un sistema lineal.
\end{observacion}

\begin{lecturas}
\item OpenStax \emph{College Algebra}~\cite{openstax}, §1.4--1.5 (polinomios y factorización), §2.5 «Quadratic Equations», §2.6 «Other Types of Equations».
\item OpenStax~\cite{openstax}, §5.4 «Dividing Polynomials» y §5.5 «Zeros of Polynomial Functions» (teoremas del resto, del factor y de las raíces racionales).
\item OpenStax~\cite{openstax}, §7.1--7.2 (sistemas lineales de dos y tres incógnitas) y §7.8 «Solving Systems with Cramer's Rule».
\end{lecturas}

\section*{Ejercicios del tema 3}
\addcontentsline{toc}{section}{Ejercicios del tema 3}

\begin{ejercicio}{1}{eje:3-basico}
Sea $P(x)=2x^3-x+5$. (a) Calcule $P(-1)$. (b) ¿Qué grado tiene $P(x)^2\cdot(x^2+1)$? (c) Escriba una ecuación de segundo grado con coeficientes enteros cuyas raíces sean $3$ y $-\frac12$.
\end{ejercicio}

\begin{ejercicio}{1}{eje:3-discriminante}
Sin resolverlas, determine cuántas soluciones reales tiene cada ecuación: (a) $x^2-4x+5=0$; (b) $4x^2-12x+9=0$; (c) $2x^2+x-3=0$.
\end{ejercicio}

\begin{ejercicio}{2}{eje:3-division}
Divida $P(x)=2x^4-3x^3+x-5$ entre $D(x)=x^2-x+1$ y compruebe el resultado.
\fuente{adaptado de OpenStax \emph{College Algebra}~\cite{openstax}, §5.4.}
\end{ejercicio}

\begin{ejercicio}{2}{eje:3-ruffini}
Factorice con la regla de Ruffini: (a) $x^3-6x^2+11x-6$; (b) $x^4-5x^2+4$.
\fuente{adaptado de OpenStax \emph{College Algebra}~\cite{openstax}, §5.4--5.5.}
\end{ejercicio}

\begin{ejercicio}{2}{eje:3-resto}
(a) Halle $k$ para que $x^3+kx^2-4x+4$ sea divisible por $x-2$. (b) Halle, sin dividir, el resto de $x^{100}-1$ entre $x+1$ y el de $x^{101}+3$ entre $x+1$.
\end{ejercicio}

\begin{ejercicio}{2}{eje:3-racionales}
Halle todas las raíces de $2x^3-3x^2-3x+2$ usando el \cref{teo:raices-racionales}.
\fuente{adaptado de OpenStax \emph{College Algebra}~\cite{openstax}, §5.5 («Rational Zero Theorem»).}
\end{ejercicio}

\begin{ejercicio}{2}{eje:3-vieta}
Sean $x_1,x_2$ las raíces de $x^2-5x+3=0$. Sin calcularlas, halle $x_1^2+x_2^2$, $\dfrac1{x_1}+\dfrac1{x_2}$ y $x_1^3+x_2^3$.
\end{ejercicio}

\begin{ejercicio}{2}{eje:3-ecuaciones}
Resuelva: (a) $x^4-13x^2+36=0$; (b) $\sqrt{x+5}=x-1$; (c) $\dfrac1{x-1}+\dfrac2{x+1}=1$.
\fuente{adaptado de OpenStax \emph{College Algebra}~\cite{openstax}, §2.6.}
\end{ejercicio}

\begin{ejercicio}{2}{eje:3-inecuaciones}
Resuelva: (a) $x^2-x-6<0$; (b) $\dfrac{(x-1)(x+2)}{x-3}\ge0$.
\end{ejercicio}

\begin{ejercicio}{2}{eje:3-sistemas}
(a) Discuta y resuelva, según el valor de $k$, el sistema $x+2y=4$, $3x+ky=6$. (b) Resuelva $x+y+z=6$, $2x-y+z=3$, $x+2y-z=2$ por un orden de eliminación distinto del del \cref{ej:gauss3}.
\fuente{adaptado de OpenStax \emph{College Algebra}~\cite{openstax}, §7.1--7.2 y §7.8.}
\end{ejercicio}

\begin{ejercicio}{3}{eje:3-identidad}
Demuestre que si dos polinomios $P$ y $Q$ de grado a lo sumo $n$ toman el mismo valor en $n+1$ puntos distintos, entonces $P=Q$ (tienen los mismos coeficientes).
\end{ejercicio}

\begin{ejercicio}{3}{eje:3-diferencia-potencias}
Demuestre que, para todo $n\in\N$,
\[
x^n-a^n=(x-a)(x^{n-1}+ax^{n-2}+\dots+a^{n-2}x+a^{n-1}).
\]
Deduzca que $x-a$ divide a $x^n-a^n$ y que, si $n$ es impar, $x+a$ divide a $x^n+a^n$.
\end{ejercicio}

\begin{ejercicio}{3}{eje:3-impares}
Sean $b$ y $c$ enteros impares. Demuestre que la ecuación $x^2+bx+c=0$ no tiene soluciones racionales.
\end{ejercicio}

\begin{ejercicio}{3}{eje:3-det-cero}
Demuestre que si $ad-bc=0$, el sistema $ax+by=e$, $cx+dy=f$ o bien no tiene solución o bien tiene infinitas. (Suponga que algún coeficiente $a,b,c,d$ es no nulo.)
\end{ejercicio}

\begin{ejercicio}{3}{eje:3-cuadratica-positiva}
Sea $a>0$ y $f(x)=ax^2+bx+c$. Demuestre que $f(x)\ge c-\dfrac{b^2}{4a}$ para todo $x$, con igualdad solo en $x=-\dfrac{b}{2a}$, y deduzca que $f(x)\ge0$ para todo $x\in\R$ si y solo si $b^2-4ac\le0$.
\end{ejercicio}

% =====================================================================
\chapter{Sumatorios, productorios y binomio de Newton}
\label{tema:4}
% =====================================================================

La probabilidad de una secuencia es un producto de probabilidades condicionadas, su logaritmo es una suma, y la pérdida de entrenamiento es una media sobre millones de tokens. Manejar con soltura la notación $\sum$ y $\prod$ es, por tanto, imprescindible. Este tema sigue el capítulo~2 «Sums» de Graham, Knuth y Patashnik~\cite{gkp} y el capítulo~14 «Sums and Asymptotics» de Lehman, Leighton y Meyer~\cite{mcs}; el binomio de Newton sigue el §9.6 de OpenStax~\cite{openstax}.

\section{Notación y principio de inducción}

\begin{definicion}[Sumatorio y productorio]
Dados enteros $m\le n$ y números $a_m,\dots,a_n$,
\[
\sum_{k=m}^{n}a_k=a_m+a_{m+1}+\dots+a_n,\qquad \prod_{k=m}^{n}a_k=a_m\,a_{m+1}\cdots a_n.
\]
La letra $k$ es un \emph{índice mudo}: puede sustituirse por cualquier otra sin cambiar el valor. Si $n<m$, la suma es \emph{vacía} y vale~0, y el producto vacío vale~1. Más en general, $\sum_{k\in S}a_k$ denota la suma sobre un conjunto finito $S$ de índices. El \emph{factorial} es $n!=\prod_{k=1}^{n}k$, con $0!=1$.
\end{definicion}

Formalmente, la suma se define por recurrencia: $\sum_{k=m}^{m-1}a_k=0$ y $\sum_{k=m}^{n}a_k=\sum_{k=m}^{n-1}a_k+a_n$. Por eso las fórmulas para sumas se demuestran de forma natural por inducción.

\begin{teorema}[Principio de inducción]\label{teo:induccion}
Sea $P(n)$ una afirmación sobre los enteros $n\ge n_0$. Si $P(n_0)$ es cierta y, para cada $n\ge n_0$, $P(n)$ implica $P(n+1)$, entonces $P(n)$ es cierta para todo $n\ge n_0$.
\end{teorema}

Este principio es equivalente al principio del buen orden (todo subconjunto no vacío de $\N$ tiene mínimo); ambos se estudian en~\cite[cap.~2 y~5]{mcs} y~\cite[cap.~10]{hammack}.

\section{Propiedades de los sumatorios}

\begin{proposicion}\label{prop:sumas}
Para números $a_k,b_k$ y una constante $c$:
\begin{enumerate}
  \item linealidad: $\sum_{k=m}^n(ca_k+b_k)=c\sum_{k=m}^na_k+\sum_{k=m}^nb_k$;
  \item partición: si $m\le p<n$, $\sum_{k=m}^na_k=\sum_{k=m}^pa_k+\sum_{k=p+1}^na_k$;
  \item suma de una constante: $\sum_{k=m}^nc=(n-m+1)\,c$;
  \item cambio de índice: $\sum_{k=m}^na_k=\sum_{j=m+r}^{n+r}a_{j-r}$ para todo entero $r$, y $\sum_{k=0}^na_k=\sum_{k=0}^na_{n-k}$.
\end{enumerate}
\end{proposicion}

\begin{proof}
(i)--(iii) son consecuencia de la asociatividad, la conmutatividad y la distributividad de los axiomas de cuerpo, aplicadas un número finito de veces (formalmente, por inducción sobre el número de términos). (iv) En ambos casos se suman exactamente los mismos números en distinto orden: con $j=k+r$, cuando $k$ recorre $m,\dots,n$, $j$ recorre $m+r,\dots,n+r$; con $j=n-k$, cuando $k$ recorre $0,\dots,n$, $j$ recorre $n,\dots,0$.
\end{proof}

No hay reglas análogas para el producto: en general $\sum a_kb_k\neq\bigl(\sum a_k\bigr)\bigl(\sum b_k\bigr)$.

\section{Sumas clásicas}

\begin{teorema}[Suma aritmética]\label{teo:aritmetica}
Para todo $n\in\N$, $\displaystyle\sum_{k=1}^nk=\frac{n(n+1)}2$. En general, si $a_k=a_1+(k-1)d$, entonces $\sum_{k=1}^na_k=\dfrac{n(a_1+a_n)}{2}$.
\end{teorema}

\begin{proof}
Sea $S=\sum_{k=1}^nk$. Por la \cref{prop:sumas}(iv), $S=\sum_{k=1}^n(n+1-k)$. Sumando ambas expresiones, $2S=\sum_{k=1}^n(n+1)=n(n+1)$. La fórmula general se obtiene igual, porque $a_k+a_{n+1-k}=a_1+a_n$ para todo $k$.
\end{proof}

\begin{teorema}[Suma geométrica]\label{teo:geometrica}
Para $r\neq1$ y $n\in\N_0$, $\displaystyle\sum_{k=0}^nr^k=\frac{1-r^{n+1}}{1-r}$. Si $r=1$, la suma vale $n+1$.
\end{teorema}

\begin{proof}
Por el \emph{método de perturbación}~\cite[§2.3]{gkp}: sea $S=\sum_{k=0}^nr^k$. Entonces $rS=\sum_{k=1}^{n+1}r^k$ y, restando, se cancelan todos los términos salvo dos: $S-rS=1-r^{n+1}$. Como $r\neq1$, se divide por $1-r$.
\end{proof}

\begin{proposicion}[Sumas telescópicas]\label{prop:telescopica}
Para cualesquiera números $a_m,\dots,a_{n+1}$,
\[
\sum_{k=m}^n(a_{k+1}-a_k)=a_{n+1}-a_m.
\]
\end{proposicion}

\begin{proof}
Por linealidad y cambio de índice, $\sum_{k=m}^na_{k+1}-\sum_{k=m}^na_k=\sum_{j=m+1}^{n+1}a_j-\sum_{k=m}^na_k=a_{n+1}-a_m$, ya que los términos con índice entre $m+1$ y $n$ aparecen en ambas sumas.
\end{proof}

\begin{ejemplo}\label{ej:telescopica}
Como $\frac1{k(k+1)}=\frac1k-\frac1{k+1}$, se tiene $\displaystyle\sum_{k=1}^n\frac1{k(k+1)}=1-\frac1{n+1}$.
\end{ejemplo}

\begin{teorema}[Sumas de potencias]\label{teo:potencias}
Para todo $n\in\N$,
\[
\sum_{k=1}^nk^2=\frac{n(n+1)(2n+1)}6,\qquad\sum_{k=1}^nk^3=\Bigl(\frac{n(n+1)}2\Bigr)^2.
\]
\end{teorema}

\begin{proof}
Para la primera, se suma la identidad $(k+1)^3-k^3=3k^2+3k+1$ para $k=1,\dots,n$. El miembro izquierdo es telescópico y vale $(n+1)^3-1$; el derecho vale $3S_2+3\frac{n(n+1)}2+n$, donde $S_2=\sum k^2$. Despejando,
\[
3S_2=(n+1)^3-1-\tfrac32n(n+1)-n=\tfrac{(n+1)}{2}\bigl[2(n+1)^2-3n-2\bigr]=\tfrac{(n+1)n(2n+1)}{2}.
\]
La segunda fórmula se demuestra por inducción en el \cref{eje:4-cubos}.
\end{proof}

\begin{ejemplo}[Medias exponenciales]\label{ej:ema}
El optimizador Adam (módulo~7) mantiene la media móvil $m_t=\beta m_{t-1}+(1-\beta)g_t$ con $m_0=0$ y $0<\beta<1$. Desplegando la recurrencia, $m_t=(1-\beta)\sum_{k=1}^{t}\beta^{t-k}g_k$. Si todos los gradientes valen $g$, el \cref{teo:geometrica} da $m_t=(1-\beta)g\sum_{j=0}^{t-1}\beta^{j}=(1-\beta^t)\,g$: la media está sesgada hacia cero. Con $\beta=0{,}9$ y $t=10$, $m_{10}\approx0{,}651\,g$. Adam divide por $1-\beta^t$ para corregir ese sesgo.
\end{ejemplo}

\section{Sumas dobles}

Una suma doble $\sum_{i}\sum_{j}a_{ij}$ suma los elementos de una tabla. Si los índices recorren un rectángulo, el orden de suma puede intercambiarse:
\[
\sum_{i=1}^m\sum_{j=1}^na_{ij}=\sum_{j=1}^n\sum_{i=1}^ma_{ij}.
\]
Si recorren un triángulo, hay que reescribir los límites~\cite[§2.4]{gkp}:
\[
\sum_{i=1}^n\sum_{j=i}^na_{ij}=\sum_{1\le i\le j\le n}a_{ij}=\sum_{j=1}^n\sum_{i=1}^ja_{ij}.
\]

\begin{proposicion}\label{prop:producto-sumas}
$\displaystyle\Bigl(\sum_{i=1}^ma_i\Bigr)\Bigl(\sum_{j=1}^nb_j\Bigr)=\sum_{i=1}^m\sum_{j=1}^na_ib_j$. En particular, $\displaystyle\Bigl(\sum_{i=1}^na_i\Bigr)^2=\sum_{i=1}^na_i^2+2\sum_{1\le i<j\le n}a_ia_j$.
\end{proposicion}

\begin{proof}
Por linealidad, $\bigl(\sum_ia_i\bigr)B=\sum_ia_iB$ con $B=\sum_jb_j$, y de nuevo $a_iB=\sum_ja_ib_j$. Para la segunda, con $b=a$ se separan en la suma doble los términos con $i=j$, los de $i<j$ y los de $i>j$; los dos últimos grupos son iguales por simetría.
\end{proof}

\begin{ejemplo}[Coste de la atención causal]\label{ej:causal}
En un Transformer con máscara causal (módulo~13), el token $i$ atiende a los tokens $j\le i$. El número de pares $(i,j)$ con $1\le j\le i\le n$ es $\sum_{i=1}^n\sum_{j=1}^i1=\sum_{i=1}^ni=\frac{n(n+1)}2$, aproximadamente la mitad de los $n^2$ pares de la atención completa.
\end{ejemplo}

\section{Productorios y binomio de Newton}

Los productorios cumplen $\prod_k(a_kb_k)=\prod_ka_k\prod_kb_k$, $\prod_{k=1}^nc=c^n$ y, si todos los $a_k$ son positivos, $\ln\prod_ka_k=\sum_k\ln a_k$ (\cref{tema:5}).

\begin{definicion}[Números combinatorios]
Para enteros $0\le k\le n$, $\displaystyle\binom nk=\frac{n!}{k!\,(n-k)!}$, y $\binom nk=0$ si $k<0$ o $k>n$.
\end{definicion}

\begin{lema}[Identidad de Pascal]\label{lem:pascal}
Para $1\le k\le n$, $\displaystyle\binom nk+\binom n{k-1}=\binom{n+1}k$.
\end{lema}

\begin{proof}
Con denominador común $k!\,(n+1-k)!$:
\[
\frac{n!}{k!(n-k)!}+\frac{n!}{(k-1)!(n-k+1)!}=\frac{n!\,\bigl[(n+1-k)+k\bigr]}{k!\,(n+1-k)!}=\frac{(n+1)!}{k!\,(n+1-k)!}.\qedhere
\]
\end{proof}

\begin{teorema}[Binomio de Newton]\label{teo:binomio}
Para todos $a,b\in\R$ y $n\in\N_0$, $\displaystyle(a+b)^n=\sum_{k=0}^n\binom nka^{n-k}b^k$.
\end{teorema}

\begin{proof}
Por inducción sobre $n$. Para $n=0$ ambos miembros valen~1. Supóngase cierto para $n$. Entonces
\[
(a+b)^{n+1}=(a+b)\sum_{k=0}^n\binom nka^{n-k}b^k=\sum_{k=0}^n\binom nka^{n+1-k}b^k+\sum_{k=1}^{n+1}\binom n{k-1}a^{n+1-k}b^k,
\]
donde en la segunda suma se ha cambiado $k$ por $k-1$. Usando el convenio $\binom n{-1}=\binom n{n+1}=0$, ambas sumas pueden extenderse a $k=0,\dots,n+1$, y por el \cref{lem:pascal} el coeficiente de $a^{n+1-k}b^k$ es $\binom nk+\binom n{k-1}=\binom{n+1}k$ (para $k=0$ y $k=n+1$ vale~1, que también coincide).
\end{proof}

\begin{corolario}\label{cor:binomio}
$\displaystyle\sum_{k=0}^n\binom nk=2^n$ y, si $n\ge1$, $\displaystyle\sum_{k=0}^n(-1)^k\binom nk=0$.
\end{corolario}

\begin{proof}
Basta tomar $a=b=1$ y $a=1$, $b=-1$ en el \cref{teo:binomio}.
\end{proof}

\begin{observacion}[Errores frecuentes]
\begin{enumerate}[label=(\alph*)]
\item $\sum_{k=1}^n c=nc$, no $c$.
\item $\sum a_kb_k\neq\sum a_k\sum b_k$ y $\bigl(\sum a_k\bigr)^2\neq\sum a_k^2$ (\cref{prop:producto-sumas}).
\item Errores de uno en los límites: $\sum_{k=m}^n$ tiene $n-m+1$ términos, no $n-m$.
\item La fórmula geométrica exige $r\neq1$; al cambiar el índice hay que cambiar a la vez los límites y el término general.
\end{enumerate}
\end{observacion}

\begin{observacion}[Conexión con el curso]
Por la regla de la cadena (módulo~1), la probabilidad de una secuencia es $\prod_tP(x_t\mid x_{<t})$ y su logaritmo es $\sum_t\ln P(x_t\mid x_{<t})$: los productorios se convierten en sumatorios. La serie geométrica reaparece en el retorno descontado $\sum_t\gamma^tr_t$ del aprendizaje por refuerzo (módulo~21), y los números combinatorios en el estimador pass@$k$ del módulo~22.
\end{observacion}

\begin{lecturas}
\item Graham, Knuth y Patashnik~\cite{gkp}, §2.1 «Notation», §2.3 «Manipulation of Sums» (método de perturbación) y §2.4 «Multiple Sums».
\item Lehman, Leighton y Meyer~\cite{mcs}, §14.1 «The Value of an Annuity», §14.2 «Sums of Powers», §14.5 «Products» y §14.6 «Double Trouble».
\item OpenStax \emph{College Algebra}~\cite{openstax}, §9.2--9.4 (sucesiones aritméticas y geométricas, series) y §9.6 «Binomial Theorem».
\item Hammack~\cite{hammack}, cap.~10 «Mathematical Induction».
\end{lecturas}

\section*{Ejercicios del tema 4}
\addcontentsline{toc}{section}{Ejercicios del tema 4}

\begin{ejercicio}{1}{eje:4-notacion}
Escriba con notación de sumatorio: (a) $1+3+5+\dots+99$; (b) $2+4+8+\dots+1024$; (c) $\frac12-\frac13+\frac14-\dots+\frac1{100}$.
\end{ejercicio}

\begin{ejercicio}{1}{eje:4-calcular}
Calcule: (a) $\sum_{k=1}^4k^2$; (b) $\prod_{k=1}^5\frac k{k+1}$; (c) $\sum_{k=4}^{0}q_k$ con el convenio de este tema; (d) $\sum_{k=1}^{10}7$.
\fuente{adaptado de Graham, Knuth y Patashnik~\cite{gkp}, cap.~2, ejercicios de calentamiento (convenio sobre sumas vacías).}
\end{ejercicio}

\begin{ejercicio}{1}{eje:4-errores}
Encuentre el error en cada «igualdad»: (a) $\sum_{k=1}^n3=3$; (b) $\sum_{k=1}^2a_kb_k=(a_1+a_2)(b_1+b_2)$; (c) $\bigl(\sum_{k=1}^2a_k\bigr)^2=a_1^2+a_2^2$. Dé un contraejemplo numérico de (b).
\end{ejercicio}

\begin{ejercicio}{2}{eje:4-aritmeticas}
(a) Calcule $\sum_{k=1}^{100}(3k-2)$. (b) Calcule la suma de todos los múltiplos de~7 comprendidos entre 100 y 1000.
\fuente{adaptado de OpenStax \emph{College Algebra}~\cite{openstax}, §9.2 y §9.4.}
\end{ejercicio}

\begin{ejercicio}{2}{eje:4-geometricas}
Calcule: (a) $\sum_{k=0}^{9}3\cdot2^k$; (b) $\sum_{k=1}^{n}\bigl(\frac12\bigr)^k$; (c) $\sum_{k=0}^{n}(-1)^kx^k$ para $x\neq-1$.
\fuente{adaptado de OpenStax \emph{College Algebra}~\cite{openstax}, §9.3--9.4.}
\end{ejercicio}

\begin{ejercicio}{2}{eje:4-telescopicas}
Calcule en forma cerrada: (a) $\sum_{k=1}^n\frac1{k(k+2)}$ y su valor para $n=10$; (b) $\sum_{k=1}^n\ln\bigl(1+\frac1k\bigr)$; (c) $\sum_{k=1}^{99}\frac1{\sqrt k+\sqrt{k+1}}$.
\end{ejercicio}

\begin{ejercicio}{2}{eje:4-indice}
(a) Escriba $\sum_{k=3}^{n+2}(k-2)^2$ como una suma que empiece en~1 y calcúlela. (b) Calcule $\sum_{k=1}^nk(k+1)$ en forma cerrada usando las fórmulas del tema y compruébela para $n=10$.
\end{ejercicio}

\begin{ejercicio}{2}{eje:4-dobles}
Calcule: (a) $\sum_{i=1}^3\sum_{j=1}^4(i+2j)$; (b) $\sum_{i=1}^n\sum_{j=1}^ij$; (c) intercambie el orden de suma en $\sum_{i=1}^n\sum_{j=i}^n\frac1j$ y calcule su valor.
\fuente{adaptado de Lehman, Leighton y Meyer~\cite{mcs}, §14.6 «Double Trouble».}
\end{ejercicio}

\begin{ejercicio}{2}{eje:4-binomio}
(a) Desarrolle $(2x-1)^5$. (b) Halle el coeficiente de $x^4$ en $\bigl(x+\frac2x\bigr)^8$. (c) Calcule $\sum_{k=0}^{6}\binom6k2^k$ sin sumar término a término.
\fuente{adaptado de OpenStax \emph{College Algebra}~\cite{openstax}, §9.6.}
\end{ejercicio}

\begin{ejercicio}{2}{eje:4-productos}
Calcule: (a) $\prod_{k=2}^n\bigl(1-\frac1{k^2}\bigr)$ y su valor para $n=10$; (b) $\prod_{k=1}^n2^k$; (c) $\ln\prod_{k=1}^{n}\frac{k+1}{k}$.
\fuente{adaptado de Lehman, Leighton y Meyer~\cite{mcs}, §14.5 «Products».}
\end{ejercicio}

\begin{ejercicio}{3}{eje:4-cubos}
Demuestre por inducción que $\displaystyle\sum_{k=1}^nk^3=\Bigl(\frac{n(n+1)}2\Bigr)^2$ para todo $n\in\N$.
\fuente{adaptado de Lehman, Leighton y Meyer~\cite{mcs}, problema~2.17 (allí se pide por el principio del buen orden).}
\end{ejercicio}

\begin{ejercicio}{3}{eje:4-perturbacion}
Use el método de perturbación para demostrar que, si $r\neq1$,
\[
\sum_{k=0}^nkr^k=\frac{r-(n+1)r^{n+1}+nr^{n+2}}{(1-r)^2}.
\]
\fuente{adaptado de Graham, Knuth y Patashnik~\cite{gkp}, §2.3.}
\end{ejercicio}

\begin{ejercicio}{3}{eje:4-kbinom}
Demuestre que $k\binom nk=n\binom{n-1}{k-1}$ para $1\le k\le n$ y deduzca que $\sum_{k=0}^nk\binom nk=n\,2^{n-1}$.
\end{ejercicio}

\begin{ejercicio}{3}{eje:4-basilea}
Demuestre por inducción que $\displaystyle\sum_{k=1}^n\frac1{k^2}\le2-\frac1n$ para todo $n\in\N$.
\fuente{adaptado de Hammack~\cite{hammack}, cap.~10, ejercicios.}
\end{ejercicio}

\begin{ejercicio}{3}{eje:4-varianza}
Sean $a_1,\dots,a_n\in\R$. Demuestre que
\[
\sum_{i=1}^n\sum_{j=1}^n(a_i-a_j)^2=2n\sum_{i=1}^na_i^2-2\Bigl(\sum_{i=1}^na_i\Bigr)^2
\]
y deduzca que $\bigl(\sum_ia_i\bigr)^2\le n\sum_ia_i^2$, con igualdad si y solo si todos los $a_i$ son iguales.
\end{ejercicio}

% =====================================================================
\chapter{Exponenciales y logaritmos}
\label{tema:5}
% =====================================================================

El logaritmo convierte productos en sumas y potencias en productos. Esa propiedad lo hace omnipresente en el curso: la log-verosimilitud, la entropía, la perplejidad y la escala de las leyes de potencia son logaritmos. Este tema repasa las propiedades de exponenciales y logaritmos siguiendo el capítulo~6 de OpenStax~\cite{openstax}, introduce la función log-suma-exp siguiendo a Blanchard, Higham y Higham~\cite{blanchard} y fija las unidades de información (bits y nats) como en el capítulo~2 de Cover y Thomas~\cite{ct}. La construcción rigurosa de $e^x$ y $\ln x$ requiere el cálculo integral~\cite[cap.~18]{spivak}; aquí se admiten sus propiedades básicas.

\section{Potencias de exponente real y la función exponencial}

Para $a>0$, la potencia $a^r$ está definida para $r$ racional (\cref{tema:1}) y se extiende a todo $x\in\R$ de forma que la función $x\mapsto a^x$ es continua. Se admiten las siguientes propiedades, que extienden la \cref{prop:exponentes}.

\begin{proposicion}[Función exponencial]\label{prop:exp}
Sea $a>0$. Para todos $x,y\in\R$: $a^x>0$; $a^{x+y}=a^xa^y$; $(a^x)^y=a^{xy}$; $a^{-x}=1/a^x$. Si $a>1$, la función $x\mapsto a^x$ es estrictamente creciente; si $0<a<1$, es estrictamente decreciente. Si $a\neq1$, toma cada valor positivo exactamente una vez.
\end{proposicion}

El número $e=\lim_{n\to\infty}(1+\frac1n)^n\approx2{,}71828$ es la base natural; más en general, $e^x=\lim_{n\to\infty}(1+\frac xn)^n$ para todo $x$~\cite[cap.~18]{spivak}.

\begin{proposicion}\label{prop:ex-cota}
Para todo $x\in\R$, $e^x\ge1+x$.
\end{proposicion}

\begin{proof}
Si $x\le-1$, el miembro derecho es $\le0<e^x$. Si $x>-1$, para todo $n\in\N$ se tiene $x/n>-1$, y la desigualdad de Bernoulli (\cref{teo:bernoulli}) da $(1+\frac xn)^n\ge1+x$. Pasando al límite, que conserva las desigualdades no estrictas, $e^x\ge1+x$.
\end{proof}

\section{Logaritmos}

\begin{definicion}[Logaritmo]\label{def:log}
Sean $a>0$, $a\neq1$ y $x>0$. El \emph{logaritmo en base $a$} de $x$ es el único $y\in\R$ tal que $a^y=x$; se escribe $y=\log_ax$. Se usan $\ln=\log_e$ (logaritmo natural), $\log_2$ (binario) y $\log_{10}$ (decimal).
\end{definicion}

Por definición, $a^{\log_ax}=x$ para $x>0$ y $\log_a(a^y)=y$ para todo $y$. El logaritmo solo está definido para argumentos positivos.

\begin{teorema}[Propiedades de los logaritmos]\label{teo:log}
Sean $a,b>0$ distintos de~1, $x,y>0$ y $r\in\R$. Entonces:
\begin{enumerate}
  \item $\log_a(xy)=\log_ax+\log_ay$ y $\log_a(x/y)=\log_ax-\log_ay$;
  \item $\log_a(x^r)=r\log_ax$;
  \item cambio de base: $\log_ax=\dfrac{\log_bx}{\log_ba}$; en particular, $\log_ab\cdot\log_ba=1$;
  \item si $a>1$, $x<y$ si y solo si $\log_ax<\log_ay$.
\end{enumerate}
\end{teorema}

\begin{proof}
Sean $u=\log_ax$ y $v=\log_ay$, es decir, $a^u=x$ y $a^v=y$. (i) $a^{u+v}=a^ua^v=xy$, luego $\log_a(xy)=u+v$; el cociente es análogo con $a^{u-v}=x/y$. (ii) $a^{ru}=(a^u)^r=x^r$, luego $\log_a(x^r)=ru$. (iii) Tomando $\log_b$ en $a^u=x$ y usando (ii): $u\log_ba=\log_bx$; como $a\neq1$, $\log_ba\neq0$ y se despeja $u$. Con $x=b$ se obtiene $\log_ab=1/\log_ba$. (iv) Por la \cref{prop:exp}, $t\mapsto a^t$ es estrictamente creciente si $a>1$, así que $u<v$ si y solo si $a^u<a^v$.
\end{proof}

\begin{corolario}\label{cor:lnx}
Para todo $x>0$, $\ln x\le x-1$, con igualdad si y solo si $x=1$. En consecuencia, $\ln x\ge1-\frac1x$.
\end{corolario}

\begin{proof}
Por la \cref{prop:ex-cota} con $\ln x$ en lugar de $x$: $x=e^{\ln x}\ge1+\ln x$. Aplicándolo a $1/x$: $-\ln x=\ln\frac1x\le\frac1x-1$. (La igualdad solo en $x=1$ se sigue de que $e^t>1+t$ para $t\neq0$, que se demuestra en el módulo de cálculo.)
\end{proof}

\section{Ecuaciones exponenciales y logarítmicas}

Para resolver una ecuación exponencial se toman logaritmos o se cambia la incógnita ($t=a^x$); en las logarítmicas se usan las propiedades del \cref{teo:log} para llegar a $\log_aA=\log_aB$, que equivale a $A=B$, y se descartan las soluciones que hacen negativo o nulo algún argumento.

\begin{ejemplo}\label{ej:ecuacion-exp}
$9^x-4\cdot3^x+3=0$. Con $t=3^x>0$, queda $t^2-4t+3=0$, de raíces $t=1$ y $t=3$, que dan $x=0$ y $x=1$.
\end{ejemplo}

\begin{ejemplo}\label{ej:ecuacion-log}
$\log_3(x+1)+\log_3(x-1)=1$. Exige $x>1$. Por el \cref{teo:log}, $\log_3(x^2-1)=1$, es decir, $x^2-1=3$ y $x=\pm2$. Solo $x=2$ cumple $x>1$.
\end{ejemplo}

\section{La función log-suma-exp}

\begin{definicion}[Log-suma-exp]\label{def:lse}
Para $z=(z_1,\dots,z_n)\in\R^n$, $\displaystyle\LSE(z)=\ln\sum_{i=1}^ne^{z_i}$. La función \emph{softmax} es $\softmax(z)_i=e^{z_i}/\sum_je^{z_j}=e^{z_i-\LSE(z)}$.
\end{definicion}

\begin{proposicion}\label{prop:lse}
Para todo $z\in\R^n$ y $c\in\R$, con $M=\max_iz_i$:
\begin{enumerate}
  \item invariancia por traslación: $\LSE(z_1+c,\dots,z_n+c)=\LSE(z)+c$; en consecuencia, la softmax no cambia al sumar $c$ a todas las componentes;
  \item cotas: $M\le\LSE(z)\le M+\ln n$.
\end{enumerate}
\end{proposicion}

\begin{proof}
(i) $\ln\sum_ie^{z_i+c}=\ln\bigl(e^c\sum_ie^{z_i}\bigr)=c+\LSE(z)$ por el \cref{teo:log}(i). Para la softmax, el factor $e^c$ aparece en el numerador y en el denominador. (ii) La suma $\sum_ie^{z_i}$ contiene el término $e^M$ y todos sus términos son positivos y menores o iguales que $e^M$, así que $e^M\le\sum_ie^{z_i}\le ne^M$. Tomando $\ln$, que es creciente, se obtienen las cotas.
\end{proof}

La \cref{prop:lse}(i) con $c=-M$ da el \emph{cálculo estable}: $\LSE(z)=M+\ln\sum_ie^{z_i-M}$. Todas las exponenciales son ahora $\le1$, de modo que no hay desbordamiento, y al menos una vale exactamente~1, de modo que el logaritmo no recibe un cero.

\begin{ejemplo}\label{ej:lse-estable}
Con $z=(1000,1000)$, el cálculo directo de $e^{1000}$ desborda en doble precisión (el mayor número representable es del orden de $1{,}8\cdot10^{308}$, mientras que $e^{1000}\approx10^{434}$). El cálculo estable da $\LSE(z)=1000+\ln(1+1)=1000+\ln2\approx1000{,}693$, y la softmax es $(\frac12,\frac12)$.
\end{ejemplo}

\section{Escalas logarítmicas, bits y nats}

Una cantidad que varía en muchos órdenes de magnitud se representa mejor por su logaritmo. Si $y=Cx^{k}$, entonces $\log y=\log C+k\log x$: en ejes doblemente logarítmicos una \emph{ley de potencia} es una recta de pendiente $k$, que es como se presentan las leyes de escala del módulo~16. Si $y=Ca^{x}$, entonces $\log y$ es lineal en $x$ (escala semilogarítmica).

\begin{definicion}[Unidades de información]
La \emph{información} de un suceso de probabilidad $p\in(0,1]$ es $-\log_2p$ \emph{bits} o $-\ln p$ \emph{nats}~\cite[cap.~2]{ct}. Por el cambio de base, $1\ \text{nat}=\frac1{\ln2}\ \text{bits}\approx1{,}4427$ bits.
\end{definicion}

\begin{ejemplo}[Productos de probabilidades]\label{ej:underflow}
Un texto de 500 tokens en el que el modelo asigna probabilidad $10^{-3}$ a cada token tiene probabilidad $10^{-1500}$, muy por debajo del menor número positivo representable en doble precisión (del orden de $10^{-308}$, o $5\cdot10^{-324}$ con números subnormales): el producto se redondea a~0. Su logaritmo, en cambio, es $500\ln10^{-3}\approx-3453{,}9$ nats, perfectamente representable. Por eso los modelos trabajan con log-probabilidades y la regla de la cadena se aplica como suma de logaritmos.
\end{ejemplo}

\begin{observacion}[Errores frecuentes]
\begin{enumerate}[label=(\alph*)]
\item $\log(x+y)\neq\log x+\log y$; la propiedad es $\log(xy)=\log x+\log y$.
\item $(\ln x)^2\neq\ln(x^2)=2\ln x$, y $\ln x/\ln y\neq\ln(x/y)$.
\item $a^{x+y}=a^xa^y$, no $a^x+a^y$.
\item Al resolver una ecuación logarítmica hay que comprobar que los argumentos son positivos (\cref{ej:ecuacion-log}).
\item Confundir $\log$ (que en muchos libros de aprendizaje automático y en Python significa $\ln$) con $\log_{10}$.
\end{enumerate}
\end{observacion}

\begin{observacion}[Conexión con el curso]
La log-verosimilitud, la entropía cruzada y la perplejidad del módulo~1 son sumas de logaritmos; el cambio de base convierte nats en bits (módulo~6). La función log-suma-exp es el denominador de la softmax en logaritmos: aparece en la entropía cruzada (módulo~8) y en la atención (módulo~13), y las cotas de la \cref{prop:lse} explican por qué la softmax con temperatura tiende al máximo cuando la temperatura tiende a cero (\cref{eje:5-temperatura}). La desigualdad $\ln x\le x-1$ es el paso clave de la desigualdad de Gibbs (módulo~6).
\end{observacion}

\begin{lecturas}
\item OpenStax \emph{College Algebra}~\cite{openstax}, §6.1 «Exponential Functions», §6.3 «Logarithmic Functions», §6.5 «Logarithmic Properties», §6.6 «Exponential and Logarithmic Equations» y §6.7 «Exponential and Logarithmic Models».
\item Spivak~\cite{spivak}, cap.~18 «The Logarithm and Exponential Functions», para la construcción rigurosa.
\item Blanchard, Higham y Higham~\cite{blanchard}, secciones introductorias, sobre el cálculo estable de log-suma-exp y softmax.
\item Cover y Thomas~\cite{ct}, cap.~2, para bits, nats y la desigualdad de la información.
\end{lecturas}

\section*{Ejercicios del tema 5}
\addcontentsline{toc}{section}{Ejercicios del tema 5}

\begin{ejercicio}{1}{eje:5-calcular}
Calcule sin calculadora: $\log_264$, $\log_3\frac1{81}$, $\log_{10}0{,}001$, $\ln e^7$, $\log_48$ y $10^{\log_{10}5}$.
\fuente{adaptado de OpenStax \emph{College Algebra}~\cite{openstax}, §6.3.}
\end{ejercicio}

\begin{ejercicio}{1}{eje:5-vf}
Razone si son verdaderas o falsas para $a,b,x>0$: (a) $\ln(a+b)=\ln a+\ln b$; (b) $\log_2(8^x)=3x$; (c) $\ln\frac1x=-\ln x$; (d) $(\ln x)^2=2\ln x$.
\end{ejercicio}

\begin{ejercicio}{1}{eje:5-unidades}
(a) ¿Cuántos bits son 8 nats? (b) ¿Cuánta información, en bits y en nats, tiene un suceso de probabilidad $1/1024$? (c) ¿Cuántos bits hacen falta como mínimo para numerar los tokens de un vocabulario de $50\,257$ tokens?
\end{ejercicio}

\begin{ejercicio}{2}{eje:5-simplificar}
Simplifique: (a) $\log_212-\log_23+\log_2\frac18$; (b) $\frac12\ln16+\ln\frac e4$; (c) $\log_23\cdot\log_38$; (d) $\log_530$ expresado con logaritmos naturales, y su valor aproximado.
\fuente{adaptado de OpenStax \emph{College Algebra}~\cite{openstax}, §6.5.}
\end{ejercicio}

\begin{ejercicio}{2}{eje:5-exponenciales}
Resuelva: (a) $2^{x+1}=3^x$; (b) $4^x-3\cdot2^x-4=0$; (c) $e^{2x}-5e^x+6=0$.
\fuente{adaptado de OpenStax \emph{College Algebra}~\cite{openstax}, §6.6.}
\end{ejercicio}

\begin{ejercicio}{2}{eje:5-logaritmicas}
Resuelva: (a) $\log_2x+\log_2(x-2)=3$; (b) $\ln(x+1)-\ln(x-1)=\ln3$.
\fuente{adaptado de OpenStax \emph{College Algebra}~\cite{openstax}, §6.6.}
\end{ejercicio}

\begin{ejercicio}{2}{eje:5-inecuaciones}
Resuelva: (a) $0{,}5^x>8$; (b) $\log_3(2x-1)\le2$.
\end{ejercicio}

\begin{ejercicio}{2}{eje:5-lse}
(a) Calcule $\LSE(z)$ y $\softmax(z)$ para $z=(2,1,0,-1)$ y compruebe las cotas de la \cref{prop:lse}. (b) Calcule de forma estable $\LSE(1000,999,998)$.
\end{ejercicio}

\begin{ejercicio}{2}{eje:5-perplejidad}
Un modelo asigna a un texto de $N$ tokens una log-verosimilitud media de $-2{,}3$ nats por token. (a) Exprese esa cantidad en bits por token. (b) Calcule la perplejidad $e^{2{,}3}$ y su equivalente $2^{H}$ con $H$ en bits. (c) ¿Qué probabilidad constante por token daría la misma perplejidad?
\end{ejercicio}

\begin{ejercicio}{2}{eje:5-ley-potencia}
Una pérdida sigue la ley de potencia $L(N)=(N_c/N)^{\alpha}$ con $\alpha=0{,}076$. (a) ¿En qué factor se multiplica $L$ al multiplicar $N$ por~10? (b) ¿Por qué factor hay que multiplicar $N$ para reducir $L$ a la mitad? (c) ¿Qué pendiente tiene la recta $\log L$ frente a $\log N$?
\fuente{adaptado de OpenStax \emph{College Algebra}~\cite{openstax}, §6.7 (modelos exponenciales y logarítmicos).}
\end{ejercicio}

\begin{ejercicio}{3}{eje:5-log23}
Demuestre que $\log_23$ es irracional.
\end{ejercicio}

\begin{ejercicio}{3}{eje:5-temperatura}
Sea $z\in\R^n$ con máximo $M$, y $T>0$. Demuestre que $M\le T\,\LSE(z/T)\le M+T\ln n$ y deduzca que $T\,\LSE(z/T)\to M$ cuando $T\to0^{+}$.
\end{ejercicio}

\begin{ejercicio}{3}{eje:5-lse-media}
Demuestre, usando la desigualdad de las medias, que $\LSE(z)\ge\frac1n\sum_{i=1}^nz_i+\ln n$ para todo $z\in\R^n$. ¿Cuándo se da la igualdad?
\end{ejercicio}

\begin{ejercicio}{3}{eje:5-gibbs}
Sean $p=(p_1,\dots,p_n)$ y $q=(q_1,\dots,q_n)$ distribuciones de probabilidad con todos los $p_i,q_i>0$. Demuestre, usando el \cref{cor:lnx}, que $\sum_ip_i\ln\frac{q_i}{p_i}\le0$, es decir, $-\sum_ip_i\ln q_i\ge-\sum_ip_i\ln p_i$.
\fuente{adaptado de Cover y Thomas~\cite{ct}, cap.~2 (desigualdad de la información).}
\end{ejercicio}

\begin{ejercicio}{3}{eje:5-e-creciente}
Aplicando la desigualdad de las medias a los $n+1$ números $1+\frac1n,\dots,1+\frac1n,\,1$, demuestre que $\bigl(1+\frac1n\bigr)^n<\bigl(1+\frac1{n+1}\bigr)^{n+1}$ para todo $n\in\N$.
\end{ejercicio}

% =====================================================================
\chapter{Sistemas de numeración y coma flotante}
\label{tema:6}
% =====================================================================

Un ordenador solo almacena bits. Este tema estudia cómo se representan con ellos los enteros sin signo, los enteros con signo y los números reales, y qué se pierde por el camino. La representación posicional se fundamenta en la división euclídea de enteros; la parte de máquina sigue el capítulo~2 «Representing and Manipulating Information» de Bryant y O'Hallaron~\cite{csapp} y el artículo de Goldberg~\cite{goldberg}; los formatos de coma flotante son los de la norma IEEE~754~\cite{ieee754}.

\section{Representación posicional}

\begin{teorema}[Representación en base $b$]\label{teo:base}
Sea $b\ge2$ un entero. Todo $n\in\N$ se escribe de manera única como
\[
n=\sum_{k=0}^{m}d_kb^k=(d_md_{m-1}\dots d_1d_0)_b,\qquad d_k\in\{0,1,\dots,b-1\},\ d_m\neq0.
\]
Los $d_k$ son las \emph{cifras} de $n$ en base $b$.
\end{teorema}

\begin{proof}
Se usa la división euclídea de enteros: para $n\in\N_0$ existen únicos $q,r$ con $n=bq+r$ y $0\le r<b$.

\emph{Existencia}, por inducción fuerte sobre $n$. Si $n<b$, basta tomar $m=0$ y $d_0=n$. Si $n\ge b$, se escribe $n=bq+r$ con $0\le r<b$; entonces $1\le q<n$ y, por la hipótesis de inducción, $q=\sum_{k=0}^{m'}e_kb^k$. Luego $n=r+\sum_{k=0}^{m'}e_kb^{k+1}$, que es una representación de $n$ con $d_0=r$ y $d_{k+1}=e_k$.

\emph{Unicidad}, también por inducción fuerte. Si $n=\sum_kd_kb^k=\sum_kd'_kb^k$, entonces $n=b\bigl(\sum_{k\ge1}d_kb^{k-1}\bigr)+d_0$ con $0\le d_0<b$, y lo mismo con las $d'_k$. Por la unicidad del resto, $d_0=d'_0$, y por la del cociente, $\sum_{k\ge1}d_kb^{k-1}=\sum_{k\ge1}d'_kb^{k-1}$, un número menor que $n$ al que se aplica la hipótesis de inducción (si es~0, no quedan más cifras no nulas).
\end{proof}

La demostración es un algoritmo: las cifras de $n$ en base $b$ son los restos de dividir sucesivamente entre $b$, leídos de abajo arriba. En sentido contrario, el valor de $(d_m\dots d_0)_b$ se obtiene con el \emph{esquema de Horner}, $(\cdots((d_mb+d_{m-1})b+d_{m-2})\cdots)b+d_0$, que es la regla de Ruffini del \cref{tema:3} evaluando en $x=b$.

\begin{ejemplo}\label{ej:conversion}
$45=2\cdot22+1$, $22=2\cdot11+0$, $11=2\cdot5+1$, $5=2\cdot2+1$, $2=2\cdot1+0$, $1=2\cdot0+1$. Leyendo los restos de abajo arriba, $45=(101101)_2$. Comprobación por Horner: $((((1\cdot2+0)2+1)2+1)2+0)2+1=45$.
\end{ejemplo}

\begin{proposicion}[Número de cifras]\label{prop:cifras}
Un entero $n\ge1$ tiene exactamente $\lfloor\log_bn\rfloor+1$ cifras en base $b$. Con $m$ cifras en base $b$ se representan exactamente los $b^m$ enteros $0,1,\dots,b^m-1$.
\end{proposicion}

\begin{proof}
Si $n$ tiene $m+1$ cifras con $d_m\neq0$, entonces $b^m\le n\le\sum_{k=0}^{m}(b-1)b^k=b^{m+1}-1$ por la suma geométrica (\cref{teo:geometrica}). Tomando $\log_b$, $m\le\log_bn<m+1$, es decir, $m=\lfloor\log_bn\rfloor$. La segunda afirmación se sigue del \cref{teo:base} y de que $(b-1\ \dots\ b-1)_b=b^m-1$.
\end{proof}

Las cifras de la parte fraccionaria de un número $x\in[0,1)$ se obtienen multiplicando sucesivamente por $b$ y tomando la parte entera: $x=\sum_{k\ge1}d_kb^{-k}=(0{,}d_1d_2\dots)_b$.

\begin{proposicion}\label{prop:fraccion-binaria}
Sea $p/q$ una fracción irreducible con $0<p<q$. Su desarrollo en base~2 es finito si y solo si $q$ es una potencia de~2.
\end{proposicion}

\begin{proof}
Si el desarrollo tiene $k$ cifras, $p/q=N/2^k$ con $N$ entero, es decir, $p\,2^k=qN$. Entonces $q$ divide a $p\,2^k$ y, como no tiene factores comunes con $p$, divide a $2^k$, luego es potencia de~2. Recíprocamente, si $q=2^k$, entonces $p<2^k$ se escribe con $k$ cifras binarias, $p=(c_{k-1}\dots c_0)_2$, y $p/2^k=(0{,}c_{k-1}\dots c_0)_2$.
\end{proof}

Por ejemplo, $0{,}1=1/10$ tiene denominador $10=2\cdot5$, de modo que su desarrollo binario es infinito: $0{,}1=(0{,}0\overline{0011})_2$. Ningún ordenador binario almacena $0{,}1$ exactamente.

\section{Binario, hexadecimal y bytes}

Un \emph{bit} es una cifra binaria y un \emph{byte} son 8 bits, que representan $2^8=256$ valores. Como $16=2^4$, cada cifra \emph{hexadecimal} ($0$--$9$, $\texttt{A}$--$\texttt{F}$) corresponde a un grupo de 4 bits, y un byte se escribe con dos cifras hexadecimales, de $\texttt{0x00}$ a $\texttt{0xFF}$. Para convertir entre binario y hexadecimal basta agrupar los bits de cuatro en cuatro desde la derecha: $(1011\,0110)_2=\texttt{0xB6}=182$.

Para cantidades de memoria se usan dos sistemas de prefijos: los decimales del SI ($1\ \text{kB}=10^3$ bytes, $1\ \text{GB}=10^9$ bytes) y los binarios de la IEC ($1\ \text{KiB}=2^{10}$ bytes, $1\ \text{GiB}=2^{30}$ bytes). La diferencia es de un $7{,}4\,\%$ en el caso de los gigas.

\section{Enteros con signo: complemento a dos}

\begin{definicion}[Complemento a dos]\label{def:c2}
Una palabra de $n$ bits $w=(a_{n-1}\dots a_0)$ representa, \emph{sin signo}, el entero $U(w)=\sum_{k=0}^{n-1}a_k2^k$ y, en \emph{complemento a dos}, el entero
\[
T(w)=-a_{n-1}2^{n-1}+\sum_{k=0}^{n-2}a_k2^k=U(w)-a_{n-1}2^n.
\]
\end{definicion}

\begin{teorema}\label{teo:c2}
Para $n\ge2$:
\begin{enumerate}
  \item $T$ es una biyección entre las palabras de $n$ bits y los enteros del intervalo $[-2^{n-1},\,2^{n-1}-1]$; $T(w)<0$ si y solo si $a_{n-1}=1$;
  \item $T(w)$ y $U(w)$ difieren en un múltiplo de $2^n$;
  \item si $\bar w$ es la palabra con todos los bits invertidos, $T(\bar w)+1=-T(w)$; en particular, el opuesto de $T(w)$ se obtiene invirtiendo los bits y sumando~1, y el resultado es representable en $n$ bits salvo si $T(w)=-2^{n-1}$.
\end{enumerate}
\end{teorema}

\begin{proof}
(i) Por la \cref{prop:cifras}, $U$ es una biyección sobre $\{0,\dots,2^n-1\}$. Si $a_{n-1}=0$, $U(w)<2^{n-1}$ y $T(w)=U(w)\in[0,2^{n-1}-1]$. Si $a_{n-1}=1$, $U(w)\in[2^{n-1},2^n-1]$ y $T(w)=U(w)-2^n\in[-2^{n-1},-1]$. Cada mitad se aplica biyectivamente sobre su imagen, y las imágenes son disjuntas y cubren el intervalo. (ii) Es la igualdad $T(w)=U(w)-a_{n-1}2^n$. (iii) $U(w)+U(\bar w)=\sum_k2^k=2^n-1$, luego $T(w)+T(\bar w)=2^n-1-(a_{n-1}+\bar a_{n-1})2^n=-1$, porque exactamente uno de los dos bits de signo vale~1. Así, $T(\bar w)+1=-T(w)$ como igualdad de enteros; si $T(w)=-2^{n-1}$, el valor $-T(w)=2^{n-1}$ no está en el rango de~(i).
\end{proof}

\begin{ejemplo}\label{ej:c2}
Con $n=8$ el rango es $[-128,127]$. Como $5=(0000\,0101)_2$, se tiene $-5=(1111\,1010)_2+1=(1111\,1011)_2$; en efecto, $-128+64+32+16+8+2+1=-5$.
\end{ejemplo}

Por el \cref{teo:c2}(ii), la suma de enteros con signo se realiza con el mismo circuito que la suma sin signo, descartando el acarreo final: el resultado es correcto salvo que salga del rango, situación que se llama \emph{desbordamiento} (\cref{eje:6-desbordamiento}).

\section{Coma flotante}

\begin{definicion}[Número en coma flotante normalizado]\label{def:flotante}
Un formato binario de coma flotante con $e$ bits de exponente, sesgo $\beta=2^{e-1}-1$ y $p$ bits de fracción representa la palabra $(s,E,f)$, con $1\le E\le2^e-2$, como
\[
x=(-1)^s\,(1{,}f_1f_2\dots f_p)_2\cdot2^{E-\beta}.
\]
Los valores extremos de $E$ se reservan para el cero, los números subnormales, los infinitos y \texttt{NaN}~\cite{ieee754}.
\end{definicion}

\begin{table}[h]
\centering\small
\caption{Formatos binarios de coma flotante habituales. El épsilon es $\varepsilon=2^{-p}$, la distancia de~1 al siguiente número representable.}\label{tab:formatos}
\begin{tabular}{@{}lcccccc@{}}
\toprule
Formato & Bits & Signo & Exponente $e$ & Fracción $p$ & $\varepsilon=2^{-p}$ & Máximo finito \\
\midrule
fp64 (doble) & 64 & 1 & 11 & 52 & $2{,}2\cdot10^{-16}$ & $1{,}8\cdot10^{308}$ \\
fp32 (simple) & 32 & 1 & 8 & 23 & $1{,}2\cdot10^{-7}$ & $3{,}4\cdot10^{38}$ \\
fp16 (media) & 16 & 1 & 5 & 10 & $9{,}8\cdot10^{-4}$ & $65\,504$ \\
bf16 (bfloat16) & 16 & 1 & 8 & 7 & $7{,}8\cdot10^{-3}$ & $3{,}4\cdot10^{38}$ \\
\bottomrule
\end{tabular}
\end{table}

El formato bfloat16 conserva los 8 bits de exponente de fp32 y recorta la fracción a 7 bits: tiene el mismo rango que fp32 con mucha menos precisión, y sus 16 bits son exactamente los 16 bits superiores de un fp32.

\begin{proposicion}[Error de redondeo]\label{prop:redondeo}
Sea $\fl(x)$ el número del formato más cercano a $x$. Si $x$ está en el rango de los números normalizados, entonces
\[
\fl(x)=x(1+\delta)\qquad\text{con}\qquad\abs{\delta}\le u=2^{-(p+1)}=\frac\varepsilon2.
\]
\end{proposicion}

La demostración es el \cref{eje:6-redondeo}. El número $u$ se llama \emph{unidad de redondeo}: el error \emph{relativo} es pequeño, pero el absoluto crece con $\abs{x}$, porque la distancia entre números consecutivos en $[2^k,2^{k+1})$ es $2^{k-p}$~\cite{goldberg}.

\begin{observacion}[Errores frecuentes]
\begin{enumerate}[label=(\alph*)]
\item Confundir GB y GiB, o kB y KiB.
\item Olvidar que el rango del complemento a dos es asimétrico: en 8 bits, $-(-128)$ desborda.
\item Comparar números en coma flotante con igualdad exacta: $0{,}1+0{,}2\neq0{,}3$ en doble precisión, porque ninguno de los tres es representable.
\item Creer que fp16 y bf16 son intercambiables: el primero tiene más precisión y el segundo mucho más rango.
\end{enumerate}
\end{observacion}

\begin{observacion}[Conexión con el curso]
Los tokenizadores BPE a nivel de byte (módulo~12) parten de un vocabulario inicial de exactamente $2^8=256$ símbolos, uno por byte, y el texto se convierte antes en bytes mediante UTF-8 (\cref{eje:6-utf8}). El entrenamiento en precisión mixta con bf16 o fp16 (módulo~11) y la cuantización a enteros de 8 o 4 bits (módulo~17) son elecciones entre rango, precisión y memoria que se analizan con las herramientas de este tema.
\end{observacion}

\begin{lecturas}
\item Bryant y O'Hallaron~\cite{csapp}, §2.1 «Information Storage» (hexadecimal, bytes), §2.2 «Integer Representations» (complemento a dos, extensión de signo), §2.3 «Integer Arithmetic» (desbordamiento) y §2.4 «Floating Point».
\item Goldberg~\cite{goldberg}, sección sobre errores de redondeo (épsilon de la máquina y unidad de redondeo).
\item IEEE 754~\cite{ieee754}, para la definición de los formatos y del redondeo al par más cercano.
\end{lecturas}

\section*{Ejercicios del tema 6}
\addcontentsline{toc}{section}{Ejercicios del tema 6}

\begin{ejercicio}{1}{eje:6-a-decimal}
Pase a base~10: $(1011\,0110)_2$, $\texttt{0x2F}$, $\texttt{0xFF}$ y $(777)_8$.
\fuente{adaptado de Bryant y O'Hallaron~\cite{csapp}, §2.1, problemas de práctica de conversión.}
\end{ejercicio}

\begin{ejercicio}{1}{eje:6-contar}
(a) ¿Cuántos valores distintos pueden representarse con 8, 16 y 32 bits? (b) ¿Cuál es el rango de un entero sin signo de 8 bits y el de uno con signo en complemento a dos? (c) ¿Cuántas cifras hexadecimales ocupa una palabra de 32 bits?
\end{ejercicio}

\begin{ejercicio}{1}{eje:6-vf}
Razone si son verdaderas o falsas: (a) $0{,}1$ es exactamente representable en fp32; (b) bf16 y fp32 tienen aproximadamente el mismo máximo finito; (c) en complemento a dos, el bit más significativo indica el signo; (d) $0{,}625$ es exactamente representable en fp16.
\end{ejercicio}

\begin{ejercicio}{2}{eje:6-a-binario}
Escriba en binario y en hexadecimal: $200$, $1000$ y $4095$. Compruebe cada resultado con el esquema de Horner.
\fuente{adaptado de Bryant y O'Hallaron~\cite{csapp}, §2.1.}
\end{ejercicio}

\begin{ejercicio}{2}{eje:6-fracciones}
Escriba en base~2: (a) $0{,}625$; (b) $13/16$; (c) $1/3$; (d) $0{,}1$. Explique, con la \cref{prop:fraccion-binaria}, cuáles tienen desarrollo finito.
\end{ejercicio}

\begin{ejercicio}{2}{eje:6-c2}
Con palabras de 8 bits en complemento a dos: (a) represente $-1$, $-100$ y $127$; (b) interprete $(1000\,0001)_2$ y $(1111\,0110)_2$; (c) calcule $100+50$ con aritmética de 8 bits e interprete el resultado.
\fuente{adaptado de Bryant y O'Hallaron~\cite{csapp}, §2.2--2.3.}
\end{ejercicio}

\begin{ejercicio}{2}{eje:6-utf8}
UTF-8 codifica los puntos de código de \texttt{U+0080} a \texttt{U+07FF} con dos bytes, y los de \texttt{U+0800} a \texttt{U+FFFF} con tres bytes, según las plantillas
\[
\texttt{110xxxxx 10xxxxxx}\qquad\text{y}\qquad\texttt{1110xxxx 10xxxxxx 10xxxxxx},
\]
donde las $x$ son los bits del punto de código, completados con ceros a la izquierda. Codifique «é» (\texttt{U+00E9}) y «€» (\texttt{U+20AC}) y dé los bytes en hexadecimal.
\end{ejercicio}

\begin{ejercicio}{2}{eje:6-digitos}
Usando la \cref{prop:cifras}: (a) ¿cuántas cifras binarias tiene $10^9$? (b) ¿cuántas cifras decimales tiene $2^{100}$? (Use $\log_{10}2\approx0{,}30103$.) (c) ¿Cuántas cifras hexadecimales tiene $2^{32}-1$?
\end{ejercicio}

\begin{ejercicio}{2}{eje:6-decodificar}
(a) Interprete las palabras fp32 $\texttt{0x3F800000}$ y $\texttt{0xC0A00000}$. (b) Codifique $0{,}75$ y $3$ en bf16 y dé el resultado en hexadecimal.
\fuente{adaptado de Bryant y O'Hallaron~\cite{csapp}, §2.4.}
\end{ejercicio}

\begin{ejercicio}{2}{eje:6-espaciado}
En bf16 ($p=7$) con redondeo al más cercano (en caso de empate, al de última cifra par): (a) ¿cuál es la distancia entre números consecutivos en $[2^k,2^{k+1})$? (b) Calcule $\fl(256+1)$, $\fl(\fl(256+1)+1)$ y $\fl(256+\fl(1+1))$. (c) ¿Cuál es el menor entero positivo no representable en bf16? ¿Y en fp32? (d) Compruebe que el máximo finito de fp16 es $65\,504$.
\end{ejercicio}

\begin{ejercicio}{2}{eje:6-memoria}
Un modelo tiene $7\cdot10^{9}$ parámetros. Calcule la memoria que ocupan sus pesos en fp32, bf16, int8 y con 4 bits por parámetro, en GB y, para fp32, también en GiB.
\end{ejercicio}

\begin{ejercicio}{3}{eje:6-divisibilidad}
Demuestre que un entero es divisible por~3 (respectivamente, por~9) si y solo si la suma de sus cifras decimales lo es. Enuncie y demuestre el criterio análogo para la divisibilidad por~15 en base~16. (Indicación: \cref{eje:3-diferencia-potencias}.)
\end{ejercicio}

\begin{ejercicio}{3}{eje:6-extension}
Demuestre que la \emph{extensión de signo} conserva el valor en complemento a dos: si $w$ es una palabra de $n$ bits y $w'$ la de $n+1$ bits que se obtiene repitiendo a la izquierda el bit de signo, entonces $T(w')=T(w)$.
\fuente{adaptado de Bryant y O'Hallaron~\cite{csapp}, §2.2.}
\end{ejercicio}

\begin{ejercicio}{3}{eje:6-decimal-finito}
Demuestre que una fracción irreducible $p/q$ tiene desarrollo decimal finito si y solo si $q=2^a5^b$ con $a,b\in\N_0$.
\end{ejercicio}

\begin{ejercicio}{3}{eje:6-desbordamiento}
Sean $x,y$ enteros en $[-2^{n-1},2^{n-1}-1]$ y sea $s$ el resultado de la suma en complemento a dos, es decir, el único entero del rango que difiere de $x+y$ en un múltiplo de $2^n$. Demuestre: (a) si $x$ e $y$ tienen signos distintos, no hay desbordamiento ($s=x+y$); (b) hay desbordamiento si y solo si $x$ e $y$ tienen el mismo signo y $s$ tiene el signo contrario.
\fuente{adaptado de Bryant y O'Hallaron~\cite{csapp}, §2.3.}
\end{ejercicio}

\begin{ejercicio}{3}{eje:6-redondeo}
Demuestre la \cref{prop:redondeo}: si $2^k\le\abs{x}<2^{k+1}$ y el formato tiene $p$ bits de fracción, entonces $\abs{\fl(x)-x}\le2^{k-p-1}$ y $\abs{\fl(x)-x}/\abs{x}\le2^{-(p+1)}$.
\fuente{adaptado de Goldberg~\cite{goldberg}, sección sobre errores de redondeo.}
\end{ejercicio}


\appendix
% =====================================================================
\chapter{Soluciones de los ejercicios}
\label{ap:soluciones}
% =====================================================================

Cada solución se presenta en cinco partes: datos, planteamiento, desarrollo, resultado y comprobación. En las demostraciones, los datos son las hipótesis, el planteamiento es la estrategia y la comprobación es un caso particular o una verificación independiente. Los valores aproximados se dan con cuatro cifras significativas salvo que se indique otra cosa.

% ---------------------------------------------------------------------
\section{Tema 1}

\begin{solucion}{eje:1-clasificar}
\datos los números $-3$, $0{,}75$, $0{,}\overline3$, $\sqrt9$, $\sqrt8$, $\pi$, $1/7$, $0$.
\planteamiento escribir cada número como fracción si es posible (\cref{prop:decimal-periodico}).
\pasos
\begin{enumerate}[label=\arabic*.]
  \item $-3\in\Z$, no es natural. $0\in\Z$ y no está en $\N$ con el convenio de estos apuntes.
  \item $0{,}75=3/4$ y $0{,}\overline3=1/3$ son racionales no enteros; $1/7\in\Q$.
  \item $\sqrt9=3\in\N$.
  \item $\sqrt8=2\sqrt2$ es irracional: si fuera racional, $\sqrt2=\sqrt8/2$ también lo sería, en contra del \cref{teo:raiz2}. $\pi$ es irracional (resultado conocido, demostrado en~\cite[cap.~16]{spivak}).
\end{enumerate}
\resultado \textbf{$\N$: $\sqrt9$. $\Z$: además $-3$ y $0$. $\Q$: además $0{,}75$, $0{,}\overline3$ y $1/7$. $\R\setminus\Q$: $\sqrt8$ y $\pi$.} Todos son reales.
\comprobacion las inclusiones $\N\subset\Z\subset\Q$ hacen que cada número de un conjunto pertenezca también a los siguientes.
\end{solucion}

\begin{solucion}{eje:1-axiomas}
\datos $a+x=a+y$.
\planteamiento identificar en cada paso qué propiedad se usa.
\pasos
\begin{enumerate}[label=\arabic*.]
  \item (1) neutro de la suma (C3), con la conmutatividad (C2) para escribir $0+x=x+0$.
  \item (2) opuesto (C4), con (C2): $(-a)+a=a+(-a)=0$.
  \item (3) asociatividad (C1).
  \item (4) la hipótesis $a+x=a+y$ (no es un axioma).
  \item (5) asociatividad (C1). (6) opuesto (C4). (7) neutro (C3).
\end{enumerate}
\resultado \textbf{(1) C3; (2) C4; (3) C1; (4) hipótesis; (5) C1; (6) C4; (7) C3} (con C2 donde se cambia el orden).
\comprobacion es la ley de cancelación de la suma; con $y=0$ se obtiene la \cref{prop:cuerpo}(i).
\end{solucion}

\begin{solucion}{eje:1-vf}
\planteamiento buscar un contraejemplo o una demostración por reducción al absurdo.
\pasos
\begin{enumerate}[label=(\alph*)]
  \item Falsa: $\sqrt2$ y $-\sqrt2$ son irracionales y su suma es $0\in\Q$.
  \item Verdadera: si $r\neq0$ es racional, $x$ irracional y $rx=s\in\Q$, entonces $x=s/r\in\Q$, contradicción.
  \item Verdadera: si $a<b$ son racionales, $(a+b)/2$ es racional y está entre ambos (\cref{eje:1-epsilon}).
\end{enumerate}
\resultado \textbf{(a) falsa; (b) verdadera; (c) verdadera.}
\comprobacion en (c), con $a=\frac13$, $b=\frac12$: $\frac{5}{12}$ cumple $\frac4{12}<\frac5{12}<\frac6{12}$.
\end{solucion}

\begin{solucion}{eje:1-fracciones}
\datos $0{,}1\overline6$, $2{,}\overline{45}$, $0{,}\overline{142857}$.
\planteamiento multiplicar por potencias de~10 para eliminar el periodo (\cref{ej:periodico}).
\pasos
\begin{enumerate}[label=(\alph*)]
  \item $x=0{,}1\overline6$: $100x=16{,}\overline6$ y $10x=1{,}\overline6$; restando, $90x=15$, luego $x=15/90=1/6$.
  \item $x=2{,}\overline{45}$: $100x-x=245{,}\overline{45}-2{,}\overline{45}=243$, luego $x=243/99=27/11$.
  \item $x=0{,}\overline{142857}$: $10^6x-x=142857$, luego $x=142857/999999=1/7$ (dividiendo numerador y denominador por $142857$).
\end{enumerate}
\resultado \textbf{(a) $1/6$; (b) $27/11$; (c) $1/7$.}
\comprobacion $1\div6=0{,}1666\ldots$, $27\div11=2{,}4545\ldots$ y $1\div7=0{,}142857\,142857\ldots$
\end{solucion}

\begin{solucion}{eje:1-potencias}
\planteamiento escribir todo como potencias de primos y aplicar la \cref{prop:exponentes}.
\pasos
\begin{enumerate}[label=(\alph*)]
  \item $\dfrac{2^{-3}\cdot2^{4}}{2^{-3}}=2^{-3+4+3}=2^4=16$.
  \item $27^{2/3}=(3^3)^{2/3}=3^2=9$; $16^{-3/4}=(2^4)^{-3/4}=2^{-3}=\frac18$. Producto: $\frac98$.
  \item $\sqrt{50}=5\sqrt2$, $\sqrt{18}=3\sqrt2$, $\sqrt8=2\sqrt2$; total $(5-3+2)\sqrt2=4\sqrt2$.
\end{enumerate}
\resultado \textbf{(a) $16$; (b) $9/8$; (c) $4\sqrt2$.}
\comprobacion (c): $7{,}0711-4{,}2426+2{,}8284=5{,}6569\approx4\cdot1{,}4142$.
\end{solucion}

\begin{solucion}{eje:1-racionalizar}
\planteamiento multiplicar por el conjugado, o por el factor que completa una potencia exacta.
\pasos
\begin{enumerate}[label=(\alph*)]
  \item $\dfrac{1}{\sqrt3-\sqrt2}\cdot\dfrac{\sqrt3+\sqrt2}{\sqrt3+\sqrt2}=\dfrac{\sqrt3+\sqrt2}{3-2}=\sqrt3+\sqrt2$.
  \item $\dfrac{2(\sqrt5-1)}{(\sqrt5+1)(\sqrt5-1)}=\dfrac{2(\sqrt5-1)}{4}=\dfrac{\sqrt5-1}2$.
  \item $\dfrac{6}{2^{1/3}}\cdot\dfrac{2^{2/3}}{2^{2/3}}=\dfrac{6\cdot2^{2/3}}{2}=3\sqrt[3]4$.
\end{enumerate}
\resultado \textbf{(a) $\sqrt3+\sqrt2$; (b) $(\sqrt5-1)/2$; (c) $3\sqrt[3]4$.}
\comprobacion (b): $2/(1+2{,}2361)=0{,}6180$ y $(2{,}2361-1)/2=0{,}6180$.
\end{solucion}

\begin{solucion}{eje:1-ordenar}
\datos $3{,}1415<\pi<3{,}1416$.
\planteamiento los tres números son positivos; por el \cref{eje:1-potencia-n}, basta comparar sus cuadrados.
\pasos
\begin{enumerate}[label=\arabic*.]
  \item $(\sqrt2+\sqrt3)^2=5+2\sqrt6$ y $(\sqrt{10})^2=10$. Como $\sqrt6<2{,}5$ (porque $6<6{,}25$), $5+2\sqrt6<10$, luego $\sqrt2+\sqrt3<\sqrt{10}$.
  \item $\pi^2<3{,}1416^2=9{,}86965\ldots<9{,}8697$. Por otro lado, $\sqrt6>2{,}449$ porque $2{,}449^2=5{,}997601<6$, luego $5+2\sqrt6>9{,}898>9{,}8697>\pi^2$, y $\pi<\sqrt2+\sqrt3$.
\end{enumerate}
\resultado \textbf{$\pi<\sqrt2+\sqrt3<\sqrt{10}$.}
\comprobacion con calculadora: $3{,}1416<3{,}1463<3{,}1623$.
\end{solucion}

\begin{solucion}{eje:1-cientifica}
\planteamiento agrupar mantisas y potencias de~10 por separado.
\pasos
\begin{enumerate}[label=\arabic*.]
  \item Mantisas: $3{,}2\cdot4/1{,}6=12{,}8/1{,}6=8$. Potencias: $10^{5-8-(-2)}=10^{-1}$.
  \item (a) $8\cdot10^{-1}$.
  \item (b) $7\cdot10^{10}\cdot2=1{,}4\cdot10^{11}$ bytes $=1{,}4\cdot10^{11}/10^9=140$ GB.
\end{enumerate}
\resultado \textbf{(a) $8\cdot10^{-1}=0{,}8$; (b) $1{,}4\cdot10^{11}$ bytes, es decir, $140$ GB.}
\comprobacion (a): $3{,}2\cdot10^5\cdot4\cdot10^{-8}=1{,}28\cdot10^{-2}$, y $1{,}28\cdot10^{-2}/1{,}6\cdot10^{-2}=0{,}8$.
\end{solucion}

\begin{solucion}{eje:1-intervalos}
\datos $A=[-2,3)$, $B=(1,5]$.
\planteamiento dibujar ambos intervalos en la recta y respetar qué extremos se incluyen.
\pasos
\begin{enumerate}[label=\arabic*.]
  \item $A\cap B$: $x>1$ y $x<3$; $(1,3)$.
  \item $A\cup B$: desde $-2$ (incluido) hasta $5$ (incluido), sin huecos porque los intervalos se solapan; $[-2,5]$.
  \item $A\setminus B$: puntos de $A$ con $x\le1$; $[-2,1]$.
  \item $\R\setminus A$: $x<-2$ o $x\ge3$; $(-\infty,-2)\cup[3,\infty)$.
\end{enumerate}
\resultado \textbf{$(1,3)$; $[-2,5]$; $[-2,1]$; $(-\infty,-2)\cup[3,\infty)$.}
\comprobacion $1\in A\setminus B$ porque $1\in A$ y $1\notin B$; $3\notin A$, luego $3\in\R\setminus A$.
\end{solucion}

\begin{solucion}{eje:1-ecuacion-producto}
\datos $x^3=x$.
\planteamiento pasar todo a un miembro, factorizar y aplicar la \cref{prop:cuerpo}(iv).
\pasos
\begin{enumerate}[label=\arabic*.]
  \item $x^3-x=0$, es decir, $x(x^2-1)=x(x-1)(x+1)=0$.
  \item Por la \cref{prop:cuerpo}(iv), aplicada dos veces, algún factor es cero: $x=0$, $x=1$ o $x=-1$.
  \item Si se divide por $x$ se obtiene $x^2=1$, que pierde la solución $x=0$: la división solo es válida si $x\neq0$.
\end{enumerate}
\resultado \textbf{$x\in\{-1,0,1\}$.} Dividir por $x$ pierde $x=0$.
\comprobacion $(-1)^3=-1$, $0^3=0$, $1^3=1$.
\end{solucion}

\begin{solucion}{eje:1-quebrados}
\datos $b,d\neq0$; $a/b=ab^{-1}$; $(bd)^{-1}=b^{-1}d^{-1}$.
\planteamiento escribir cada cociente como producto por el inverso y usar (C1), (C2), (C6).
\pasos
\begin{enumerate}[label=\arabic*.]
  \item $bd\neq0$ por la \cref{prop:cuerpo}(iv), así que $(bd)^{-1}$ existe.
  \item $\frac{ad+bc}{bd}=(ad+bc)b^{-1}d^{-1}=adb^{-1}d^{-1}+bcb^{-1}d^{-1}$ (C6).
  \item Reordenando (C1, C2): $adb^{-1}d^{-1}=ab^{-1}(dd^{-1})=ab^{-1}$ y $bcb^{-1}d^{-1}=cd^{-1}(bb^{-1})=cd^{-1}$, por (C5) y (C3).
  \item Luego $\frac{ad+bc}{bd}=ab^{-1}+cd^{-1}=\frac ab+\frac cd$.
  \item Producto: $\frac ab\cdot\frac cd=ab^{-1}cd^{-1}=(ac)(b^{-1}d^{-1})=(ac)(bd)^{-1}=\frac{ac}{bd}$.
\end{enumerate}
\resultado \textbf{Ambas fórmulas se deducen de los axiomas de cuerpo.}
\comprobacion $\frac12+\frac13=\frac{1\cdot3+2\cdot1}{6}=\frac56$ y $\frac12\cdot\frac23=\frac26=\frac13$.
\end{solucion}

\begin{solucion}{eje:1-negativo}
\datos axiomas (O1)--(O3) y la \cref{prop:orden}.
\planteamiento reducir cada desigualdad a la pertenencia de una diferencia a $P$.
\pasos
\begin{enumerate}[label=(\alph*)]
  \item $c<0$ significa $-c\in P$, y $a<b$ significa $b-a\in P$. Por (O3), $(b-a)(-c)\in P$, es decir, $ac-bc\in P$, que es $bc<ac$.
  \item Sea $a>0$. Como $aa^{-1}=1\neq 0$, $a^{-1}\neq0$. Si fuera $a^{-1}<0$, por (a) con $0<a$ y $c=a^{-1}$ se tendría $0\cdot a^{-1}>a\cdot a^{-1}$, es decir, $0>1$, en contra de la \cref{prop:orden}(iv). Por tricotomía, $a^{-1}>0$.
  \item Si $0<a<b$, entonces $ab>0$ por (O3) y $(ab)^{-1}>0$ por (b). Multiplicando $a<b$ por $(ab)^{-1}$ (\cref{prop:orden}(iii)): $a(ab)^{-1}<b(ab)^{-1}$, es decir, $b^{-1}<a^{-1}$.
\end{enumerate}
\resultado \textbf{Las tres propiedades se deducen de los axiomas de orden.}
\comprobacion con $a=2$, $b=4$: $\frac14<\frac12$; con $c=-1$, $2<4$ da $-2>-4$.
\end{solucion}

\begin{solucion}{eje:1-potencia-n}
\datos $0\le a<b$.
\planteamiento inducción sobre $n$, multiplicando desigualdades de números no negativos.
\pasos
\begin{enumerate}[label=\arabic*.]
  \item Caso $n=1$: es la hipótesis.
  \item Paso: supóngase $a^n<b^n$. Como $a\ge0$, $a^{n+1}=a^n\cdot a\le a^n\cdot b$ (si $a=0$ es una igualdad $0=0\le\dots$; si $a>0$, se usa $a<b$ y $a^n\ge0$).
  \item Como $b>0$, $a^nb<b^nb=b^{n+1}$ por la \cref{prop:orden}(iii).
  \item Luego $a^{n+1}\le a^nb<b^{n+1}$.
  \item Recíproco para $n=2$: si $a,b\ge0$ y $a^2<b^2$, no puede ser $b\le a$, porque entonces $b^2\le a^2$ (por lo anterior si $b<a$, o igualdad si $b=a$). Por tricotomía, $a<b$.
\end{enumerate}
\resultado \textbf{$a^n<b^n$ para todo $n\in\N$; y para no negativos, $a^2<b^2\Rightarrow a<b$.}
\comprobacion la hipótesis $a\ge0$ es necesaria: $-3<2$ pero $(-3)^2>2^2$.
\end{solucion}

\begin{solucion}{eje:1-raiz3}
\datos se quiere probar que $\sqrt3\notin\Q$ y $\sqrt2+\sqrt3\notin\Q$.
\planteamiento reducción al absurdo, como en el \cref{teo:raiz2}, con el análogo del \cref{lem:par} para el~3.
\pasos
\begin{enumerate}[label=\arabic*.]
  \item Lema: si $3\nmid n$, entonces $n=3k\pm1$ y $n^2=9k^2\pm6k+1=3(3k^2\pm2k)+1$, que no es múltiplo de~3. Por contrarrecíproco, si $3\mid n^2$, entonces $3\mid n$.
  \item Si $\sqrt3=p/q$ irreducible, $p^2=3q^2$, luego $3\mid p$; con $p=3r$, $9r^2=3q^2$, $q^2=3r^2$ y $3\mid q$. Contradicción.
  \item Si $\sqrt2+\sqrt3=s\in\Q$, entonces $s\neq0$ y, elevando al cuadrado, $s^2=5+2\sqrt6$, luego $\sqrt6=(s^2-5)/2\in\Q$.
  \item Pero $\sqrt6$ es irracional: si $\sqrt6=p/q$ irreducible, $p^2=6q^2$ es par, $p=2r$, $2r^2=3q^2$, luego $q^2$ es par (porque $3q^2$ es par y 3 es impar) y $q$ es par. Contradicción.
\end{enumerate}
\resultado \textbf{$\sqrt3$ y $\sqrt2+\sqrt3$ son irracionales.}
\comprobacion el argumento falla, como debe, para $\sqrt4$: $p^2=4q^2$ no conduce a contradicción porque $p/q=2/1$ es irreducible con $q$ impar.
\end{solucion}

\begin{solucion}{eje:1-epsilon}
\datos (a) $a<b$; (b) $x\le y+\varepsilon$ para todo $\varepsilon>0$.
\planteamiento (a) operar con la \cref{prop:orden}; (b) reducción al absurdo eligiendo un $\varepsilon$ concreto.
\pasos
\begin{enumerate}[label=\arabic*.]
  \item (a) Sumando $a$ a $a<b$: $2a<a+b$. Multiplicando por $\frac12>0$ (positivo por el \cref{eje:1-negativo}(b), ya que $2=1+1>0$): $a<\frac{a+b}2$. Análogamente, sumando $b$: $a+b<2b$ y $\frac{a+b}2<b$.
  \item (b) Supóngase $x>y$. Entonces $\varepsilon=\frac{x-y}2>0$ y, por hipótesis, $x\le y+\frac{x-y}2=\frac{x+y}2$.
  \item Pero por (a) con $y<x$, $\frac{x+y}2<x$. Contradicción; luego $x\le y$.
\end{enumerate}
\resultado \textbf{(a) el punto medio está estrictamente entre $a$ y $b$; (b) $x\le y$.}
\comprobacion en (b), la conclusión no puede mejorarse a $x<y$: con $x=y$ la hipótesis se cumple.
\end{solucion}
% ---------------------------------------------------------------------
\section{Tema 2}

\begin{solucion}{eje:2-calcular}
\planteamiento aplicar la \cref{def:abs} según el signo de lo que hay dentro de las barras.
\pasos
\begin{enumerate}[label=\arabic*.]
  \item $\abs{-3}=3$ y $\abs{2-7}=\abs{-5}=5$; la diferencia es $-2$.
  \item $\pi-4<0$, luego $\abs{\pi-4}=4-\pi$.
  \item Si $x<1$, $x-1<0$ y $\abs{x-1}=1-x$.
\end{enumerate}
\resultado \textbf{$-2$; $4-\pi$; $1-x$.}
\comprobacion $4-\pi\approx0{,}858>0$, como debe ser un valor absoluto.
\end{solucion}

\begin{solucion}{eje:2-vf}
\planteamiento contrastar con el \cref{teo:triangular}, su \cref{cor:triangular-inversa} y la \cref{prop:abs}.
\pasos
\begin{enumerate}[label=(\alph*)]
  \item Falsa: con $a=0$, $b=1$, $\abs{a-b}=1>-1=\abs a-\abs b$. Lo correcto es $\abs{a-b}\ge\abs a-\abs b$.
  \item Verdadera: es el caso de igualdad del \cref{teo:triangular}.
  \item Falsa para $x<0$: $\sqrt{(-2)^2}=2\neq-2$. Lo correcto es $\sqrt{x^2}=\abs x$.
\end{enumerate}
\resultado \textbf{(a) falsa; (b) verdadera; (c) falsa.}
\comprobacion (b): $\abs{3+5}=8=3+5$ con $3\cdot5>0$; $\abs{3-5}=2\neq8$ con $3\cdot(-5)<0$.
\end{solucion}

\begin{solucion}{eje:2-traducir}
\planteamiento un intervalo $(c-r,c+r)$ es $\abs{x-c}<r$, con $c$ el punto medio y $r$ el radio.
\pasos
\begin{enumerate}[label=(\alph*)]
  \item Centro $(2+8)/2=5$, radio $3$: $\abs{x-5}<3$.
  \item Centro $(-1+4)/2=1{,}5$, radio $2{,}5$, extremos incluidos: $\abs{x-1{,}5}\le2{,}5$.
  \item $\abs{x-3}<0{,}01$.
\end{enumerate}
\resultado \textbf{(a) $\abs{x-5}<3$; (b) $\abs{x-1{,}5}\le2{,}5$; (c) $\abs{x-3}<0{,}01$.}
\comprobacion (b): $\abs{-1-1{,}5}=2{,}5$ y $\abs{4-1{,}5}=2{,}5$, los extremos cumplen la igualdad.
\end{solucion}

\begin{solucion}{eje:2-inec-simple}
\planteamiento (a) \cref{prop:abs}(iv); (b) $\abs A=\abs B$ si y solo si $A=B$ o $A=-B$.
\pasos
\begin{enumerate}[label=\arabic*.]
  \item (a) $-7<3x-2<7$; sumando 2: $-5<3x<9$; dividiendo por $3>0$: $-\frac53<x<3$.
  \item (b) Caso $2x-1=x+3$: $x=4$.
  \item Caso $2x-1=-(x+3)$: $3x=-2$, $x=-\frac23$.
\end{enumerate}
\resultado \textbf{(a) $\bigl(-\frac53,3\bigr)$; (b) $x=4$ o $x=-\frac23$.}
\comprobacion (b): $\abs{7}=\abs{7}$ y $\abs{-\frac73}=\abs{\frac73}$.
\end{solucion}

\begin{solucion}{eje:2-inec-casos}
\datos $\abs{x+1}\ge2x-4$.
\planteamiento separar los casos $x\ge-1$ y $x<-1$.
\pasos
\begin{enumerate}[label=\arabic*.]
  \item Si $x\ge-1$: $x+1\ge2x-4$, es decir, $x\le5$. Solución en este caso: $[-1,5]$.
  \item Si $x<-1$: $-x-1\ge2x-4$, es decir, $3\ge3x$, $x\le1$, que se cumple siempre en este caso. Solución: $(-\infty,-1)$.
  \item Unión: $(-\infty,5]$.
\end{enumerate}
\resultado \textbf{$x\in(-\infty,5]$.}
\comprobacion $x=5$: $6\ge6$; $x=6$: $7\ge8$ es falso; $x=-10$: $9\ge-24$.
\end{solucion}

\begin{solucion}{eje:2-inec-suma}
\datos $\abs{x-1}+\abs{x+2}\le5$; puntos de corte $-2$ y $1$.
\planteamiento estudiar los tres intervalos como en el \cref{ej:abs-casos}.
\pasos
\begin{enumerate}[label=\arabic*.]
  \item $x\ge1$: $2x+1\le5$, $x\le2$; da $[1,2]$.
  \item $-2\le x<1$: $3\le5$, siempre; da $[-2,1)$.
  \item $x<-2$: $-2x-1\le5$, $x\ge-3$; da $[-3,-2)$.
  \item Unión: $[-3,2]$.
\end{enumerate}
\resultado \textbf{$[-3,2]$}, que es exactamente el complementario de la solución estricta $(-\infty,-3)\cup(2,\infty)$ de la desigualdad $>5$; los puntos $-3$ y $2$ (suma de distancias igual a~5) pertenecen a ambos conjuntos «$\le5$» y «$\ge5$».
\comprobacion $x=2$: $1+4=5$; $x=-3$: $4+1=5$; $x=0$: $1+2=3$.
\end{solucion}

\begin{solucion}{eje:2-amgm-aplic}
\planteamiento aplicar el \cref{teo:amgm} a números cuyo producto (o suma) es constante.
\pasos
\begin{enumerate}[label=(\alph*)]
  \item $x+\frac4x\ge2\sqrt{x\cdot\frac4x}=4$, con igualdad si $x=\frac4x$, es decir, $x=2$.
  \item Lados $a,b>0$ con $2a+2b=20$, es decir, $a+b=10$. Área $ab\le\bigl(\frac{a+b}2\bigr)^2=25$, con igualdad si $a=b=5$.
  \item $\frac{x+y+z}3\ge\sqrt[3]{xyz}=\sqrt[3]8=2$, luego $x+y+z\ge6$, con igualdad si $x=y=z=2$.
\end{enumerate}
\resultado \textbf{(a) mínimo $4$ en $x=2$; (b) el cuadrado de lado~5, área~25; (c) mínimo $6$ en $x=y=z=2$.}
\comprobacion (a): $x=1$ da $5$ y $x=4$ da $5$, ambos mayores que $4$.
\end{solucion}

\begin{solucion}{eje:2-bernoulli-num}
\planteamiento aplicar el \cref{teo:bernoulli} con $x=0{,}01$, $n=100$ y con $x=-0{,}01$, $n=50$.
\pasos
\begin{enumerate}[label=\arabic*.]
  \item $1{,}01^{100}=(1+0{,}01)^{100}\ge1+100\cdot0{,}01=2$.
  \item $0{,}99^{50}=(1-0{,}01)^{50}\ge1-50\cdot0{,}01=0{,}5$; se puede aplicar porque $-0{,}01\ge-1$.
  \item Valores reales: $1{,}01^{100}\approx2{,}705$ y $0{,}99^{50}\approx0{,}6050$.
\end{enumerate}
\resultado \textbf{$1{,}01^{100}\ge2$ y $0{,}99^{50}\ge0{,}5$}; las cotas son correctas pero no ajustadas.
\comprobacion la cota de Bernoulli es la recta tangente en $x=0$; se aleja del valor real cuando $nx$ no es pequeño.
\end{solucion}

\begin{solucion}{eje:2-cs-num}
\planteamiento (a) evaluar ambos miembros; (b) aplicar el \cref{teo:cs} a $(2,3)$ y $(x,y)$.
\pasos
\begin{enumerate}[label=\arabic*.]
  \item (a) $\sum a_ib_i=4-2+6=8$; $\sum a_i^2=1+4+9=14$; $\sum b_i^2=16+1+4=21$. Se cumple $64\le294$.
  \item (b) $(2x+3y)^2\le(2^2+3^2)(x^2+y^2)=13\cdot13=169$, luego $2x+3y\le13$.
  \item La igualdad en Cauchy--Schwarz se da cuando $(x,y)$ es proporcional a $(2,3)$: $(x,y)=t(2,3)$ con $13t^2=13$; para el máximo, $t=1$.
\end{enumerate}
\resultado \textbf{(a) $64\le294$; (b) máximo $13$, alcanzado en $(2,3)$.}
\comprobacion $2^2+3^2=13$ y $2\cdot2+3\cdot3=13$.
\end{solucion}

\begin{solucion}{eje:2-triangular-n}
\datos $x_1,\dots,x_n\in\R$.
\planteamiento inducción sobre $n$ usando el \cref{teo:triangular} en el paso.
\pasos
\begin{enumerate}[label=\arabic*.]
  \item Caso $n=1$: $\abs{x_1}\le\abs{x_1}$.
  \item Paso: sea $S=x_1+\dots+x_n$ y supóngase $\abs S\le\sum_{i=1}^n\abs{x_i}$.
  \item Por el \cref{teo:triangular}, $\abs{S+x_{n+1}}\le\abs S+\abs{x_{n+1}}$.
  \item Por la hipótesis de inducción, $\abs S+\abs{x_{n+1}}\le\sum_{i=1}^{n+1}\abs{x_i}$.
\end{enumerate}
\resultado \textbf{$\abs{\sum_{i=1}^nx_i}\le\sum_{i=1}^n\abs{x_i}$ para todo $n$.}
\comprobacion $\abs{3-5+1}=1\le9=3+5+1$.
\end{solucion}

\begin{solucion}{eje:2-maxmin}
\planteamiento separar los casos $x\ge y$ y $x<y$.
\pasos
\begin{enumerate}[label=\arabic*.]
  \item Si $x\ge y$: $\abs{x-y}=x-y$; $\frac{x+y+(x-y)}2=x=\max(x,y)$ y $\frac{x+y-(x-y)}2=y=\min(x,y)$.
  \item Si $x<y$: $\abs{x-y}=y-x$; $\frac{x+y+(y-x)}2=y=\max(x,y)$ y $\frac{x+y-(y-x)}2=x=\min(x,y)$.
\end{enumerate}
\resultado \textbf{Las dos fórmulas son ciertas para todos $x,y$.}
\comprobacion $x=2$, $y=7$: $\frac{9+5}2=7$ y $\frac{9-5}2=2$. Además $\max+\min=x+y$, como debe ser.
\end{solucion}

\begin{solucion}{eje:2-cadena-medias}
\datos $0<a<b$.
\planteamiento probar las tres desigualdades por separado.
\pasos
\begin{enumerate}[label=\arabic*.]
  \item $a<\sqrt{ab}$: ambos son positivos y $a^2<ab$ (multiplicando $a<b$ por $a>0$); por el \cref{eje:1-potencia-n}, $a<\sqrt{ab}$.
  \item $\sqrt{ab}<\frac{a+b}2$: por la \cref{prop:amgm2}, con desigualdad estricta porque $a\neq b$.
  \item $\frac{a+b}2<b$: es el \cref{eje:1-epsilon}(a).
\end{enumerate}
\resultado \textbf{$a<\sqrt{ab}<\frac{a+b}2<b$.}
\comprobacion $a=1$, $b=4$: $1<2<2{,}5<4$.
\end{solucion}

\begin{solucion}{eje:2-media-armonica}
\datos $a,b>0$; $G=\sqrt{ab}$, $A=\frac{a+b}2$.
\planteamiento reducir cada desigualdad a la \cref{prop:amgm2} o a un cuadrado no negativo.
\pasos
\begin{enumerate}[label=\arabic*.]
  \item Armónica--geométrica: $\frac{2ab}{a+b}=\frac{G^2}{A}\le\frac{G^2}{G}=G$, porque $A\ge G>0$ (\cref{eje:1-negativo}(c)). Igualdad si y solo si $A=G$, es decir, $a=b$.
  \item Geométrica--aritmética: es la \cref{prop:amgm2}.
  \item Aritmética--cuadrática: ambos miembros son positivos, así que basta comparar cuadrados: $\frac{a^2+b^2}2-\frac{(a+b)^2}4=\frac{2a^2+2b^2-a^2-2ab-b^2}4=\frac{(a-b)^2}4\ge0$, con igualdad si y solo si $a=b$.
\end{enumerate}
\resultado \textbf{$\frac{2ab}{a+b}\le\sqrt{ab}\le\frac{a+b}2\le\sqrt{\frac{a^2+b^2}2}$}, con igualdades si y solo si $a=b$.
\comprobacion $a=1$, $b=4$: $1{,}6\le2\le2{,}5\le2{,}915$.
\end{solucion}

\begin{solucion}{eje:2-bernoulli-fuerte}
\datos $x\ge0$, $n\in\N$.
\planteamiento inducción como en el \cref{teo:bernoulli}, descartando solo términos no negativos.
\pasos
\begin{enumerate}[label=\arabic*.]
  \item $n=1$: $1+x\ge1+x+0$.
  \item Paso: si $(1+x)^n\ge1+nx+\frac{n(n-1)}2x^2$, multiplicando por $1+x>0$:
  \[
  (1+x)^{n+1}\ge1+(n+1)x+\Bigl(n+\frac{n(n-1)}2\Bigr)x^2+\frac{n(n-1)}2x^3.
  \]
  \item $n+\frac{n(n-1)}2=\frac{n(n+1)}2=\frac{(n+1)n}2$ y el término en $x^3$ es $\ge0$; se obtiene el enunciado para $n+1$.
  \item Con $x=1$: $2^n\ge1+n+\frac{n(n-1)}2>\frac{n(n-1)}2$.
\end{enumerate}
\resultado \textbf{$(1+x)^n\ge1+nx+\frac{n(n-1)}2x^2$ para $x\ge0$, y $2^n>\frac{n(n-1)}2$.}
\comprobacion $n=3$, $x=1$: $8\ge1+3+3=7$.
\end{solucion}

\begin{solucion}{eje:2-suma-inversos}
\datos $x_1,\dots,x_n>0$.
\planteamiento aplicar el \cref{teo:cs} a $a_i=\sqrt{x_i}$, $b_i=1/\sqrt{x_i}$.
\pasos
\begin{enumerate}[label=\arabic*.]
  \item $\sum_ia_ib_i=\sum_i1=n$; $\sum_ia_i^2=\sum_ix_i$; $\sum_ib_i^2=\sum_i1/x_i$.
  \item Cauchy--Schwarz: $n^2\le\bigl(\sum_ix_i\bigr)\bigl(\sum_i1/x_i\bigr)$.
  \item Igualdad: si todos los $x_i$ valen $c$, el producto es $nc\cdot\frac nc=n^2$. Recíprocamente, por la desigualdad de las medias, $\sum x_i\ge n\sqrt[n]{\prod x_i}$ y $\sum\frac1{x_i}\ge n\big/\sqrt[n]{\prod x_i}$; si el producto vale $n^2$, ambas son igualdades y, por el \cref{teo:amgm}, todos los $x_i$ son iguales.
\end{enumerate}
\resultado \textbf{$\bigl(\sum x_i\bigr)\bigl(\sum 1/x_i\bigr)\ge n^2$, con igualdad si y solo si todos los $x_i$ coinciden.}
\comprobacion $x=(1,2)$: $3\cdot1{,}5=4{,}5\ge4$.
\end{solucion}
% ---------------------------------------------------------------------
\section{Tema 3}

\begin{solucion}{eje:3-basico}
\datos $P(x)=2x^3-x+5$; raíces pedidas $3$ y $-\frac12$.
\planteamiento evaluar; usar $\deg(PQ)=\deg P+\deg Q$; construir el polinomio a partir de sus factores (\cref{prop:vieta}).
\pasos
\begin{enumerate}[label=\arabic*.]
  \item (a) $P(-1)=2(-1)-(-1)+5=-2+1+5=4$.
  \item (b) $\deg(P^2(x^2+1))=3+3+2=8$.
  \item (c) $(x-3)\bigl(x+\frac12\bigr)=0$; multiplicando por~2 para tener coeficientes enteros: $(x-3)(2x+1)=2x^2-5x-3=0$.
\end{enumerate}
\resultado \textbf{(a) $4$; (b) $8$; (c) $2x^2-5x-3=0$.}
\comprobacion (c) por Vieta: suma $\frac52=3-\frac12$, producto $-\frac32=3\cdot(-\frac12)$.
\end{solucion}

\begin{solucion}{eje:3-discriminante}
\planteamiento calcular $\Delta=b^2-4ac$ (\cref{teo:cuadratica}).
\pasos
\begin{enumerate}[label=(\alph*)]
  \item $\Delta=16-20=-4<0$: ninguna solución real.
  \item $\Delta=144-144=0$: una solución doble, $x=\frac{12}8=\frac32$.
  \item $\Delta=1+24=25>0$: dos soluciones, $x=\frac{-1\pm5}{4}$, es decir, $1$ y $-\frac32$.
\end{enumerate}
\resultado \textbf{(a) ninguna; (b) una (doble); (c) dos.}
\comprobacion (b) $4x^2-12x+9=(2x-3)^2$.
\end{solucion}

\begin{solucion}{eje:3-division}
\datos $P=2x^4-3x^3+0x^2+x-5$, $D=x^2-x+1$.
\planteamiento algoritmo de la demostración del \cref{teo:division}: eliminar en cada paso el término de mayor grado.
\pasos
\begin{enumerate}[label=\arabic*.]
  \item $2x^4/x^2=2x^2$; $P-2x^2D=-x^3-2x^2+x-5$.
  \item $-x^3/x^2=-x$; restando $-xD=-x^3+x^2-x$: queda $-3x^2+2x-5$.
  \item $-3x^2/x^2=-3$; restando $-3D=-3x^2+3x-3$: queda $-x-2$, de grado $1<2$.
\end{enumerate}
\resultado \textbf{$Q=2x^2-x-3$, $R=-x-2$.}
\comprobacion $(x^2-x+1)(2x^2-x-3)=2x^4-3x^3+0x^2+2x-3$; sumando $R$: $2x^4-3x^3+x-5=P$.
\end{solucion}

\begin{solucion}{eje:3-ruffini}
\planteamiento probar divisores del término independiente (\cref{teo:raices-racionales}) y dividir por Ruffini.
\pasos
\begin{enumerate}[label=\arabic*.]
  \item (a) $P(1)=1-6+11-6=0$. Ruffini con $1$: coeficientes $1,-5,6$, es decir, $x^2-5x+6=(x-2)(x-3)$.
  \item (b) Con $t=x^2$: $t^2-5t+4=(t-1)(t-4)$, luego $(x^2-1)(x^2-4)=(x-1)(x+1)(x-2)(x+2)$. Por Ruffini con $x=1$ (coeficientes $1,0,-5,0,4$): cociente $x^3+x^2-4x-4$, resto~0; y así sucesivamente.
\end{enumerate}
\resultado \textbf{(a) $(x-1)(x-2)(x-3)$; (b) $(x-1)(x+1)(x-2)(x+2)$.}
\comprobacion (a) en $x=0$: $(-1)(-2)(-3)=-6$, igual al término independiente.
\end{solucion}

\begin{solucion}{eje:3-resto}
\planteamiento teorema del resto (\cref{teo:resto}): el resto entre $x-a$ es $P(a)$.
\pasos
\begin{enumerate}[label=\arabic*.]
  \item (a) $P(2)=8+4k-8+4=4k+4=0$, luego $k=-1$.
  \item (b) $x+1=x-(-1)$: el resto de $x^{100}-1$ es $(-1)^{100}-1=0$; el de $x^{101}+3$ es $(-1)^{101}+3=2$.
\end{enumerate}
\resultado \textbf{(a) $k=-1$; (b) restos $0$ y $2$.}
\comprobacion (a) $x^3-x^2-4x+4=(x-1)(x^2-4)=(x-1)(x-2)(x+2)$, divisible por $x-2$.
\end{solucion}

\begin{solucion}{eje:3-racionales}
\datos $P(x)=2x^3-3x^2-3x+2$.
\planteamiento por el \cref{teo:raices-racionales}, las raíces racionales son de la forma $p/q$ con $p\mid2$, $q\mid2$: candidatas $\pm1,\pm2,\pm\frac12$.
\pasos
\begin{enumerate}[label=\arabic*.]
  \item $P(1)=-2$, $P(-1)=0$. Ruffini con $-1$: coeficientes $2,-5,2$.
  \item $2x^2-5x+2=0$: $x=\frac{5\pm3}4$, es decir, $2$ y $\frac12$.
\end{enumerate}
\resultado \textbf{Raíces $-1$, $\frac12$ y $2$; $P(x)=(x+1)(2x-1)(x-2)$.}
\comprobacion Vieta: suma $-1+\frac12+2=\frac32=-\frac{-3}2$; producto $-1=(-1)^3\cdot\frac22$.
\end{solucion}

\begin{solucion}{eje:3-vieta}
\datos $x_1+x_2=5$, $x_1x_2=3$ (\cref{prop:vieta}).
\planteamiento expresar cada cantidad mediante la suma $s$ y el producto $p$.
\pasos
\begin{enumerate}[label=\arabic*.]
  \item $x_1^2+x_2^2=s^2-2p=25-6=19$.
  \item $\frac1{x_1}+\frac1{x_2}=\frac{s}{p}=\frac53$.
  \item $x_1^3+x_2^3=s^3-3ps=125-45=80$.
\end{enumerate}
\resultado \textbf{$19$, $\frac53$ y $80$.}
\comprobacion las raíces son $\frac{5\pm\sqrt{13}}2\approx4{,}303$ y $0{,}6972$; $4{,}303^2+0{,}6972^2\approx18{,}52+0{,}486\approx19{,}0$.
\end{solucion}

\begin{solucion}{eje:3-ecuaciones}
\planteamiento (a) cambio $t=x^2$; (b) elevar al cuadrado y comprobar; (c) multiplicar por $(x-1)(x+1)$ con $x\neq\pm1$.
\pasos
\begin{enumerate}[label=\arabic*.]
  \item (a) $t^2-13t+36=0$, $t=4$ o $t=9$; $x=\pm2$ o $x=\pm3$.
  \item (b) $x+5=(x-1)^2=x^2-2x+1$, luego $x^2-3x-4=0$ y $x=4$ o $x=-1$. Comprobación: $x=4$ da $\sqrt9=3=4-1$, válida; $x=-1$ da $\sqrt4=2\neq-2$, extraña.
  \item (c) $(x+1)+2(x-1)=x^2-1$, es decir, $x^2-3x=0$, $x=0$ o $x=3$; ninguna anula los denominadores.
\end{enumerate}
\resultado \textbf{(a) $\pm2,\pm3$; (b) $x=4$; (c) $x=0$ o $x=3$.}
\comprobacion (c) $x=3$: $\frac12+\frac24=1$; $x=0$: $-1+2=1$.
\end{solucion}

\begin{solucion}{eje:3-inecuaciones}
\planteamiento factorizar y hacer la tabla de signos (\cref{ej:tabla-signos}).
\pasos
\begin{enumerate}[label=\arabic*.]
  \item (a) $x^2-x-6=(x-3)(x+2)$; es negativo entre las raíces: $(-2,3)$.
  \item (b) Puntos críticos $-2$, $1$, $3$. Signos de $\frac{(x-1)(x+2)}{x-3}$: en $(-\infty,-2)$, $\frac{(-)(-)}{(-)}<0$; en $(-2,1)$, $\frac{(-)(+)}{(-)}>0$; en $(1,3)$, $\frac{(+)(+)}{(-)}<0$; en $(3,\infty)$, $>0$.
  \item Se incluyen los ceros del numerador ($-2$ y $1$) y se excluye $3$.
\end{enumerate}
\resultado \textbf{(a) $(-2,3)$; (b) $[-2,1]\cup(3,\infty)$.}
\comprobacion (b) $x=0$: $\frac{(-1)(2)}{-3}=\frac23\ge0$; $x=2$: $\frac{1\cdot4}{-1}<0$.
\end{solucion}

\begin{solucion}{eje:3-sistemas}
\planteamiento (a) regla de Cramer (\cref{teo:cramer}) y estudio aparte de $\delta=0$; (b) reducción eliminando primero $z$.
\pasos
\begin{enumerate}[label=\arabic*.]
  \item (a) $\delta=1\cdot k-2\cdot3=k-6$. Si $k\neq6$: $x=\frac{4k-2\cdot6}{k-6}=\frac{4k-12}{k-6}$, $y=\frac{1\cdot6-4\cdot3}{k-6}=\frac{-6}{k-6}$.
  \item Si $k=6$: la segunda ecuación es $3x+6y=6$, es decir, $x+2y=2$, incompatible con $x+2y=4$.
  \item (b) Sumando la primera y la tercera: $2x+3y=8$. Sumando la segunda y la tercera: $3x+y=5$.
  \item De la última, $y=5-3x$; sustituyendo, $2x+15-9x=8$, $x=1$, $y=2$, y de la primera $z=3$.
\end{enumerate}
\resultado \textbf{(a) si $k\neq6$, $(x,y)=\bigl(\frac{4k-12}{k-6},\frac{-6}{k-6}\bigr)$; si $k=6$, incompatible. (b) $(1,2,3)$.}
\comprobacion (a) $k=0$: $(2,1)$, y $2+2=4$, $6+0=6$. (b) coincide con el \cref{ej:gauss3}.
\end{solucion}

\begin{solucion}{eje:3-identidad}
\datos $\deg P,\deg Q\le n$; $P(c_i)=Q(c_i)$ para $n+1$ valores distintos $c_0,\dots,c_n$.
\planteamiento aplicar el \cref{cor:raices} a $R=P-Q$.
\pasos
\begin{enumerate}[label=\arabic*.]
  \item $R=P-Q$ es nulo o tiene grado a lo sumo $n$.
  \item $R(c_i)=0$ para $n+1$ valores distintos.
  \item Si $R$ no fuera nulo, tendría como mucho $n$ raíces (\cref{cor:raices}): contradicción. Luego $R=0$ y $P=Q$.
\end{enumerate}
\resultado \textbf{$P=Q$ coeficiente a coeficiente.}
\comprobacion el número $n+1$ no puede rebajarse: $P=x^2$ y $Q=x$ ($n=2$) coinciden en los dos puntos $0$ y $1$, y son distintos.
\end{solucion}

\begin{solucion}{eje:3-diferencia-potencias}
\datos $n\in\N$.
\planteamiento desarrollar el producto: la suma es telescópica.
\pasos
\begin{enumerate}[label=\arabic*.]
  \item $(x-a)\sum_{k=0}^{n-1}a^kx^{n-1-k}=\sum_{k=0}^{n-1}a^kx^{n-k}-\sum_{k=0}^{n-1}a^{k+1}x^{n-1-k}$.
  \item Con $j=k+1$ en la segunda suma, esta es $\sum_{j=1}^{n}a^jx^{n-j}$. Al restar se cancelan los términos con $1\le j\le n-1$ y queda $x^n-a^n$ (\cref{prop:telescopica}).
  \item Por tanto $x-a$ divide a $x^n-a^n$.
  \item Si $n$ es impar, sustituyendo $a$ por $-a$: $x^n-(-a)^n=x^n+a^n$ es divisible por $x-(-a)=x+a$.
\end{enumerate}
\resultado \textbf{La identidad es cierta; $x-a\mid x^n-a^n$ siempre y $x+a\mid x^n+a^n$ para $n$ impar.}
\comprobacion $n=3$: $x^3+8=(x+2)(x^2-2x+4)$.
\end{solucion}

\begin{solucion}{eje:3-impares}
\datos $b,c$ impares.
\planteamiento por el \cref{teo:raices-racionales} con $a_n=1$, toda raíz racional es entera; se estudia la paridad.
\pasos
\begin{enumerate}[label=\arabic*.]
  \item Si $p/q$ es raíz racional irreducible, $q\mid1$, luego la raíz es un entero $r$.
  \item Si $r$ es par: $r^2$ y $br$ son pares y $c$ impar, así que $r^2+br+c$ es impar.
  \item Si $r$ es impar: $r^2$, $br$ y $c$ son impares y su suma es impar.
  \item En ambos casos $r^2+br+c\neq0$, contradicción.
\end{enumerate}
\resultado \textbf{$x^2+bx+c=0$ no tiene raíces racionales.}
\comprobacion $x^2+3x+1=0$ tiene raíces $\frac{-3\pm\sqrt5}2$, irracionales.
\end{solucion}

\begin{solucion}{eje:3-det-cero}
\datos $ad-bc=0$; algún coeficiente no nulo; por simetría se supone $a\neq0$.
\planteamiento mostrar que la segunda ecuación, en su parte izquierda, es múltiplo de la primera.
\pasos
\begin{enumerate}[label=\arabic*.]
  \item Sea $\lambda=c/a$. Entonces $c=\lambda a$ y $d=\frac{bc}a=\lambda b$.
  \item La segunda ecuación es $\lambda(ax+by)=f$. Si $(x,y)$ resuelve la primera, debe cumplirse $\lambda e=f$.
  \item Si $f\neq\lambda e$, no hay solución.
  \item Si $f=\lambda e$, la segunda ecuación es consecuencia de la primera, cuyas soluciones son $x=\frac{e-by}a$ con $y$ libre: infinitas.
\end{enumerate}
\resultado \textbf{Con $\delta=0$ el sistema es incompatible o compatible indeterminado.}
\comprobacion el sistema del \cref{eje:3-sistemas} con $k=6$ tiene $\lambda=3$, $\lambda e=12\neq6=f$: incompatible.
\end{solucion}

\begin{solucion}{eje:3-cuadratica-positiva}
\datos $a>0$, $f(x)=ax^2+bx+c$.
\planteamiento completar el cuadrado como en el \cref{teo:cuadratica}.
\pasos
\begin{enumerate}[label=\arabic*.]
  \item $f(x)=a\bigl(x+\frac b{2a}\bigr)^2+c-\frac{b^2}{4a}$.
  \item Como $a>0$ y un cuadrado es $\ge0$, $f(x)\ge c-\frac{b^2}{4a}$, con igualdad si y solo si $x=-\frac b{2a}$.
  \item $f\ge0$ en todo $\R$ si y solo si su mínimo es $\ge0$: $c-\frac{b^2}{4a}\ge0$; multiplicando por $4a>0$, $4ac-b^2\ge0$.
\end{enumerate}
\resultado \textbf{Mínimo $c-\frac{b^2}{4a}$ en $x=-\frac b{2a}$; $f\ge0$ si y solo si $b^2-4ac\le0$.}
\comprobacion $f=x^2-2x+5$: mínimo en $x=1$ con valor $4=5-\frac44$, y $\Delta=-16\le0$.
\end{solucion}
% ---------------------------------------------------------------------
\section{Tema 4}

\begin{solucion}{eje:4-notacion}
\planteamiento encontrar el término general y los límites del índice.
\pasos
\begin{enumerate}[label=(\alph*)]
  \item Impares $2k-1$; $2k-1=99$ da $k=50$: $\sum_{k=1}^{50}(2k-1)$.
  \item Potencias $2^k$; $2^{10}=1024$: $\sum_{k=1}^{10}2^k$.
  \item Términos $(-1)^k/k$ para $k=2,\dots,100$: $\sum_{k=2}^{100}\frac{(-1)^k}{k}$.
\end{enumerate}
\resultado \textbf{(a) $\sum_{k=1}^{50}(2k-1)$; (b) $\sum_{k=1}^{10}2^k$; (c) $\sum_{k=2}^{100}(-1)^k/k$.}
\comprobacion (a) por el \cref{teo:aritmetica}, $\frac{50(1+99)}2=2500=50^2$, la conocida suma de los primeros impares.
\end{solucion}

\begin{solucion}{eje:4-calcular}
\planteamiento evaluar término a término y aplicar los convenios de la suma vacía y de la suma de una constante.
\pasos
\begin{enumerate}[label=(\alph*)]
  \item $1+4+9+16=30$.
  \item $\frac12\cdot\frac23\cdot\frac34\cdot\frac45\cdot\frac56=\frac16$: cada numerador se cancela con el denominador anterior.
  \item El límite superior es menor que el inferior: suma vacía, que vale~$0$.
  \item $\sum_{k=1}^{10}7=10\cdot7=70$ (\cref{prop:sumas}(iii)).
\end{enumerate}
\resultado \textbf{(a) $30$; (b) $1/6$; (c) $0$; (d) $70$.}
\comprobacion (a) con el \cref{teo:potencias}: $\frac{4\cdot5\cdot9}6=30$.
\end{solucion}

\begin{solucion}{eje:4-errores}
\planteamiento comparar con la \cref{prop:sumas} y la \cref{prop:producto-sumas}.
\pasos
\begin{enumerate}[label=(\alph*)]
  \item $\sum_{k=1}^n3=3n$: se suman $n$ treses.
  \item $(a_1+a_2)(b_1+b_2)=a_1b_1+a_2b_2+a_1b_2+a_2b_1$: faltan los términos cruzados. Contraejemplo: $a=b=(1,1)$ da $2$ frente a $4$.
  \item $(a_1+a_2)^2=a_1^2+a_2^2+2a_1a_2$: falta el doble producto.
\end{enumerate}
\resultado \textbf{(a) vale $3n$; (b) y (c) olvidan los términos cruzados.}
\comprobacion (c) con $a=(1,2)$: $9\neq5$, y $5+2\cdot2=9$.
\end{solucion}

\begin{solucion}{eje:4-aritmeticas}
\planteamiento linealidad y \cref{teo:aritmetica}.
\pasos
\begin{enumerate}[label=\arabic*.]
  \item (a) $\sum(3k-2)=3\cdot\frac{100\cdot101}2-2\cdot100=15150-200=14950$.
  \item (b) El primer múltiplo de~7 mayor que 100 es $105=7\cdot15$; el último menor que 1000 es $994=7\cdot142$.
  \item Hay $142-15+1=128$ términos. Suma: $\frac{128(105+994)}2=64\cdot1099=70336$.
\end{enumerate}
\resultado \textbf{(a) $14950$; (b) $70336$.}
\comprobacion (b) $7\sum_{k=15}^{142}k=7\bigl(\frac{142\cdot143}2-\frac{14\cdot15}2\bigr)=7(10153-105)=70336$.
\end{solucion}

\begin{solucion}{eje:4-geometricas}
\planteamiento \cref{teo:geometrica}.
\pasos
\begin{enumerate}[label=(\alph*)]
  \item $3\sum_{k=0}^92^k=3\cdot\frac{2^{10}-1}{2-1}=3\cdot1023=3069$.
  \item $\sum_{k=1}^n(\frac12)^k=\sum_{k=0}^n(\frac12)^k-1=\frac{1-2^{-(n+1)}}{1/2}-1=2-2^{-n}-1=1-2^{-n}$.
  \item Razón $r=-x\neq1$: $\sum_{k=0}^n(-x)^k=\frac{1-(-x)^{n+1}}{1+x}$.
\end{enumerate}
\resultado \textbf{(a) $3069$; (b) $1-2^{-n}$; (c) $\dfrac{1-(-x)^{n+1}}{1+x}$.}
\comprobacion (b) $n=2$: $\frac12+\frac14=\frac34=1-\frac14$. (c) $n=1$: $1-x=\frac{1-x^2}{1+x}$.
\end{solucion}

\begin{solucion}{eje:4-telescopicas}
\planteamiento escribir cada término como diferencia y aplicar la \cref{prop:telescopica}.
\pasos
\begin{enumerate}[label=\arabic*.]
  \item (a) $\frac1{k(k+2)}=\frac12\bigl(\frac1k-\frac1{k+2}\bigr)$. Al sumar sobreviven $\frac11+\frac12$ con signo $+$ y $\frac1{n+1}+\frac1{n+2}$ con signo $-$: $\frac12\bigl(\frac32-\frac1{n+1}-\frac1{n+2}\bigr)$.
  \item Para $n=10$: $\frac12\bigl(\frac32-\frac1{11}-\frac1{12}\bigr)=\frac12\cdot\frac{198-12-11}{132}=\frac{175}{264}$.
  \item (b) $\ln(1+\frac1k)=\ln(k+1)-\ln k$; suma $\ln(n+1)-\ln1=\ln(n+1)$.
  \item (c) $\frac1{\sqrt k+\sqrt{k+1}}=\frac{\sqrt{k+1}-\sqrt k}{(k+1)-k}=\sqrt{k+1}-\sqrt k$; suma $\sqrt{100}-\sqrt1=9$.
\end{enumerate}
\resultado \textbf{(a) $\frac34-\frac{2n+3}{2(n+1)(n+2)}$, que vale $\frac{175}{264}$ para $n=10$; (b) $\ln(n+1)$; (c) $9$.}
\comprobacion (a) $n=1$: $\frac13$ y $\frac34-\frac5{12}=\frac4{12}=\frac13$.
\end{solucion}

\begin{solucion}{eje:4-indice}
\planteamiento (a) cambio $j=k-2$ (\cref{prop:sumas}(iv)); (b) $k(k+1)=k^2+k$ y \cref{teo:potencias}.
\pasos
\begin{enumerate}[label=\arabic*.]
  \item (a) Si $k=3$, $j=1$; si $k=n+2$, $j=n$: $\sum_{j=1}^nj^2=\frac{n(n+1)(2n+1)}6$.
  \item (b) $\sum_{k=1}^n(k^2+k)=\frac{n(n+1)(2n+1)}6+\frac{n(n+1)}2=\frac{n(n+1)(2n+1+3)}6=\frac{n(n+1)(n+2)}3$.
  \item Para $n=10$: $\frac{10\cdot11\cdot12}3=440$.
\end{enumerate}
\resultado \textbf{(a) $\sum_{j=1}^nj^2=\frac{n(n+1)(2n+1)}6$; (b) $\frac{n(n+1)(n+2)}3$, que vale $440$ para $n=10$.}
\comprobacion (b) $n=10$ directamente: $385+55=440$.
\end{solucion}

\begin{solucion}{eje:4-dobles}
\planteamiento sumar primero el índice interior; en (c), reescribir el dominio triangular.
\pasos
\begin{enumerate}[label=\arabic*.]
  \item (a) $\sum_{j=1}^4(i+2j)=4i+2\cdot10=4i+20$; $\sum_{i=1}^3(4i+20)=24+60=84$.
  \item (b) $\sum_{j=1}^ij=\frac{i(i+1)}2$; $\sum_{i=1}^n\frac{i(i+1)}2=\frac{n(n+1)(n+2)}6$ por el \cref{eje:4-indice}.
  \item (c) El dominio es $1\le i\le j\le n$: $\sum_{j=1}^n\sum_{i=1}^j\frac1j=\sum_{j=1}^n\frac jj=n$.
\end{enumerate}
\resultado \textbf{(a) $84$; (b) $\frac{n(n+1)(n+2)}6$; (c) $n$.}
\comprobacion (c) $n=2$: $(\frac11+\frac12)+\frac12=2$.
\end{solucion}

\begin{solucion}{eje:4-binomio}
\planteamiento \cref{teo:binomio} con los valores adecuados de $a$ y $b$.
\pasos
\begin{enumerate}[label=\arabic*.]
  \item (a) Con $a=2x$, $b=-1$ y $\binom5k=1,5,10,10,5,1$: $32x^5-80x^4+80x^3-40x^2+10x-1$.
  \item (b) Término general $\binom8kx^{8-k}(2/x)^k=\binom8k2^kx^{8-2k}$; $8-2k=4$ da $k=2$: $\binom82\cdot4=28\cdot4=112$.
  \item (c) $\sum_k\binom6k1^{6-k}2^k=(1+2)^6=729$.
\end{enumerate}
\resultado \textbf{(a) $32x^5-80x^4+80x^3-40x^2+10x-1$; (b) $112$; (c) $729$.}
\comprobacion (a) en $x=1$: $32-80+80-40+10-1=1=(2-1)^5$.
\end{solucion}

\begin{solucion}{eje:4-productos}
\planteamiento factorizar y cancelar (producto telescópico); exponentes con el \cref{teo:aritmetica}.
\pasos
\begin{enumerate}[label=\arabic*.]
  \item (a) $1-\frac1{k^2}=\frac{k-1}k\cdot\frac{k+1}k$. Entonces $\prod_{k=2}^n\frac{k-1}k=\frac1n$ y $\prod_{k=2}^n\frac{k+1}k=\frac{n+1}2$; el producto es $\frac{n+1}{2n}$, que para $n=10$ vale $\frac{11}{20}$.
  \item (b) $\prod2^k=2^{\sum k}=2^{n(n+1)/2}$.
  \item (c) $\ln\prod\frac{k+1}k=\sum(\ln(k+1)-\ln k)=\ln(n+1)$.
\end{enumerate}
\resultado \textbf{(a) $\frac{n+1}{2n}$ ($\frac{11}{20}$ para $n=10$); (b) $2^{n(n+1)/2}$; (c) $\ln(n+1)$.}
\comprobacion (a) $n=2$: $1-\frac14=\frac34=\frac{2+1}{2\cdot2}$.
\end{solucion}

\begin{solucion}{eje:4-cubos}
\datos $S_3(n)=\sum_{k=1}^nk^3$.
\planteamiento inducción sobre $n$ (\cref{teo:induccion}).
\pasos
\begin{enumerate}[label=\arabic*.]
  \item $n=1$: $1=\bigl(\frac{1\cdot2}2\bigr)^2$.
  \item Paso: si $S_3(n)=\frac{n^2(n+1)^2}4$, entonces $S_3(n+1)=\frac{n^2(n+1)^2}4+(n+1)^3=\frac{(n+1)^2(n^2+4n+4)}4$.
  \item Como $n^2+4n+4=(n+2)^2$, $S_3(n+1)=\bigl(\frac{(n+1)(n+2)}2\bigr)^2$, que es el enunciado para $n+1$.
\end{enumerate}
\resultado \textbf{$\sum_{k=1}^nk^3=\bigl(\frac{n(n+1)}2\bigr)^2$ para todo $n$.}
\comprobacion $n=4$: $1+8+27+64=100=10^2$.
\end{solucion}

\begin{solucion}{eje:4-perturbacion}
\datos $S=\sum_{k=0}^nkr^k$, $r\neq1$.
\planteamiento método de perturbación: calcular $S_{n+1}$ de dos formas~\cite[§2.3]{gkp}.
\pasos
\begin{enumerate}[label=\arabic*.]
  \item Separando el último término: $S+(n+1)r^{n+1}=\sum_{k=0}^{n+1}kr^k$.
  \item Separando el primero ($k=0$, que vale~0) y cambiando $k=j+1$: $\sum_{k=0}^{n+1}kr^k=\sum_{j=0}^n(j+1)r^{j+1}=rS+r\sum_{j=0}^nr^j$.
  \item Igualando y usando el \cref{teo:geometrica}: $S(1-r)=r\frac{1-r^{n+1}}{1-r}-(n+1)r^{n+1}$.
  \item Dividiendo por $1-r$ y reduciendo a común denominador:
  \[
  S=\frac{r-r^{n+2}-(n+1)r^{n+1}+(n+1)r^{n+2}}{(1-r)^2}=\frac{r-(n+1)r^{n+1}+nr^{n+2}}{(1-r)^2}.
  \]
\end{enumerate}
\resultado \textbf{$\sum_{k=0}^nkr^k=\dfrac{r-(n+1)r^{n+1}+nr^{n+2}}{(1-r)^2}$.}
\comprobacion $n=3$, $r=2$: $0+2+8+24=34$ y $\frac{2-4\cdot16+3\cdot32}{1}=34$.
\end{solucion}

\begin{solucion}{eje:4-kbinom}
\datos $1\le k\le n$.
\planteamiento desarrollar los factoriales; después, linealidad y \cref{cor:binomio}.
\pasos
\begin{enumerate}[label=\arabic*.]
  \item $k\binom nk=\frac{k\,n!}{k!(n-k)!}=\frac{n\,(n-1)!}{(k-1)!\,((n-1)-(k-1))!}=n\binom{n-1}{k-1}$.
  \item El término $k=0$ vale~0, así que $\sum_{k=0}^nk\binom nk=n\sum_{k=1}^n\binom{n-1}{k-1}$.
  \item Con $j=k-1$: $n\sum_{j=0}^{n-1}\binom{n-1}j=n\,2^{n-1}$.
\end{enumerate}
\resultado \textbf{$k\binom nk=n\binom{n-1}{k-1}$ y $\sum_kk\binom nk=n2^{n-1}$.}
\comprobacion $n=3$: $0+3+6+3=12=3\cdot4$.
\end{solucion}

\begin{solucion}{eje:4-basilea}
\datos $n\in\N$.
\planteamiento inducción; el paso se reduce a una desigualdad entre fracciones.
\pasos
\begin{enumerate}[label=\arabic*.]
  \item $n=1$: $1\le2-1$.
  \item Paso: si $\sum_{k=1}^n\frac1{k^2}\le2-\frac1n$, entonces $\sum_{k=1}^{n+1}\frac1{k^2}\le2-\frac1n+\frac1{(n+1)^2}$.
  \item Basta ver que $-\frac1n+\frac1{(n+1)^2}\le-\frac1{n+1}$, es decir, $\frac1{n+1}+\frac1{(n+1)^2}\le\frac1n$.
  \item El miembro izquierdo es $\frac{n+2}{(n+1)^2}$, y $\frac{n+2}{(n+1)^2}\le\frac1n$ equivale a $n^2+2n\le n^2+2n+1$, que es cierto.
\end{enumerate}
\resultado \textbf{$\sum_{k=1}^n\frac1{k^2}\le2-\frac1n$ para todo $n$}; en particular la suma está acotada por~2.
\comprobacion $n=3$: $1+\frac14+\frac19=1{,}3611\le1{,}6667$.
\end{solucion}

\begin{solucion}{eje:4-varianza}
\datos $a_1,\dots,a_n\in\R$; $S_1=\sum a_i$, $S_2=\sum a_i^2$.
\planteamiento desarrollar el cuadrado y aplicar la \cref{prop:producto-sumas}.
\pasos
\begin{enumerate}[label=\arabic*.]
  \item $\sum_i\sum_j(a_i-a_j)^2=\sum_i\sum_ja_i^2-2\sum_i\sum_ja_ia_j+\sum_i\sum_ja_j^2$.
  \item El primer término es $nS_2$ (no depende de $j$), el tercero también es $nS_2$, y el central es $2S_1^2$.
  \item Total: $2nS_2-2S_1^2$.
  \item El miembro izquierdo es suma de cuadrados, luego $\ge0$: $S_1^2\le nS_2$.
  \item Igualdad si y solo si todos los $(a_i-a_j)^2$ son nulos, es decir, todos los $a_i$ son iguales.
\end{enumerate}
\resultado \textbf{$\sum_{i,j}(a_i-a_j)^2=2nS_2-2S_1^2\ge0$ y $\bigl(\sum a_i\bigr)^2\le n\sum a_i^2$.}
\comprobacion $a=(1,2,3)$: los pares ordenados dan $2(1+4+1)=12$, y $2\cdot3\cdot14-2\cdot36=84-72=12$. Dividiendo por $2n^2$ se obtiene que la varianza $\frac{S_2}n-\bigl(\frac{S_1}n\bigr)^2$ es no negativa (módulo~5).
\end{solucion}
% ---------------------------------------------------------------------
\section{Tema 5}

\begin{solucion}{eje:5-calcular}
\planteamiento escribir el argumento como potencia de la base (\cref{def:log}).
\pasos
\begin{enumerate}[label=\arabic*.]
  \item $64=2^6$, luego $\log_264=6$. $\frac1{81}=3^{-4}$, luego $\log_3\frac1{81}=-4$. $0{,}001=10^{-3}$, luego $\log_{10}0{,}001=-3$.
  \item $\ln e^7=7$.
  \item $8=2^3=(4^{1/2})^3=4^{3/2}$, luego $\log_48=\frac32$.
  \item $10^{\log_{10}5}=5$ por definición.
\end{enumerate}
\resultado \textbf{$6$, $-4$, $-3$, $7$, $\frac32$ y $5$.}
\comprobacion $\log_48$ por cambio de base: $\frac{\log_28}{\log_24}=\frac32$.
\end{solucion}

\begin{solucion}{eje:5-vf}
\planteamiento contrastar con el \cref{teo:log}.
\pasos
\begin{enumerate}[label=(\alph*)]
  \item Falsa: con $a=b=1$, $\ln2\neq0=\ln1+\ln1$.
  \item Verdadera: $\log_2(8^x)=x\log_28=3x$.
  \item Verdadera: $\ln\frac1x=\ln1-\ln x=-\ln x$.
  \item Falsa: con $x=e$, $(\ln e)^2=1\neq2$. Lo correcto es $\ln(x^2)=2\ln x$.
\end{enumerate}
\resultado \textbf{(a) falsa; (b) verdadera; (c) verdadera; (d) falsa.}
\comprobacion (d) con $x=1$ ambos miembros valen 0: un caso particular no basta para afirmar una identidad.
\end{solucion}

\begin{solucion}{eje:5-unidades}
\planteamiento $1\ \text{nat}=1/\ln2$ bits; información $-\log_2p$ bits; para numerar $K$ objetos hacen falta $\lceil\log_2K\rceil$ bits (\cref{prop:cifras}).
\pasos
\begin{enumerate}[label=(\alph*)]
  \item $8/\ln2=8/0{,}6931\approx11{,}54$ bits.
  \item $-\log_2(2^{-10})=10$ bits; en nats, $10\ln2\approx6{,}931$.
  \item $2^{15}=32\,768<50\,257\le65\,536=2^{16}$: hacen falta $16$ bits.
\end{enumerate}
\resultado \textbf{(a) $11{,}54$ bits; (b) $10$ bits $=6{,}931$ nats; (c) $16$ bits.}
\comprobacion $\log_250\,257\approx15{,}62$, cuyo entero superior es~16.
\end{solucion}

\begin{solucion}{eje:5-simplificar}
\planteamiento \cref{teo:log}.
\pasos
\begin{enumerate}[label=(\alph*)]
  \item $\log_2\frac{12}3+\log_22^{-3}=\log_24-3=2-3=-1$.
  \item $\frac12\ln16=\ln4$; $\ln\frac e4=1-\ln4$; suma $1$.
  \item $\log_23\cdot\log_38=\log_23\cdot\frac{\log_28}{\log_23}=3$.
  \item $\log_530=\frac{\ln30}{\ln5}\approx\frac{3{,}401}{1{,}609}\approx2{,}113$.
\end{enumerate}
\resultado \textbf{(a) $-1$; (b) $1$; (c) $3$; (d) $\ln30/\ln5\approx2{,}113$.}
\comprobacion (d) $5^2=25<30<125=5^3$, luego el valor está entre 2 y 3.
\end{solucion}

\begin{solucion}{eje:5-exponenciales}
\planteamiento (a) tomar $\ln$; (b) y (c) cambio de variable $t=2^x$ o $t=e^x$, con $t>0$.
\pasos
\begin{enumerate}[label=\arabic*.]
  \item (a) $(x+1)\ln2=x\ln3$, luego $x(\ln3-\ln2)=\ln2$ y $x=\frac{\ln2}{\ln3-\ln2}\approx1{,}710$.
  \item (b) $t^2-3t-4=(t-4)(t+1)=0$; $t=-1$ se descarta porque $t>0$; $2^x=4$, $x=2$.
  \item (c) $t^2-5t+6=(t-2)(t-3)=0$; $x=\ln2$ o $x=\ln3$.
\end{enumerate}
\resultado \textbf{(a) $x=\frac{\ln2}{\ln(3/2)}\approx1{,}710$; (b) $x=2$; (c) $x=\ln2\approx0{,}6931$ o $x=\ln3\approx1{,}099$.}
\comprobacion (a) $2^{2{,}710}\approx6{,}541\approx3^{1{,}710}$.
\end{solucion}

\begin{solucion}{eje:5-logaritmicas}
\planteamiento agrupar logaritmos, igualar argumentos y comprobar el dominio (\cref{ej:ecuacion-log}).
\pasos
\begin{enumerate}[label=\arabic*.]
  \item (a) Dominio $x>2$. $\log_2(x(x-2))=3$, $x^2-2x-8=0$, $x=4$ o $x=-2$; solo $x=4$ está en el dominio.
  \item (b) Dominio $x>1$. $\ln\frac{x+1}{x-1}=\ln3$, $x+1=3x-3$, $x=2$.
\end{enumerate}
\resultado \textbf{(a) $x=4$; (b) $x=2$.}
\comprobacion (a) $\log_24+\log_22=2+1=3$; (b) $\ln3-\ln1=\ln3$.
\end{solucion}

\begin{solucion}{eje:5-inecuaciones}
\planteamiento usar la monotonía (\cref{teo:log}(iv)): con base menor que 1 el sentido se invierte.
\pasos
\begin{enumerate}[label=\arabic*.]
  \item (a) $0{,}5^x=2^{-x}>8=2^3$; como $2^t$ es creciente, $-x>3$, es decir, $x<-3$.
  \item (b) Dominio $2x-1>0$, $x>\frac12$. Como $\log_3$ es creciente, $2x-1\le3^2=9$, $x\le5$.
\end{enumerate}
\resultado \textbf{(a) $(-\infty,-3)$; (b) $\bigl(\frac12,5\bigr]$.}
\comprobacion (a) $x=-4$: $0{,}5^{-4}=16>8$. (b) $x=5$: $\log_39=2$.
\end{solucion}

\begin{solucion}{eje:5-lse}
\datos $z=(2,1,0,-1)$, $M=2$, $n=4$.
\planteamiento cálculo estable: $\LSE(z)=M+\ln\sum_ie^{z_i-M}$ (\cref{prop:lse}).
\pasos
\begin{enumerate}[label=\arabic*.]
  \item $\sum_ie^{z_i-2}=1+e^{-1}+e^{-2}+e^{-3}\approx1+0{,}3679+0{,}1353+0{,}0498=1{,}553$.
  \item $\LSE(z)\approx2+\ln1{,}553\approx2+0{,}4402=2{,}440$.
  \item $\softmax(z)_i=e^{z_i-2}/1{,}553$: $(0{,}6439;\ 0{,}2369;\ 0{,}0871;\ 0{,}0321)$.
  \item Cotas: $2\le2{,}440\le2+\ln4\approx3{,}386$.
  \item (b) $M=1000$: $1000+\ln(1+e^{-1}+e^{-2})\approx1000+\ln1{,}503=1000{,}408$.
\end{enumerate}
\resultado \textbf{(a) $\LSE\approx2{,}440$, softmax $\approx(0{,}6439;0{,}2369;0{,}0871;0{,}0321)$; (b) $\approx1000{,}408$.}
\comprobacion las componentes de la softmax suman $1{,}0000$.
\end{solucion}

\begin{solucion}{eje:5-perplejidad}
\datos log-verosimilitud media $-2{,}3$ nats por token, es decir, $H=2{,}3$ nats.
\planteamiento cambio de base y definición de perplejidad $e^{H}$ ($H$ en nats) o $2^{H}$ ($H$ en bits).
\pasos
\begin{enumerate}[label=\arabic*.]
  \item (a) $2{,}3/\ln2\approx3{,}318$ bits por token.
  \item (b) $e^{2{,}3}\approx9{,}974$; $2^{3{,}318}=e^{3{,}318\ln2}=e^{2{,}3}$, el mismo valor.
  \item (c) Si cada token tiene probabilidad $q$, la log-verosimilitud media es $\ln q=-2{,}3$, luego $q=e^{-2{,}3}\approx0{,}1003$.
\end{enumerate}
\resultado \textbf{(a) $3{,}318$ bits/token; (b) perplejidad $\approx9{,}97$; (c) $q\approx0{,}100$.}
\comprobacion la perplejidad es $1/q$: $1/0{,}1003\approx9{,}97$, el modelo duda como si eligiera entre unos 10 tokens equiprobables.
\end{solucion}

\begin{solucion}{eje:5-ley-potencia}
\datos $L(N)=(N_c/N)^{0{,}076}$.
\planteamiento cociente $L(\lambda N)/L(N)=\lambda^{-\alpha}$ y logaritmos.
\pasos
\begin{enumerate}[label=\arabic*.]
  \item (a) $10^{-0{,}076}=e^{-0{,}076\ln10}\approx e^{-0{,}175}\approx0{,}8395$.
  \item (b) $\lambda^{-0{,}076}=\frac12$, luego $\lambda=2^{1/0{,}076}=e^{\ln2/0{,}076}\approx e^{9{,}120}\approx9{,}14\cdot10^3$.
  \item (c) $\log L=\alpha\log N_c-\alpha\log N$: pendiente $-0{,}076$.
\end{enumerate}
\resultado \textbf{(a) $\times0{,}8395$ (baja un $16\,\%$); (b) unas $9{,}14\cdot10^3$ veces más; (c) pendiente $-0{,}076$.}
\comprobacion (b) $\log_{10}\lambda=\log_{10}2/0{,}076\approx3{,}961$, es decir, casi cuatro órdenes de magnitud.
\end{solucion}

\begin{solucion}{eje:5-log23}
\datos $\log_23>0$ porque $3>1$.
\planteamiento reducción al absurdo y paridad.
\pasos
\begin{enumerate}[label=\arabic*.]
  \item Si $\log_23=p/q$ con $p,q\in\N$, entonces $2^{p/q}=3$ y, elevando a $q$, $2^p=3^q$.
  \item $2^p$ es par (porque $p\ge1$) y $3^q$ es impar (producto de impares).
  \item Contradicción.
\end{enumerate}
\resultado \textbf{$\log_23$ es irracional.}
\comprobacion $\log_23\approx1{,}585$; las aproximaciones racionales $3/2$ y $8/5$ dan $2^3=8\neq9=3^2$ y $2^8=256\neq243=3^5$.
\end{solucion}

\begin{solucion}{eje:5-temperatura}
\datos $z\in\R^n$, $M=\max_iz_i$, $T>0$.
\planteamiento aplicar la \cref{prop:lse}(ii) al vector $z/T$, cuyo máximo es $M/T$.
\pasos
\begin{enumerate}[label=\arabic*.]
  \item $\frac MT\le\LSE(z/T)\le\frac MT+\ln n$.
  \item Multiplicando por $T>0$: $M\le T\LSE(z/T)\le M+T\ln n$.
  \item Cuando $T\to0^+$, $T\ln n\to0$; por el criterio del sándwich, $T\LSE(z/T)\to M$.
\end{enumerate}
\resultado \textbf{$M\le T\LSE(z/T)\le M+T\ln n$, y el límite es $M$}: con temperatura baja, log-suma-exp se comporta como el máximo y la softmax se concentra en el logit mayor.
\comprobacion $z=(1,0)$, $T=0{,}1$: $0{,}1\ln(e^{10}+1)\approx1{,}0000045$, entre $1$ y $1+0{,}1\ln2\approx1{,}069$.
\end{solucion}

\begin{solucion}{eje:5-lse-media}
\datos $z\in\R^n$; $\bar z=\frac1n\sum z_i$.
\planteamiento aplicar el \cref{teo:amgm} a los números positivos $e^{z_i}$.
\pasos
\begin{enumerate}[label=\arabic*.]
  \item $\frac1n\sum_ie^{z_i}\ge\sqrt[n]{\prod_ie^{z_i}}=e^{\frac1n\sum_iz_i}=e^{\bar z}$.
  \item Multiplicando por $n$ y tomando $\ln$ (creciente): $\LSE(z)\ge\ln n+\bar z$.
  \item Igualdad si y solo si todos los $e^{z_i}$ son iguales, es decir, todos los $z_i$ son iguales.
\end{enumerate}
\resultado \textbf{$\LSE(z)\ge\bar z+\ln n$, con igualdad si y solo si $z$ es constante.}
\comprobacion con $z=(2,1,0,-1)$: $\bar z+\ln4\approx0{,}5+1{,}386=1{,}886\le2{,}440$.
\end{solucion}

\begin{solucion}{eje:5-gibbs}
\datos $p_i,q_i>0$, $\sum p_i=\sum q_i=1$.
\planteamiento aplicar $\ln x\le x-1$ (\cref{cor:lnx}) con $x=q_i/p_i$ y sumar con pesos $p_i$.
\pasos
\begin{enumerate}[label=\arabic*.]
  \item $\ln\frac{q_i}{p_i}\le\frac{q_i}{p_i}-1$; multiplicando por $p_i>0$: $p_i\ln\frac{q_i}{p_i}\le q_i-p_i$.
  \item Sumando en $i$: $\sum_ip_i\ln\frac{q_i}{p_i}\le\sum_iq_i-\sum_ip_i=1-1=0$.
  \item Como $\ln\frac{q_i}{p_i}=\ln q_i-\ln p_i$, la desigualdad es $\sum_ip_i\ln q_i\le\sum_ip_i\ln p_i$; cambiando de signo se obtiene la forma del enunciado.
\end{enumerate}
\resultado \textbf{$-\sum_ip_i\ln q_i\ge-\sum_ip_i\ln p_i$}: la entropía cruzada es al menos la entropía (desigualdad de Gibbs).
\comprobacion $p=(\frac12,\frac12)$, $q=(\frac14,\frac34)$: $-\frac12(\ln\frac14+\ln\frac34)\approx0{,}8370\ge\ln2\approx0{,}6931$.
\end{solucion}

\begin{solucion}{eje:5-e-creciente}
\datos $n\in\N$; los $n+1$ números $a_1=\dots=a_n=1+\frac1n$ y $a_{n+1}=1$.
\planteamiento \cref{teo:amgm} con desigualdad estricta, porque los números no son todos iguales.
\pasos
\begin{enumerate}[label=\arabic*.]
  \item Media aritmética: $\frac{n(1+\frac1n)+1}{n+1}=\frac{n+2}{n+1}=1+\frac1{n+1}$.
  \item Media geométrica: $\bigl((1+\frac1n)^n\cdot1\bigr)^{1/(n+1)}$.
  \item Como $1+\frac1n\neq1$, la desigualdad es estricta: $\bigl(1+\frac1n\bigr)^{n/(n+1)}<1+\frac1{n+1}$.
  \item Elevando a $n+1$ (números positivos, \cref{eje:1-potencia-n}): $\bigl(1+\frac1n\bigr)^n<\bigl(1+\frac1{n+1}\bigr)^{n+1}$.
\end{enumerate}
\resultado \textbf{La sucesión $(1+\frac1n)^n$ es estrictamente creciente}, lo que es un paso para probar que converge al número~$e$.
\comprobacion $n=1,2,3,4$: $2<2{,}25<2{,}370<2{,}441$.
\end{solucion}
% ---------------------------------------------------------------------
\section{Tema 6}

\begin{solucion}{eje:6-a-decimal}
\planteamiento sumar cada cifra por su potencia de la base (\cref{teo:base}).
\pasos
\begin{enumerate}[label=\arabic*.]
  \item $(1011\,0110)_2=128+32+16+4+2=182$.
  \item $\texttt{0x2F}=2\cdot16+15=47$; $\texttt{0xFF}=15\cdot16+15=255$.
  \item $(777)_8=7\cdot64+7\cdot8+7=511$.
\end{enumerate}
\resultado \textbf{$182$, $47$, $255$ y $511$.}
\comprobacion $(1011\,0110)_2$ agrupado en cuartetos es $\texttt{B}\,\texttt{6}$, y $\texttt{0xB6}=11\cdot16+6=182$; $(777)_8=8^3-1=511$, como predice la \cref{prop:cifras}.
\end{solucion}

\begin{solucion}{eje:6-contar}
\planteamiento \cref{prop:cifras} y \cref{teo:c2}.
\pasos
\begin{enumerate}[label=\arabic*.]
  \item (a) $2^8=256$, $2^{16}=65\,536$, $2^{32}=4\,294\,967\,296$.
  \item (b) Sin signo: $[0,255]$. Complemento a dos: $[-2^7,2^7-1]=[-128,127]$.
  \item (c) Cada cifra hexadecimal son 4 bits: $32/4=8$ cifras.
\end{enumerate}
\resultado \textbf{(a) $256$, $65\,536$, $4\,294\,967\,296$; (b) $[0,255]$ y $[-128,127]$; (c) $8$.}
\comprobacion ambos rangos de (b) contienen $256$ enteros.
\end{solucion}

\begin{solucion}{eje:6-vf}
\planteamiento \cref{prop:fraccion-binaria}, \cref{tab:formatos} y \cref{teo:c2}.
\pasos
\begin{enumerate}[label=(\alph*)]
  \item Falsa: $1/10$ tiene denominador no potencia de~2, su desarrollo binario es infinito.
  \item Verdadera: ambos tienen 8 bits de exponente; máximos $\approx3{,}39\cdot10^{38}$ y $3{,}40\cdot10^{38}$.
  \item Verdadera: $T(w)<0$ si y solo si $a_{n-1}=1$.
  \item Verdadera: $0{,}625=5/8=(0{,}101)_2=(1{,}01)_2\cdot2^{-1}$ necesita solo 2 bits de fracción.
\end{enumerate}
\resultado \textbf{(a) falsa; (b) verdadera; (c) verdadera; (d) verdadera.}
\comprobacion (b) $(2-2^{-7})2^{127}\approx3{,}390\cdot10^{38}$ y $(2-2^{-23})2^{127}\approx3{,}403\cdot10^{38}$.
\end{solucion}

\begin{solucion}{eje:6-a-binario}
\planteamiento divisiones sucesivas entre 2 (\cref{ej:conversion}) y agrupación en cuartetos.
\pasos
\begin{enumerate}[label=\arabic*.]
  \item $200=128+64+8=(1100\,1000)_2=\texttt{0xC8}$.
  \item $1000=512+256+128+64+32+8=(11\,1110\,1000)_2=\texttt{0x3E8}$.
  \item $4095=2^{12}-1=(1111\,1111\,1111)_2=\texttt{0xFFF}$.
\end{enumerate}
\resultado $200=\texttt{0xC8}$, $1000=\texttt{0x3E8}$ y $4095=\texttt{0xFFF}$, con las expresiones binarias de los pasos 1--3.
\comprobacion Horner: $12\cdot16+8=200$; $(3\cdot16+14)\cdot16+8=1000$; $(15\cdot16+15)\cdot16+15=4095$.
\end{solucion}

\begin{solucion}{eje:6-fracciones}
\planteamiento multiplicar sucesivamente por 2 y anotar la parte entera.
\pasos
\begin{enumerate}[label=(\alph*)]
  \item $0{,}625\to1{,}25$ (1), $0{,}25\to0{,}5$ (0), $0{,}5\to1$ (1): $(0{,}101)_2$.
  \item $13/16=(1101)_2/2^4=(0{,}1101)_2$.
  \item $1/3\to2/3$ (0), $2/3\to4/3$ (1), $1/3\to\dots$: $(0{,}\overline{01})_2$.
  \item $0{,}1\to0{,}2$ (0), $0{,}4$ (0), $0{,}8$ (0), $1{,}6$ (1), $1{,}2$ (1), $0{,}4$ (0), \dots: $(0{,}0\overline{0011})_2$.
\end{enumerate}
\resultado \textbf{(a) $0{,}101$; (b) $0{,}1101$; (c) $0{,}\overline{01}$; (d) $0{,}0\overline{0011}$.} Son finitos (a) y (b), de denominadores $8$ y $16$; (c) y (d) tienen denominadores $3$ y $10$, que no son potencias de~2.
\comprobacion (c) $\sum_{k\ge1}4^{-k}=\frac{1/4}{1-1/4}=\frac13$.
\end{solucion}

\begin{solucion}{eje:6-c2}
\planteamiento \cref{def:c2} y \cref{teo:c2}(iii): opuesto $=$ invertir y sumar~1.
\pasos
\begin{enumerate}[label=\arabic*.]
  \item (a) $1=(0000\,0001)_2$, invertido $1111\,1110$, más 1: $-1=(1111\,1111)_2$.
  \item $100=(0110\,0100)_2$, invertido $1001\,1011$, más 1: $-100=(1001\,1100)_2$. $127=(0111\,1111)_2$.
  \item (b) $(1000\,0001)_2=-128+1=-127$; $(1111\,0110)_2=-128+64+32+16+4+2=-10$.
  \item (c) $100+50=150>127$: desbordamiento. Se obtiene $150-256=-106$, es decir, $(1001\,0110)_2$.
\end{enumerate}
\resultado \textbf{(a) $1111\,1111$, $1001\,1100$, $0111\,1111$; (b) $-127$ y $-10$; (c) el resultado de 8 bits es $-106$ (desbordamiento).}
\comprobacion $(1001\,1100)_2=-128+16+8+4=-100$; $(1001\,0110)_2=-128+16+4+2=-106$.
\end{solucion}

\begin{solucion}{eje:6-utf8}
\planteamiento escribir el punto de código en binario con 11 o 16 bits y repartir los bits en las plantillas.
\pasos
\begin{enumerate}[label=\arabic*.]
  \item «é»: $\texttt{0xE9}=(000\,1110\,1001)_2$ con 11 bits; se reparte como $00011\mid101001$.
  \item Bytes: $110\,00011=\texttt{0xC3}$ y $10\,101001=\texttt{0xA9}$.
  \item «€»: $\texttt{0x20AC}=(0010\,0000\,1010\,1100)_2$; se reparte como $0010\mid000010\mid101100$.
  \item Bytes: $1110\,0010=\texttt{0xE2}$, $10\,000010=\texttt{0x82}$, $10\,101100=\texttt{0xAC}$.
\end{enumerate}
\resultado «é» se codifica con los bytes \texttt{C3 A9} y «€» con \texttt{E2 82 AC}.
\comprobacion al decodificar \texttt{C3 A9}, los bits útiles $00011$ y $101001$ dan $3\cdot64+41=233=\texttt{0xE9}$.
\end{solucion}

\begin{solucion}{eje:6-digitos}
\planteamiento número de cifras $=\lfloor\log_bn\rfloor+1$ (\cref{prop:cifras}).
\pasos
\begin{enumerate}[label=\arabic*.]
  \item (a) $2^{29}=536\,870\,912\le10^9<2^{30}=1\,073\,741\,824$: $\lfloor\log_210^9\rfloor=29$ y hay $30$ cifras.
  \item (b) $\log_{10}2^{100}=100\cdot0{,}30103=30{,}103$: $31$ cifras.
  \item (c) $2^{32}-1=(\texttt{FFFFFFFF})_{16}$: $8$ cifras.
\end{enumerate}
\resultado \textbf{(a) $30$; (b) $31$; (c) $8$.}
\comprobacion (b) $2^{100}=1\,267\,650\,600\,228\,229\,401\,496\,703\,205\,376$ tiene 31 cifras.
\end{solucion}

\begin{solucion}{eje:6-decodificar}
\planteamiento separar signo, exponente y fracción (\cref{def:flotante}); sesgo $127$ en fp32 y bf16.
\pasos
\begin{enumerate}[label=\arabic*.]
  \item $\texttt{0x3F800000}=0\ 01111111\ 000\dots0$: $s=0$, $E=127$, $f=0$; valor $1{,}0\cdot2^0=1$.
  \item $\texttt{0xC0A00000}=1\ 10000001\ 0100\dots0$: $s=1$, $E=129$, $f=(0{,}01)_2=0{,}25$; valor $-1{,}25\cdot2^{2}=-5$.
  \item $0{,}75=(1{,}1)_2\cdot2^{-1}$: $E=126=(01111110)_2$, $f=1000000$. Palabra $0\,01111110\,1000000=(0011\,1111\,0100\,0000)_2=\texttt{0x3F40}$.
  \item $3=(1{,}1)_2\cdot2^{1}$: $E=128=(10000000)_2$. Palabra $0\,10000000\,1000000=\texttt{0x4040}$.
\end{enumerate}
\resultado (a) $1$ y $-5$; (b) en bf16, $0{,}75$ es \texttt{0x3F40} y $3$ es \texttt{0x4040}.
\comprobacion los 16 bits superiores de los fp32 de $0{,}75$ ($\texttt{0x3F400000}$) y $3$ ($\texttt{0x40400000}$) coinciden con las palabras bf16, como se explicó tras la \cref{tab:formatos}.
\end{solucion}

\begin{solucion}{eje:6-espaciado}
\datos bf16: $p=7$; fp32: $p=23$; fp16: $e=5$, $p=10$, sesgo $15$.
\planteamiento en $[2^k,2^{k+1})$ los números son $(1{,}f)_2\cdot2^k$ con $f$ de $p$ bits.
\pasos
\begin{enumerate}[label=\arabic*.]
  \item (a) Al aumentar $f$ en una unidad del último bit, el número aumenta $2^{-p}\cdot2^k=2^{k-7}$.
  \item (b) En $[256,512)=[2^8,2^9)$ la distancia es $2^{1}=2$: los representables cercanos son $256$ y $258$. $257$ está a la misma distancia; el empate se resuelve hacia la fracción par, $256$ ($f=0000000$). Luego $\fl(256+1)=256$ y $\fl(\fl(256+1)+1)=\fl(257)=256$.
  \item $\fl(1+1)=2$ y $\fl(256+2)=258$, que es representable.
  \item (c) Hasta $2^8=256$ la distancia es $\le1$ y todos los enteros son representables; $257$ es el primero que no lo es. En fp32, por el mismo argumento, $2^{24}+1=16\,777\,217$.
  \item (d) El mayor exponente útil de fp16 es $E=2^5-2=30$, es decir, $2^{30-15}=2^{15}$, con fracción máxima $1-2^{-10}$: $(2-2^{-10})\cdot2^{15}=65\,536-32=65\,504$.
\end{enumerate}
\resultado \textbf{(a) $2^{k-7}$; (b) $256$, $256$ y $258$: la suma no es asociativa; (c) $257$ en bf16 y $16\,777\,217$ en fp32; (d) $65\,504$.}
\comprobacion (c) en bf16, $256=(1{,}0000000)_2\cdot2^8$ y $258=(1{,}0000001)_2\cdot2^8$ son consecutivos.
\end{solucion}

\begin{solucion}{eje:6-memoria}
\datos $7\cdot10^9$ parámetros; 4, 2, 1 y $0{,}5$ bytes por parámetro.
\planteamiento memoria $=$ parámetros $\times$ bytes por parámetro.
\pasos
\begin{enumerate}[label=\arabic*.]
  \item fp32: $2{,}8\cdot10^{10}$ bytes $=28$ GB; en GiB, $2{,}8\cdot10^{10}/2^{30}\approx26{,}08$ GiB.
  \item bf16: $14$ GB. int8: $7$ GB. 4 bits: $3{,}5$ GB.
\end{enumerate}
\resultado \textbf{$28$ GB ($26{,}08$ GiB), $14$ GB, $7$ GB y $3{,}5$ GB.}
\comprobacion cada reducción a la mitad de bits divide la memoria por~2; $28/26{,}08\approx1{,}074$, el factor entre GiB y GB.
\end{solucion}

\begin{solucion}{eje:6-divisibilidad}
\datos $n=\sum_kd_k10^k$.
\planteamiento por el \cref{eje:3-diferencia-potencias}, $10^k-1=(10-1)(10^{k-1}+\dots+1)$ es múltiplo de~9 y, por tanto, de~3.
\pasos
\begin{enumerate}[label=\arabic*.]
  \item $n-\sum_kd_k=\sum_kd_k(10^k-1)$, que es múltiplo de~9.
  \item Si dos enteros difieren en un múltiplo de~9 (o de~3), uno es múltiplo de~9 (o de~3) si y solo si el otro lo es.
  \item Base 16: $16^k-1=(16-1)(16^{k-1}+\dots+1)$ es múltiplo de 15, luego $n-\sum_kh_k$ es múltiplo de 15 si $n=\sum_kh_k16^k$.
  \item Criterio: $n$ es divisible por 15 (y, por tanto, por 3 y por 5) si y solo si lo es la suma de sus cifras hexadecimales.
\end{enumerate}
\resultado \textbf{Los criterios son ciertos en base 10 (para 3 y 9) y en base 16 (para 15, 3 y 5).}
\comprobacion $\texttt{0x3C}=60$: $3+12=15$, y $60=4\cdot15$.
\end{solucion}

\begin{solucion}{eje:6-extension}
\datos $w=(a_{n-1}\dots a_0)$, $w'=(a_{n-1}a_{n-1}a_{n-2}\dots a_0)$ con $n+1$ bits.
\planteamiento aplicar la \cref{def:c2} a $w'$ y separar casos según el bit de signo.
\pasos
\begin{enumerate}[label=\arabic*.]
  \item $T(w')=-a_{n-1}2^{n}+a_{n-1}2^{n-1}+\sum_{k=0}^{n-2}a_k2^k$.
  \item $-2^n+2^{n-1}=-2^{n-1}$, luego $T(w')=-a_{n-1}2^{n-1}+\sum_{k=0}^{n-2}a_k2^k=T(w)$.
\end{enumerate}
\resultado \textbf{$T(w')=T(w)$}; por inducción, la extensión a cualquier número de bits conserva el valor.
\comprobacion $-5=(1011)_2$ en 4 bits y $(1111\,1011)_2$ en 8 bits (\cref{ej:c2}).
\end{solucion}

\begin{solucion}{eje:6-decimal-finito}
\datos $p/q$ irreducible, $q\ge1$.
\planteamiento como en la \cref{prop:fraccion-binaria}, con $10=2\cdot5$ y el teorema fundamental de la aritmética.
\pasos
\begin{enumerate}[label=\arabic*.]
  \item Si el desarrollo tiene $k$ cifras decimales, $p/q=N/10^k$, luego $p\,10^k=qN$ y $q\mid p\,10^k$.
  \item Como $q$ no comparte factores primos con $p$, $q\mid10^k=2^k5^k$, y sus únicos factores primos posibles son 2 y 5: $q=2^a5^b$.
  \item Recíprocamente, si $q=2^a5^b$ y $m=\max(a,b)$, entonces $\frac pq=\frac{p\,2^{m-a}5^{m-b}}{10^m}$, que tiene a lo sumo $m$ cifras decimales tras la coma.
\end{enumerate}
\resultado \textbf{$p/q$ tiene desarrollo decimal finito si y solo si $q=2^a5^b$.}
\comprobacion $\frac3{40}=\frac3{2^35}=\frac{75}{1000}=0{,}075$; $\frac16$ tiene el factor~3 y es periódico.
\end{solucion}

\begin{solucion}{eje:6-desbordamiento}
\datos $x,y\in[-2^{n-1},2^{n-1}-1]$; $s\equiv x+y\pmod{2^n}$ con $s$ en el rango.
\planteamiento acotar $x+y$ en cada caso; hay desbordamiento si y solo si $x+y$ sale del rango.
\pasos
\begin{enumerate}[label=\arabic*.]
  \item (a) Si $x\ge0>y$: $y\le x+y<x$, luego $x+y$ está en el rango y $s=x+y$. El caso $y\ge0>x$ es simétrico.
  \item (b) Si $x,y\ge0$: $0\le x+y\le2^n-2$. Si $x+y\le2^{n-1}-1$, no hay desbordamiento y $s\ge0$. Si $x+y\ge2^{n-1}$, hay desbordamiento y $s=x+y-2^n\in[-2^{n-1},-2]$, negativo.
  \item Si $x,y<0$: $-2^n\le x+y\le-2$. Si $x+y\ge-2^{n-1}$, $s=x+y<0$. Si $x+y<-2^{n-1}$, hay desbordamiento y $s=x+y+2^n\in[0,2^{n-1}-1]$, no negativo.
  \item Juntando los casos: hay desbordamiento exactamente cuando los sumandos tienen el mismo signo y $s$ tiene el contrario.
\end{enumerate}
\resultado \textbf{Sumandos de signo distinto nunca desbordan; el desbordamiento equivale a un cambio de signo con sumandos del mismo signo.}
\comprobacion $n=8$: $100+50$ da $s=-106<0$ (\cref{eje:6-c2}).
\end{solucion}

\begin{solucion}{eje:6-redondeo}
\datos $2^k\le\abs x<2^{k+1}$; $p$ bits de fracción; $\fl$ redondea al más cercano.
\planteamiento los números representables en $[2^k,2^{k+1}]$ forman una rejilla de paso $2^{k-p}$ (\cref{eje:6-espaciado}); el error es como mucho medio paso.
\pasos
\begin{enumerate}[label=\arabic*.]
  \item Sin pérdida de generalidad $x>0$. Los números $2^k+j\,2^{k-p}$, $j=0,\dots,2^p$, son representables (para $j=2^p$ se obtiene $2^{k+1}$).
  \item $x$ está entre dos consecutivos, $c\le x\le c+2^{k-p}$, y el más cercano dista como mucho $\frac12\cdot2^{k-p}=2^{k-p-1}$.
  \item Como $\fl(x)$ es el más cercano de todo el formato, $\abs{\fl(x)-x}\le2^{k-p-1}$.
  \item Dividiendo por $\abs x\ge2^k$: $\frac{\abs{\fl(x)-x}}{\abs x}\le\frac{2^{k-p-1}}{2^k}=2^{-(p+1)}$.
  \item Escribiendo $\delta=(\fl(x)-x)/x$ se obtiene $\fl(x)=x(1+\delta)$ con $\abs\delta\le u$.
\end{enumerate}
\resultado \textbf{$\abs{\fl(x)-x}\le2^{k-p-1}$ y el error relativo es $\le u=2^{-(p+1)}$.}
\comprobacion en bf16, $x=257$: $k=8$, $p=7$; $\abs{256-257}=1=2^{8-7-1}$ y $1/257\approx0{,}003891\le2^{-8}\approx0{,}003906$.
\end{solucion}


\backmatter
% =====================================================================
\begin{thebibliography}{99}
\addcontentsline{toc}{chapter}{Bibliografía}
\markboth{Bibliografía}{}
\small

\bibitem{openstax}
J.~Abramson.
\newblock \emph{College Algebra}, 2.ª ed.
\newblock OpenStax, Rice University, 2021.
\newblock \url{https://openstax.org/details/books/college-algebra-2e}

\bibitem{blanchard}
P.~Blanchard, D.~J. Higham y N.~J. Higham.
\newblock Accurate computation of the log-sum-exp and softmax functions.
\newblock Preprint, 2019.
\newblock \url{https://arxiv.org/abs/1909.03469}

\bibitem{csapp}
R.~E. Bryant y D.~R. O'Hallaron.
\newblock \emph{Computer Systems: A Programmer's Perspective}, 3.ª ed.
\newblock Pearson, 2016.
\newblock Cap.~2, «Representing and Manipulating Information». Texto de los cursos de sistemas de Carnegie Mellon.

\bibitem{ct}
T.~M. Cover y J.~A. Thomas.
\newblock \emph{Elements of Information Theory}, 2.ª ed.
\newblock Wiley, 2006.

\bibitem{goldberg}
D.~Goldberg.
\newblock What every computer scientist should know about floating-point arithmetic.
\newblock \emph{ACM Computing Surveys}, 23(1):5--48, 1991.

\bibitem{gkp}
R.~L. Graham, D.~E. Knuth y O.~Patashnik.
\newblock \emph{Concrete Mathematics: A Foundation for Computer Science}, 2.ª ed.
\newblock Addison-Wesley, 1994.
\newblock Cap.~2, «Sums».

\bibitem{hammack}
R.~Hammack.
\newblock \emph{Book of Proof}, 3.ª ed.
\newblock Richmond (Virginia), 2018.
\newblock \url{https://richardhammack.github.io/BookOfProof/}

\bibitem{ieee754}
IEEE.
\newblock \emph{IEEE Standard for Floating-Point Arithmetic} (IEEE Std 754-2019).
\newblock IEEE, 2019.

\bibitem{mcs}
E.~Lehman, F.~T. Leighton y A.~R. Meyer.
\newblock \emph{Mathematics for Computer Science}, revisión del 6 de junio de 2018.
\newblock Massachusetts Institute of Technology. Texto del curso 6.042J.
\newblock \url{https://courses.csail.mit.edu/6.042/spring18/mcs.pdf}

\bibitem{mit6042}
A.~R. Meyer y A.~Chlipala.
\newblock 6.042J: Mathematics for Computer Science.
\newblock Massachusetts Institute of Technology, MIT OpenCourseWare, primavera de 2015.
\newblock \url{https://ocw.mit.edu/courses/6-042j-mathematics-for-computer-science-spring-2015/}

\bibitem{spivak}
M.~Spivak.
\newblock \emph{Calculus}, 4.ª ed.
\newblock Publish or Perish, 2008.
\newblock Cap.~1, «Basic Properties of Numbers», y cap.~2, «Numbers of Various Sorts».

\end{thebibliography}


\end{document}
